% concatenated from rearrangement.tex
\documentclass[11pt]{article}
\usepackage{latexsym,color,xcolor,amsmath,amsthm,amssymb,amscd,amsfonts,stmaryrd,dsfont}
\usepackage{graphicx,bbm}
\usepackage{cancel, mathtools} 
\renewcommand{\baselinestretch}{1}

\setlength{\topmargin}{-0.45 in}     %% dist from automatic 1in baseline to top of header
\setlength{\oddsidemargin}{0.3in}  %% automatic 1in baseline
\setlength{\evensidemargin}{0.3in} 
\setlength{\textheight}{9in}
\setlength{\textwidth}{6.1in} 
\setlength{\footskip}{0.55in}  %% dist(footer bo

%\setlength{\textwidth}{5.3in} %\setlength{\evensidemargin}{0.7in}
%\setlength{\oddsidemargin}{0.7in} %\setlength{\textheight}{8.3in}
%\setlengthq{\topmargin}{-0.1in} \setlength{\parskip}{2mm}
%\setlength{\baselineskip}{1.7\baselineskip}
%Bonne 8

\iffalse
\newtheorem{theo}{Theorem}
\newtheorem{lem}[theo]{Lemma}v6mh3lo33yl6i893ll3o3ol3lo3lol683lhl
\newtheorem{cor}[theo]{Corollary}
\newtheorem{prop}[theo]{Proposition}
\newtheorem{defn}[theo]{Definition}
\newtheorem{rem}{Remarks}
 
\newtheorem{note}{Note}
\fi


\newtheorem{conjecture}{Conjecture}[section]
\newtheorem{theorem}[conjecture]{Theorem}
\newtheorem{ce}[conjecture]{Counter-example}
\newtheorem{remark}[conjecture]{Remark}
    \newtheorem{lemma}[conjecture]{Lemma}
\newtheorem{proposition}[conjecture]{Proposition}
\newtheorem{corollary}[conjecture]{Corollary}
\newtheorem{ques}[conjecture]{Question}
\newtheorem{cond}[conjecture]{Condition}
\newtheorem{definition}[conjecture]{Definition}
\newtheorem{example}[conjecture]{Example}
\newtheorem{cmt}{Comment}

\newcommand\independent{\protect\mathpalette{\protect\independent}{\perp}} 
\def\independent#1#2{\mathrel{\rlap{$#1#2$}\mkern2mu{#1#2}}} 


\newcommand{\supp}{\mathrm{Supp}}
\newcommand{\Conv}{\mathrm{Conv}} 
\newcommand{\cF}{\mathcal{F}} 
%\newcommand{\cE}{\mathcal{E}} 
\newcommand{\F}{\mathbb{F}}  
\newcommand{\mR}{\mathbb{R}} 
\newcommand{\R}{\mathbb{R}} 
\newcommand{\mN}{\mathbb{N}}
\newcommand{\mZ}{\mathbb{Z}}
\newcommand{\mC}{\mathbb{C}}
\renewcommand{\P}{\mathbb{P}}
\newcommand{\mF}{\mathbb{F}}
\newcommand{\E}{\mathbb{E}}
\newcommand{\mS}{\mathbb{S}}
\newcommand{\ima}{\mathrm{Im}\:}
\newcommand{\re}{\mathrm{Re}\:}
\newcommand{\tr}{\mathrm{tr}}
\newcommand{\Corr}{\mathrm{corr}}
\newcommand{\Cov}{\mathrm{cov}}
\newcommand{\pp}{\mathbb{P}}
\DeclareMathOperator{\Var}{Var}
\newcommand{\e}{\varepsilon}
%\newcommand{\lam}{\lambda}
\newcommand{\ent}{h}
\newcommand{\info}{I}
\newcommand{\rank}{\mathrm{rank}}
\DeclareMathOperator{\diag}{diag}
\newcommand{\A}{\mathcal{A}}
\newcommand{\X}{\mathcal{X}}
\newcommand{\Y}{\mathcal{Y}}
\newcommand{\Z}{\mathcal{Z}}
\newcommand{\jj}{[j]}
\newcommand{\ii}{[i]}
\newcommand{\ic}{[i^c]}
\newcommand{\J}{[J]}
\newcommand{\C}{[J^c]}
\newcommand{\K}{[J\setminus i]}
\newcommand{\m}{[m]}
\newcommand{\1}{[1]}
\newcommand{\M}{E_m}
\newcommand{\bX}{\bar{X}}
\newcommand{\bY}{\bar{Y}}
\newcommand{\bU}{\bar{U}}
\newcommand{\mat}{\mathrm{MAT}}
\newcommand{\JJ}{\hat{I}}

%%%%%%%%%%%%


\newcommand{\de}{\mathrm{d}}
%\newcommand{\e}{\mathrm{e}}
\newcommand{\conv}{\mathrm{conv}}
\newcommand{\Med}{\mathrm{Med}}
\newcommand{\relint}{\mathrm{relint}}
\newcommand{\vol}{\mathrm{Vol}}
\newcommand{\diam}{\mathrm{diam}}
\newcommand{\card}{\mathrm{card}}
\newcommand{\interior}{\mathrm{int}}
\newcommand{\proj}{\mathrm{P}}
\newcommand{\vect}{\mathrm{vect}}
%\newcommand{\supp}{\mathrm{supp}}
\newcommand{\dist}{\mathrm{dist}}
\newcommand{\Hess}{\mathrm{Hess}}

%\newcommand{\R}{\mathbb{R}}
%\newcommand{\N}{\mathbb{N}}
\renewcommand{\Z}{\mathbb{Z}}
%\newcommand{\F}{\mathcal{F}}


\newcommand{\N}{\mathbb{N}}
\newcommand{\ps}{\psi^{\star}_{(s)}}
\date{\vspace{-5ex}}




\date{\today}





\begin{document}

\title{Rearrangements and infimum convolutions}
\author{Devraj Duggal, James Melbourne, and Cyril Roberto}
\maketitle
\begin{abstract}
We develop a general comparison result for inf-convolution operators related to rearrangement. As a consequence we derive comparison results under spherically symmetric rearrangement for Laplace and polar transforms. As a by product we simplify existing proofs related to the functional Blaschke-Santalo's inequality of Keith Ball and derive a comparison result for some parabolic PDE.
\end{abstract}





\section{Introduction}

In his seminal papers \cite{talenti,talenti76}, Talenti used spherically symmetric rearrangements to compute the optimal constant in  Sobolev inequalities (see also \cite{aubin}) and to obtain a comparison result between the solutions of the elliptic equations $-\Delta u=f$ and $-\Delta v=f^*$ (the rearranged function $f^*$ of $f$). 
These two results open wide lines of investigation in functional analysis and PDE theory. 
%
As is well-known, rearrangements proved rather fruitful in various field of mathematics that include functional analysis, PDE and Probability Theory. To give just a partial list of applications is out of reach and we refer the reader to the books \cite{lieb-loss,baerstein,leoni} for general accounts on rearrangement.
%
One key feature of rearrangements, illustrated by Talenti's papers, is that it allows one to reduce a $n$-dimensional problem to a one-dimensional one which becomes simpler and sometimes solvable.

\medskip

In this paper our aim is to proceed along the same lines. We will provide, after \cite{melbourne2019rearrangement}, some new general rearrangement comparison result related to inf-convolution operators. This will be applied to the comparison of Laplace and polar transforms and some comparison of solutions of a parabolic equation that are given by inf-convolutions. Also, after rearrangement and reduction to dimension one,  we may provide simple proofs of existing results in convex geometry, namely on the existence of maximizers in the functional form of the Blaschke-Santalo inequality and its polar counterpart, together with the best constant in the former inequality.

\medskip

In the next section (Section \ref{sec:notations-definitions}), we introduce the necessary material and in particular the notion of inf-convolution of interest for us, together with a general notion of rearrangement (that includes the classical spherically symmetric rearrangement on the Euclidean space, the Gaussian rearrangement of Ehrhard and rearrangement on the sphere, and many others).

\smallskip

Section \ref{sec:general} is dedicated to the statement and the proof of one of our main theorems (Theorem \ref{th:isoperimetric-rearrangeement}).

\smallskip

In Section  \ref{sec:transforms} we apply our main theorem to compare the action of Laplace and polar transforms on a function and its rearranged counterpart.

\smallskip

Finally Section \ref{sec:applications} collects some applications. The first one is related to solution of some parabolic equations. The second and third give a simple proof of a known result by Ben Li \cite{li} on the existence of extremizers in the functional form of Blaschke-Santalo's inequality \cite{ball86,artstein-klartag-milman,fradelizi-meyer,lehec} in convex geometry, and a proof of such an inequality through a new semi-group argument developed by Cordero et al. \cite{cordero,CFL}, after Sakamura-Tsuji \cite{nakamura}.



\section{Notations and definitions}
\label{sec:notations-definitions}

In this section we collect the necessary definitions that we will use throughout the paper. 

We assume that $(\mathcal{X},\nu)$ and $(\mathcal{Y}, \mu)$ are Polish measure spaces endowed with their Borel $\sigma$-algebra. A \emph{cost function} is a measurable function $\phi \colon \mathcal{X} \times \mathcal{Y} \to \mathbb{R}$. Observe that we relax here the usual definition of a cost function, which, in the transport theory, is usually assumed to take only non-negative values.

One of the central objects in this paper is the infimum convolution operator.

\begin{definition}[Infimum convolution] \label{def:inf-convolution}
    For $\phi$ a cost function on $\mathcal{X} \times \mathcal{Y}$, and $f: \mathcal{Y} \to \mathbb{R}$ define the \emph{infimum convolution} $Q_\phi f: \mathcal{X} \to \mathbb{R} \cup \{-\infty\}$ by
    \[
        Q_\phi f (x) \coloneqq \inf_{y \in \mathcal{Y}} \left\{ f(y) + \phi(x,y) \right\}.
    \]
\end{definition}

%In some places authors consider cost functions taking the value $\infty$. We do not consider this here for simplicity. 
In many interesting situations related to transport theory or Hamilton Jacobi equations $\mathcal{X}=\mathcal{Y}=\mathbb{R}^n$ (or some smooth manifold) and the cost function is a function of the distance $\phi(x,y)=d(x,y)^p$ for $p \geq 1$, see \textit{e.g.}\ \cite{villani}. Then the infimum convolution $Q_\phi$ amounts to the classical Hopf-Lax formula \cite{evans}.
More general cost functions were considered in the literature, related to optimal transport in a weak sense \cite{GRST} with applications in economy \cite{Backhoff-Veraguas-Pammer}. In such applications, one considers $X=\mathbb{R}^n$ and Y to be a subset of all probability measures on $\mathbb{R}^n$.
Although such general costs enter the framework of our result, we will not consider them since no associated rearrangement seems known in the corresponding literature. %\footnote{which would suppose to have some optimal isoperimetric inequality on the set of probability measures}.


In the next definition we introduce the notion of "enlargement" associated to a cost function.


\begin{definition}[Enlargement]
    For a cost function $\phi$ on $\mathcal{X} \times \mathcal{Y}$ and an $\varepsilon \in \mathbb{R}$, define for $A \subseteq \mathcal{Y}$,
    \[
        A_\varepsilon \coloneqq A_{\phi,\varepsilon} \coloneqq \{x \in \mathcal{X} : \exists y \in A \text{ s.t. }  \phi(x,y) < \varepsilon \} 
    \]
    and 
     \[
        A_{\overline{{\varepsilon}}} \coloneqq  A_{\overline{\phi,\varepsilon}} \coloneqq \{x \in \mathcal{X} : \exists y \in A \text{ s.t. }  \phi(x,y) \leq \varepsilon \}.
    \]
\end{definition}


When $\mathcal{X} = \mathcal{Y}$ and $\phi$ is a distance, then $A_\varepsilon$ and $A_{\overline{\varepsilon}}$ agree with usual notion of an $\varepsilon$-enlargement of a set in the literature. Moreover, in this case, $A_\varepsilon$ and $A_{\overline{\varepsilon}}$ are only non empty for $\varepsilon > 0$ and $\varepsilon \geq 0$ respectively.
Observe that in the definition of the enlargement of a set, we allow $\varepsilon$ to take negative values. This comes from the fact that, in some situations, the cost might take negative values.

We denote the open sets of a topological space $E$ by $\mathcal{O}(E)$ and the Borel sets (the members of the smallest $\sigma$-algebra containing $\mathcal{O}(E)$) of a topological space $E$, by $\mathcal{B}(E)$.

The second central objet of this paper is the notion of \emph{rearrangement}.

\begin{definition}[Rearrangement-Decreasing rearrangement]
    A map $*: \mathcal{B}(\mathcal{Y}) \to \mathcal{O}(\mathcal{Y})$ is a \emph{rearrangement} under a measure $\mu$ if $\mu(A) = \mu(*(A))$ and $\mu(A) \leq \mu(B)$ implies $*(A) \subseteq *(B)$.  Further for a measurable function $f: \mathcal{Y} \to [0,\infty)$, its $*$-symmetric \emph{decreasing rearrangement} is defined by
    \begin{align*}
        f^*(y) \coloneqq  \int_0^\infty \mathbbm{1}_{ *(\{ f> \lambda \})} (y) d \lambda, \qquad y \in \mathcal{Y}.
    \end{align*}
\end{definition}

For brevity of notation we write $A^* \coloneqq *(A)$.  We will further assume that the map $*$ satisfies 
\begin{align} \label{eq: decreasing rearrangment characterization}
    \{ f > \lambda \}^*  = \{ f^* > \lambda \}.
\end{align}  Note that $\{ f > \lambda \}^*  = \{ f^* > \lambda \}$ is ensured (see Proposition 2.7 of \cite{melbourne2019rearrangement}) if $*$ is a rearrangement such that $A_i \subseteq A_{i+1}$ implies 
\begin{align} \label{eq: rearrangement requirement}
        \left( \bigcup_{i=1}^\infty A_i \right)^* \ \subseteq \ \bigcup_{i=1}^\infty A_i^*.
\end{align}

In this paper we will mainly consider the spherically symmetric decreasing rearrangement on $\mathbb{R}^n$, spherical caps on the sphere, and the Ehrhard Gaussian rearrangement \cite{ehrhard}. 
For a more complete zoology of rearrangements, we refer the reader  to Section 2 of the book by Kawohl \cite{kawohl}.
For $A$ a Borel set of $\mathbb{R}^n$,
its  spherically symmetric rearrangement is the ball $A^*$ centered at $0$ with the same Lebesgue measure, while, the Gaussian rearrangement of $A$ is the half space
$A^*=\{x \in \mathbb{R}^n: x_1 \leq a\}$ with $a$ chosen so that 
$A^*$ has the same (standard) Gaussian measure as $A$. 

One can also define increasing rearrangement as follows. 

\begin{definition}[Increasing rearrangement]
    For $* \colon \mathcal{B}(\mathcal{Y}) \to \mathcal{O}(\mathcal{Y})$ a rearrangement with respect to a measure $\mu$, and a $\mu$-measurable function $f: \mathcal{Y} \to \mathbb{R}$, we define the $*$-symmetric \emph{increasing rearrangement} by,
    \[
        f_*(y) = - \log (e^{-f})^*.
    \]
\end{definition}

Note that by \eqref{eq: decreasing rearrangment characterization}, 
\begin{align} \label{eq: increasing rearrangement characterization}
    \{ f_* < \lambda \} = \{ f < \lambda \}^*
\end{align}
as $\{ f_* < \lambda \} = \{ (e^{-f})^* > e^{-\lambda} \} = \{ e^{-f} > e^{- \lambda } \}^* = \{ f < \lambda \}^*$.

As is well known, rearrangements are related to isoperimetric inequalities.
For the 3 classical isoperimetric inequalities alluded to above (Euclidean on $\mathbb{R}^n$, the sphere and the Gaussian space), the isoperimetric sets (realizing equality in the isoperimetric inequality) are nested 
in the sense that one contains the other and ,moreover, the set of extremal sets is totally ordered. This property leads to the following definition.

% \begin{definition}[Isoperimetric rearrangement]
% \label{def:isoperimetric-rearrangement}
%     We consider a map  $*: \mathcal{B}(\mathcal{Y}) \to \mathcal{O}(\mathcal{Y})$ to be an \emph{isoperimetric rearrangement} for the cost function $\phi$ and measure $\mu$, if $* (\mathcal{B}(\mathcal{Y}))$ is totally ordered and any Borel set $A \subseteq \mathcal{Y}$ and $\varepsilon \in \mathbb{R}$ satisfy
%     \[
%         \mu(A_{\phi,\varepsilon}) \geq \mu((A^*)_{\phi,\varepsilon}).
%     \]
%     In that case, we may say that $(*,\mu,\phi)$ is an isoperimetric rearrangement.
% \end{definition}



% Similarly, when the complement of the isoperimetric sets are nested, we define the isoperimetric-co-rearrangement.

% \begin{definition}[Isoperimetric co-rearrangement]\label{def:isoperimetric-co-rearrangement}
%     We consider a map  $\star: \mathcal{B}(\mathcal{Y}) \to \mathcal{O}(\mathcal{Y})$ to be an \emph{isoperimetric co-rearrangement} for the cost function $\phi$ and measure $\mu$, if $(\star (\mathcal{B}(\mathcal{Y})))^c$ -- the family of complements of rearranged sets -- is totally ordered and any Borel set $A \subseteq \mathcal{Y}$ and $\varepsilon \in \mathbb{R}$ satisfy
%     \[
%         \mu((A_{\phi,\varepsilon})^c) \leq \mu(((A^\star)_{\phi,\varepsilon})^c) .
%     \]
%     In that case, we may say that $(\star,\mu,\phi)$ is an isoperimetric co-rearrangement.
% \end{definition}


\begin{definition}
    \label{def:isoperimetric-rearrangement}
    We consider a map  $*: \mathcal{B}(\mathcal{Y}) \to \mathcal{O}(\mathcal{Y})$ to be an \emph{isoperimetric rearrangement} for the cost function $\phi$ and measure $\mu$, if $* (\mathcal{B}(\mathcal{Y}))$ is totally ordered and any Borel set $A \subseteq \mathcal{Y}$ and $\varepsilon \in \mathbb{R}$ satisfy
    \begin{align} \label{eq: rearrangement one way}
        \mu(A_{\phi,\varepsilon}) \geq \mu((A^*)_{\phi,\varepsilon}).
    \end{align}
    and
    \begin{align} \label{eq: rearrangement two way}
        \mu((A_{\phi,\varepsilon})^c) \leq \mu(((A^\star)_{\phi,\varepsilon})^c) 
    \end{align}
    In this case, we may say that $(*,\mu,\phi)$ is an isoperimetric rearrangement.
\end{definition}

In \eqref{eq: rearrangement one way} and \eqref{eq: rearrangement two way} we adopt the convention that 
$\infty \leq \infty$. 
In particular, if $\mu$ is a finite measure, then \eqref{eq: rearrangement one way} and \eqref{eq: rearrangement two way} are equivalent.  If $\mu$ is  not a finite measure, $\mu((A^*)_{\phi,\varepsilon})$ is always finite and \eqref{eq: rearrangement one way} holds, then \eqref{eq: rearrangement two way} holds trivially.  Similarly if $\mu(((A^\star)_{\phi,\varepsilon})^c) $ is always finite and \eqref{eq: rearrangement two way} holds, then $\mu((A_{\phi,\varepsilon})^c)$ is always finite and hence $\mu((A_{\phi,\varepsilon}))$ is always infinite implying \eqref{eq: rearrangement one way} trivially. 





\begin{remark} \label{rem:transfer}
If $\Phi \colon \mathcal{X} \times \mathcal{Y} \to \mathbb{R}$ is a cost function and $\alpha \colon \mathrm{Im(\phi)} \to \mathbb{R}$ is a measurable one to one map from the image $\mathrm{Im(\phi)} \subset \mathbb{R}$ of $\phi$ onto its image $\mathrm{Im}(\alpha)$, then $\psi = \alpha(\Phi)$  is a cost function on $\mathcal{X} \times \mathcal{Y}$ and $(*,\mu,\phi)$
is an isoperimetric rearrangement/co-rearrangement if and only if 
$(*,\mu,\psi)$
is an isoperimetric rearrangement/co-rearrangement 
\end{remark}


We now illustrate  these definitions 
with 3 classical rearrangements, namely the spherically symmetric rearrangement on the Euclidean space, the Gaussian rearrangement of Ehrhard, and the symmetric rearrangement on the sphere.

\subsubsection*{Spherically symmetric rearrangements on the Euclidean space}
Consider $\mathcal{X}=\mathcal{Y}=\mathbb{R}^n$ equipped with the Lebesgue measure $\lambda$, \textit{i.e.}\ $\mu=\nu=\lambda$. For sets, we interchangeably use the notation $|\cdot|$ for the Lebesgue measure.

Set $B \coloneqq \{ x : \|x\| < 1 \}$ for the open unit ball of $\mathbb{R}^n$. Here and in the sequel $\|\cdot\|$ stands for the Euclidean norm.

As mentioned earlier, given a Borel set $A$ of $\mathbb{R}^n$, its spherically symmetric rearrangement is $A^*=rB=\{x \in \mathbb{R}^n : \|x\| < r\}$, with $r$ chosen so that $|A^*|=|A|$. 

The classical isoperimetric inequality for the Lebesgue measure on $\mathbb{R}^n$ can be seen, as is well known \cite{ledoux}, as a special case of the celebrated Brunn-Minkowski inequality. This inequality may be formulated as follows \cite{gardner}: for Borel sets $A,C$ of finite measure, their Minkowsky sum $A+C \coloneqq \{a+c, \; a \in , c \in C\}$ decreases on spherically decreasing rearrangement, \textit{i.e.}\ 
$$
|A+C| \geq |A^*+C^*|.
$$
Now, consider the cost $\phi(x,y)=\|x-y\|$.
Since spherically symmetric rearrangements are centered balls (that are totally ordered), applying the Brunn-Minkoswki inequality
to $A$ of finite Lebesgue measure and the unit ball $C=B$ and noting that 
$A_\varepsilon = A + \varepsilon B$
and
$(A^*)_\varepsilon = A^* + \varepsilon B$ (that is a ball), 
$$
|A_\varepsilon| \geq |(A^*)_\varepsilon |.
$$
In particular, the spherically symmetric rearrangement is an isoperimetric rearrangement for the cost $\|x-y\|$ and the Lebesgue measure $\lambda$ as in the sense of Definition \ref{def:isoperimetric-rearrangement}.
By Remark \ref{rem:transfer} this property transfers to any cost $\alpha(\|x-y\|)$ for a proper choice of the map $\alpha$.

\subsubsection*{Ehrhard Gaussian rearrangement on the Euclidean space}
Consider again $\mathcal{X}=\mathcal{Y}=\mathbb{R}^n$ with the cost $\phi(x,y)=\|x-y\|$ but now equipped with the standard Gaussian measure  $\gamma$ with density $\varphi(x)=e^{-|x|^2/2}/(2\pi)^{n/2}$, $x \in \mathbb{R}^n$. Let $\Phi(t)=\int_{-\infty}^t e^{-s^2/2}\frac{ds}{\sqrt{2\pi}}$, $t \in \mathbb{R}$. 

Given a Borel set $A$ of $\mathbb{R}^n$, its Gaussian rearrangement is 
$A^*=\{x \in \mathbb{R}^n : x_1 < a\}$
where $a=\Phi^{-1}(\gamma(A))$ is chosen so that $\gamma(A^*)=\gamma(A)$.
Observe that the choice of the first coordinate is arbitrary and one could equivalently consider rearrangements of the form 
$A^*=\{x \in \mathbb{R}^n : \langle x , u \rangle < a\}$, for any unit vector $u$, where $\langle\cdot,\cdot \rangle$ stands for the Euclidean scalar product.

Then, the Gaussian isoperimetric inequality of Borell \cite{borell}, Sudakov and Tsireslon \cite{sudakov}
asserts that for any Borel set $A$
$$
\gamma(A_\varepsilon) \geq \Phi( \Phi^{-1}(\gamma(A)) + \varepsilon).
$$
By construction 
$\Phi( \Phi^{-1}(\gamma(A)) + \varepsilon)) = \gamma(\{x : x_1 \leq \Phi^{-1}(\gamma(A)) + \varepsilon\})$.
On the other hand, 
for $y \in (A^*)_\varepsilon$ and 
$x \in A^*=\{x: x_1 \leq \Phi^{-1}(\gamma(A))\}$ with $\|x-y\| \leq \varepsilon$, it holds that
$y_1=x_1+y_1-x_1 \leq \Phi^{-1}(\gamma(A)) + \varepsilon$. Hence 
$(A^*)_\varepsilon \subset \{x : x_1 \leq \Phi^{-1}(\gamma(A)) + \varepsilon\}$ and in turn, from the Gaussian isoperimetric inequality
$$
\gamma(A_\varepsilon) \geq \gamma((A^*)_\varepsilon).
$$
In conclusion, the Ehrhard Gaussian rearrangement is an isoperimetric rearrangement for the cost $\|x-y\|$ and the standard Gaussian measure $\gamma$. 
Again from Remark \ref{rem:transfer}, this property transfers to more general costs $\alpha(\|x-y\|)$.


\subsubsection*{Isoperimetry on the sphere}
An other classical example is given by the sphere $\mathcal{X}=\mathcal{Y}=S^{n-1}\subseteq\mathbb{R}^{n}$ equipped with the normalized uniform measure $\mu=\nu=\sigma_{n-1}$. Let $d$ denote the standard geodesic distance on the sphere. For a Borel set $A\in\mathcal{B}(S^{n-1})$, let $A_{\varepsilon}=\{x\in S^{n-1}:d(x,A)<\varepsilon\}$ be the usual enlargement related to $d$. 
By convention, we speak of spherical caps centered at the north pole (\textit{i.e.}\ the geodesic balls centered at the north pole). The rearrangement of interest associates $A\subseteq S^{n-1}$ to $A^*\subseteq S^{n-1}$ where $A^*$ is the cap centered at the north pole such that $\sigma_{n-1}(A)=\sigma_{n-1}(A^*)$.
The classical Isoperimetry result due to Lévy (see \textit{e.g.}\ \cite{ledoux}) states that 
$$
\sigma_{n-1}(A_\epsilon)\geq\sigma_{n-1}((A^*)_{\epsilon}) .
$$ 
Therefore, $(*,\sigma_{n-1},d)$ is an isoperimetric rearrangement and so is  
$(*,\sigma_{n-1},\alpha(d))$ for proper choices of the map $\alpha$ by Remark \ref{rem:transfer}.







\section{A general Argument} \label{sec:general}

Our aim is to prove the following general theorem on isoperimetric rearrangements and inf-convolution which follows and extends the chief ideas of \cite{melbourne2019rearrangement}.
At the end of the section we will give some examples of applications together with a weaker form of our general result when only a non optimal isoperimetric inequality is available.

\begin{theorem} \label{th:isoperimetric-rearrangeement}
    For $(*, \mu, \phi)$ an isoperimetric rearrangement and $\lambda \in \mathbb{R}$,
    \[
        \mu (\{ Q_\phi f < \lambda \}) \geq \mu (\{ Q_\phi f_* < \lambda \}) %.
   % \]
     %   For $(\star,\mu, \phi)$ an isoperimetric co-rearrangement, and $\lambda \in \mathbb{R}$,
   % \[
   \qquad \mbox{and} \qquad 
        \mu (\{ Q_\phi f \geq \lambda \}) 
        \leq \mu (\{ Q_\phi f_* \geq \lambda \}) .
    \]
\end{theorem}

For an infinite measure $\mu$, the conclusion of the above statement may read as $\infty \geq \infty$ (or $\infty \leq \infty$) which holds true by virtue of our convention. 

Observe that for the symmetric decreasing rearrangement and the Lebesgue measure on $\mathbb{R}^n$, that are isoperimetric rearrangement for some costs, the conclusion of the theorem involves the symmetric increasing rearranged function $f_*$ and not $f^*$.

The proof of Theorem \ref{th:isoperimetric-rearrangeement}
relies on the following decomposition lemma. The main idea here is to obtain a countable union of sets. Set $\mathbb{Q}$ for the set of rational numbers.

\begin{lemma} \label{lem: decomposition}
    For $\lambda \in \mathbb{R}$, the sublevel sets of an infimum convolution have the following decomposition,
    \[
        \{ Q_\phi f < \lambda \} = \bigcup_{q \in S(\lambda)} \{ f < q_1 \}_{\phi,q_2}
    \]
    where
    \[
        S(\lambda) \coloneqq \left\{ q=(q_1,q_2) \in \mathbb{Q}^2 : q_1 + q_2 < \lambda \right\}.
    \]
\end{lemma}

\begin{proof}
    If $x \in \{ f < q_1 \}_{\phi,q_2}$, for some $q \in S(\lambda)$ then there exists $y \in \{ f < q_1\}$ such that $\phi(x,y)< q_2$, thus
    \[
        f(y) + \phi(x,y) < q_1 + q_2 < \lambda
    \]
    so that $Q_\phi f(x) < \lambda$, giving $\bigcup_{q \in S(\lambda)} \{ f < q_1 \}_{\phi,q_2} \subseteq \{ Q_\phi f < \lambda \}$.
    For $x$ such that $Q_\phi f(x) < \lambda$, by definition there exists $y' \in \mathcal{Y}$ such that $f(y') + \phi(x,y') < \lambda$.  For $\delta > 0$, such a $y'$ belongs to $\{ f < f(y') + \delta \}$, as well as $\{ y : \phi(x,y) < \phi(x,y') + \delta \}$, thus
    \[
        x \in \{ f < f(y') + \delta \}_{\phi,\phi(x,y') + \delta} 
    \] 
    For $\delta$ chosen small enough, $f(y') + \phi(x,y') + 2 \delta < \lambda$.  Choosing $q = (q_1,q_2)$ such that $q_1 \in \mathbb{Q} \cap (f(y'), f(y') + \delta)$ and $q_2 \in (\phi(x,y'), \phi(x,y') + \delta)$ gives $x \in \{ f < q_1 \}_{\phi,q_2}$ for this $q \in S(\lambda)$, completing the proof.
\end{proof}




\begin{proof}[Proof of Theorem \ref{thm: Blaschke-Santalo rearrangement bis}]
The proof is obtained from the following computation: by Lemma \ref{lem: decomposition} and hypothesis (by definition of the isoperimetric rearrangement)        
    \begin{align*}
        \mu (\{ Q_\phi f < \lambda \} )
            &=
                \mu \left( \bigcup_{q \in S(\lambda)} \{ f < q_1 \}_{\phi,q_2} \right)
                %& (\mbox{by Lemma } \ref{lem: decomposition})
                    \\
            &\geq 
                \sup_{q \in S(\lambda)} \mu  \left( \{ f < q_1 \}_{\phi,q_2} \right)
                   \\
            &\geq 
                \sup_{q \in S(\lambda)} \mu  \left( (\{ f < q_1 \}^*)_{\phi,q_2} \right) .
                %& (\mbox{by definition of the rearrangement})     
                \\
  \end{align*}
Therefore, since the rearranged sets are totally ordered,          
    \begin{align*} 
    \mu (\{ Q_\phi f < \lambda \} )
            & \geq 
                \mu \left( \bigcup_{q \in S(\lambda)} (\{ f < q_1 \}^*)_{\phi,q_2} \right)
                %& (\mbox{by nestedness})   
                \\
            &=
                \mu \left( \bigcup_{q \in S(\lambda)} (\{ f_* < q_1 \})_{\phi,q_2} \right)
                    \\
            &=
                \mu (\{ Q_\phi f_* < \lambda \} ).
    \end{align*}
    where the first equality is \eqref{eq: increasing rearrangement characterization}, and the last equality is a second application of Lemma \ref{lem: decomposition}.

    The second part of the theorem follows the same line and uses the same arguments
       \begin{align*}
        \mu (\{ Q_\phi f \geq  \lambda \}) 
            &=
        \mu (\{ Q_\phi f <  \lambda \}^c)    
          =       \mu \left( \bigcap_{q \in S(\lambda)} \{ f < q_1 \}_{\phi,q_2}^c \right)
                %& (\mbox{by Lemma } \ref{lem: decomposition})
                    \\
            &\leq 
                \inf_{q \in S(\lambda)} \mu  \left( \{ f < q_1 \}_{\phi,q_2}^c \right)
            \leq 
                \inf_{q \in S(\lambda)} \mu  \left( (\{ f < q_1 \}^\star)_{\phi,q_2}^c \right)
                %& (\mbox{by definition of the rearrangement})     
                \\
            &=
                \mu \left( \bigcap_{q \in S(\lambda)} (\{ f < q_1 \}^*)_{\phi,q_2}^c \right)
                %& (\mbox{by nestedness})   
            =
                \mu \left( \bigcap_{q \in S(\lambda)} (\{ f_* < q_1 \})_{\phi,q_2}^c \right)
                    \\
            &=
                \mu (\{ Q_\phi f_* < \lambda \}^c) 
                 =
                 \mu (\{ Q_\phi f_* \geq \lambda \}) 
                 .
    \end{align*}
    This completes the proof of the theorem.
\end{proof}



As an illustrative example, one may consider the spherically symmetric rearrangement. It is an isoperimetric rearrangement for the cost 
$\phi(x,y)=tG(|x-y|/t)$, for $t>0$ and $G$ convex increasing (see the previous section). This special choice, and in particular the presence of the apparently superfluous parameter $t>0$, will prove useful for applications in Section \ref{sec:HJ}. 
By Theorem \ref{th:isoperimetric-rearrangeement} we can conclude that  the infimum convolution operator 
$$
Q_tf(x) \coloneqq \inf_{y \in \mathbb{R}^n} \left\{ f(y) + tG \left(\frac{|x-y|}{t} \right)  \right\}
$$
satisfies
\begin{equation} \label{eq:inf-conv}
|\{ Q_t f < \lambda \}| \geq |\{ Q_t f_* < \lambda \}| 
\end{equation}
for any $\lambda \in \mathbb{R}$.

% An other example is given by the sphere $S^{n-1}\subseteq\mathbb{R}^{n}$ equipped with the normalized uniform measure $\sigma_{n-1}$. Recall that $d$ denote the standard geodesic distance on the sphere. 
% Since
%  $(*,\sigma_{n-1},d)$ is an isoperimetric rearrangement. 
% Fix an increasing and convex function 
% $\alpha$ and define 
% $$
% Q_{\alpha}f(x)=\inf_{y\in S^{n-1}}f(y)+\alpha(d(x,y))
% $$ 
% Thanks to Remark \ref{rem:transfer},
% $(*,\sigma_{n-1}, \alpha(d))$ is an isoperimetric rearrangement. Therefore
% Theorem \ref{th:isoperimetric-rearrangeement} applies and leads to the following comparison
% $$
% \sigma_{n-1}\left(\{Q_{\alpha}f<\lambda\}\right)\geq\sigma_{n-1}\left(\{Q_{\alpha}f_{*}<\lambda\}\right) .
% $$

\medskip

We end this section with a final comment when an exact isoperimetric inequality need not hold. 
Besides the 3 classical examples given in the Section \ref{sec:notations-definitions}, there are some situations where the isoperimetric sets are known. To give a complete list of publications is out of reach. We refer the interested reader to \cite{ABFC} and references therein for recent results related to weighted measures on the upper half plane. For most of the exact isoperimetric inequalities the isoperimetric sets are totally ordered and the result of this paper apply. We will not continue in this direction and may instead consider situations where exact isoperimetric sets are not known. 
One important example is the class of strictly log-concave probability measures on $\mathbb{R}^{n}$
admitting the following density 
$$
\frac{d\mu}{dx}=e^{-V(x)}, 
\qquad x \in \mathbb{R}^n
$$ 
with $V''\geq \rho Id$, $\rho >0$, where $V''$ is a shorthand notation for the Hessian matrix of $V$, and where the inequality is understood in the quadratic sense.
Recall that 
$\Phi(t):=\int_{-\infty}^{t}\frac{1}{\sqrt{2\pi}}e^{-\frac{x^{2}}{2}}dx$, and set $\varphi=\Phi'$ for the density of the standard one dimensional Gaussian measure.
%$, \gamma(A)=\int_{A}\varphi(x)dx$$ 
It is  well known that $\mu$ satisfies the following isoperimetric inequality 
\begin{equation} \label{eq:levy}
    \mu^+(A)
    \geq 
    c
    \varphi\circ\Phi^{-1}(\mu(A)), \qquad A\in\mathcal{B}(\mathbf{R}^{n})
\end{equation} 
for some constant $c=c(\rho) \in (0,\infty)$
\cite{bakry-ledoux} (see also \cite{milman2010,ivanisvili-volberg,duggal2024curvature}).
When $V(x)=|x|^2/2-\frac{n}{2}\log(2\pi)$ the measure
$\mu$ is nothing but the $n$ dimensional standard Gaussian and \eqref{eq:levy} holds with constant $c=1$.
This is the classical Borell-Sudakov-Tsirelson Gaussian isoperimetry. 
For other  (than Gaussian) strictly log-concave probability measures, 
\eqref{eq:levy} is not optimal \cite{bobkov-houdre}.

A second situation that was considered and studied in the literature pertains to measures possessing a product structure. More precisely, these measures are of the form 
$$
\frac{d\mu_{n}}{dx}
=
\prod_{i= 1}^{n}\frac{1}{Z_{\Psi}}e^{-\Psi(|x_{i}|)}%=\prod_{i\geq 1}^{n}\rho(x_{i})
$$
where $\Psi:\mathbb{R}_{+}\rightarrow\mathbb{R}_+$ is an increasing convex function where $\Psi(0)=0$ and $\sqrt{\Psi}$ is concave (here $Z_\Psi$ is a normalization constant that turn $\mu_n$ into a probability measure). We refer the reader to \cite{roberto} for as account on isoperimetry for products of probability measures. The assumption on the potential $\Psi$ indicates that the measure is "between" exponential and Gaussian. In particular, it is not strictly convex and therefore this class of measures differ from the previous ones. 
Let $\rho=e^{-\Psi}/Z_\Psi$ and $H(t)=\int_{-\infty}^t \rho(x)dx$.
It has been shown \cite{bobkov-houdre97,BCR06,BCR07}
that there exists $c>0$ (independent of $n$) such that 
\begin{equation} \label{eq:iso-product}
\mu_{n}^+(A)\geq c \rho \circ H^{-1}(\mu_{n}(A)), \qquad  A\in\mathcal{B}(R^{n}) .
\end{equation}
Moreover, such an isoperimetric inequality for $c=1$ and $n \geq 2$ implies that $\mu$ is Gaussian \cite{bobkov-houdre}. Again, such an isoperimetric inequality is not exact and extremal isoperimetric sets are not known even for simple cases like the double sided exponential measure \cite{bobkov-houdre97,BCR06}.

Although the isoperimetric inequalities in the two classes of probability measures introduced above are not optimal, our approach on the rearrangement of the inf convolution operator is flexible enough to be applied in this situation as we explain in the following paragraph.

Observe first that the isoperimetric inequalities  
\eqref{eq:levy} and \eqref{eq:iso-product} can be reformulated in an integrated form that is more useful for our purpose. In fact, as is classical (see \textit{e.g.}\ \cite{ledoux}) \eqref{eq:levy} and \eqref{eq:iso-product} imply respectively 
$$
\mu(A_{\varepsilon})\geq\mu((A^*)_{c\varepsilon}) 
\qquad 
\mbox{and}
\qquad 
\mu_n(A_{\varepsilon})\geq\mu_n((A^*)_{c\varepsilon})
$$
where $A_{\varepsilon}:=\{x\in\mathbb{R}^{n},\exists y\in\mathbb{R}^{n}, \|x-y\|<\varepsilon\}$ and 
$A^*:=\{x\in\mathbb{R}^{n}:x_1<a\}$ with $a$ chosen such that
$\mu(A)=\mu(A^*)$ and similarly for $\mu_n$.
Now following the exact same proof of Theorem \ref{th:isoperimetric-rearrangeement} (with the cost $\phi(x,y)=t\|x-y\|$, details are left to the reader), the inf convolution operator
$$
Q_{t}f(x):=\sup_{y\in\mathbb{R}^{n}}f(y)-t\|x-y\|
$$
satisfies
$$
\mu\left(\{Q_{t}f>\lambda\}\right)
\geq
\mu\left(\{Q_{\frac{t}{c}}f_*>\lambda\}\right), \qquad t >0 
$$
and similarly for $\mu_n$. Thanks to Remark \ref{rem:transfer}, similar results hold for more general costs of the form $\phi(x,y)=\alpha(\|x-y\|)$.








\section{Spherically symmetric rearrangement of Transforms} \label{sec:transforms}

In this section we will prove a rearrangement result for some transforms that include the well known Legendre and polar transforms.

In all of this section, we consider a spherically symmetric rearrangement. In particular, $f^*$ and $A^*$ denote the spherically symmetric rearrangements of the function $f$ and of the set $A$ respectively as introduced in Section \ref{sec:notations-definitions}.



For a Borel measurable $f : \mathbb{R}^n \to \mathbb{R}$ and $\rho:\mathbb{R}\rightarrow\mathbb{R}$ an increasing function with $\rho(0)=0$, define the $\mathcal{T}$-transform of $f$,  $\mathcal{T}f:\mathbb{R}^n \to \mathbb{R} \cup\{\infty\}$
    \begin{equation*}%\label{def:legendre}
        \mathcal{T}f(x) = \sup_{y \in \mathbb{R}^n} \bigg[ \rho(\langle x , y  \rangle) -  f(y) \bigg], \qquad x \in \mathbb{R}^n
    \end{equation*}
where $\langle x , y  \rangle$ denote the Euclidean scalar product in $\mathbb{R}^n$.

Observe that when $\rho(x)=x$ the $\mathcal{T}$-transform
is the usual Legendre transform. %and the polar transform (see the end of the section). 
Such transforms were considered in the literature in relation with functional forms of the Blaschke-Santalo inequality \cite{ball86,fradelizi-meyer,lehec}. 


Recall that $|\cdot|$ stands for the Lebesgue measure.  
One of our main results about rearrangement is the following.
 
\begin{theorem} \label{thm: majorization rearrangement transform}
        If $f: \mathbb{R}^n \to [0,\infty)$ is convex, even and $f(0) = 0$, then
        \[
            | \{ \mathcal{T}f \leq \lambda \} | \leq  | \{ \mathcal{T}f_* \leq \lambda \} |,
            \qquad \lambda \in \mathbb{R}.
        \]
\end{theorem}

The proof of the theorem will be based on 
a formulation with rearrangement of the celebrated Blaschke-Santalo inequality from convex geometry that is now recalled. 
For a Borel set $K$ of $\mathbb{R}^n$, define the support function of $K$ by $h_K(x) = \sup_{y \in K} \langle x, y \rangle$, and the polar set
    \[
        K^\circ \coloneqq \{ x: h_K(x) \leq 1 \}
    \]
    when $K$ is non-empty, and $K^\circ = \mathbb{R}^n$ otherwise.
The Blaschke-Santalo inequality asserts that, for a convex symmetric (\textit{i.e.}\ $K=-K$) non-empty set $K$ with finite volume,
        \[
            |K||K^\circ| \leq |B|^2 
        \]
where we recall that $B$ is the Euclidean unit ball of $\mathbb{R}^n$.
We refer the reader to \cite{fradelizi-23} for an historical presentation and references.

Note that since $K^*$ is a ball,
$K^*=\left( \frac{|K|}{|B|}\right)^\frac{1}{n}B$
and therfeore $(K^*)^\circ = \left( \frac{|B|}{|K|}\right)^\frac{1}{n}B$ (since in general $(\lambda K)^\circ = \lambda^{-1} K^\circ$). Hence
$|K^*||(K^*)^\circ|=|B|^2$ and by the 
Blaschke-Santalo inequality (and since $|K^*|=|K|$)
$$
|K||K^\circ| \leq |K^*||(K^*)^\circ| = |K||(K^*)^\circ| 
$$
from which we infer that the Blaschke-Santalo inequality can be reformulated equivalently as a rearrangement inequality: for a symmetric convex set $K \subseteq \mathbb{R}^n$ with finite volume,
\begin{align} \label{thm: Blaschke-Santalo rearrangement bis}
    |K^\circ| \leq | (K^*)^\circ|.
\end{align}


\begin{proof}[Proof of Theorem \ref{thm: majorization rearrangement transform}]
    Let us first note that since $f$ takes its minimum $0$ at $0$, 
        \[
            \mathcal{T}f(x) = \sup_{y \in \mathbb{R}^n} \big[ \rho(\langle x,y\rangle) - f(y) \big] \geq \rho(\langle x ,0 \rangle) - f(0) = 0.
        \]
    Thus for $\lambda < 0$, $\{ \mathcal{T} f \leq \lambda \}$ is empty and there is nothing to prove.
   Hence we assume that $\lambda \geq 0$. Then with the cost function $\phi(x,y) = -\rho(\langle x,y \rangle$),
    \begin{align*}
        \{ \mathcal{T}f \leq \lambda \} 
            &=
                \{ Q_\phi f \geq - \lambda \}
                    \\
            &=
                \{ Q_\phi f < - \lambda \}^c
                    \\
            &=
               \left( \bigcup_{q \in S(-\lambda)} \{ f < q_1 \}_{\phi,q_2} \right)^c
               & (\mbox{by Lemma } \ref{lem: decomposition})
                    \\
            &=
                \bigcap_{q \in S(-\lambda)} \{ f < q_1 \}_{\phi,q_2}^c
    \end{align*}
    where 
    $S(-\lambda) = \{q = (q_1,q_2) \in \mathbb{Q}^2 :  q_1 + q_2 < - \lambda \}$.  
    Note that since $\{ f < q_1 \}$ is empty unless $q_1 > 0$, it suffices to consider the union over $q \in S_+(-\lambda)$, where $S_+(-\lambda) = S(-\lambda) \bigcap \{(x_1,x_2): x_1 > 0 \}$. 
    Summarizing, we have
    \begin{align} \label{eq: exp transform sublevel set}
        \{ \mathcal{T} f \leq \lambda \} = \bigcap_{q \in S_+(-\lambda)} \{ f < q_1 \}_{\phi,q_2}^c .
    \end{align}
    Note that for $q \in S_+(-\lambda)$, since $q_1 \geq 0$ and $q_1 + q_2 < -\lambda$, $q_2$ is necessarily negative.
    We claim that for $\varepsilon >0$, and a Borel set $K$,
     \begin{align*} %\label{eq: enlargement-transform with dot product cost}
        (K_{\phi,-\varepsilon})^c = \rho^{-1}(\varepsilon) \ K^\circ
    \end{align*} 
    where we defined 
    $$
    \rho^{-1}(\varepsilon)=\inf \{y \in\mathbb{R}:\rho(y)\geq\varepsilon\} .
    $$
    Indeed,
    \begin{align*}
        (K_{\phi,-\varepsilon})^c
            &=
                \{ x : \exists y \in K, \phi(x,y) < - \varepsilon \}^c
                    \\
            &=
                \left\{ x: \inf_{y \in K} - \rho(\langle x,y \rangle) < - \varepsilon \right\}^c
                    \\
            %&=
             %   \left\{ x : \sup_{y \in K } \rho(\langle x,y \rangle) > \varepsilon \right\}^c
             %       \\
            &=
                \left\{ x : \sup_{y \in K} \rho(\langle x,y \rangle) \leq \varepsilon \right\}
                    \\
            &=
                \left\{ x : \sup_{y \in K} \langle x,y \rangle \leq \rho^{-1}(\varepsilon) \right\}
                    \\
            &=
                \rho^{-1}(\varepsilon) \left\{ x : \sup_{y \in K} \langle x,y \rangle \leq 1 \right\}
                    \\
            &=
                \rho^{-1}(\varepsilon) \ K^\circ.
    \end{align*}
    Applying the claim to \eqref{eq: exp transform sublevel set}, we have 
    \begin{align} \label{eq:sub-level-description}
        \{ \mathcal{T} f \leq \lambda \} = \bigcap_{q \in S_+(-\lambda)} \rho^{-1}(|q_2|) \{ f < q_1 \}^\circ .
    \end{align}
    Thus,
    \begin{align*}
        |\{ \mathcal{T} f \leq \lambda \}| 
            &=
                \left|  \bigcap_{q \in S_+(-\lambda)} \rho^{-1}(|q_2|) \{ f < q_1 \}^\circ \right|
                    \\
            &\leq 
                \inf_{q \in S_+(-\lambda)} \rho^{-1}(|q_2|)^n \ \big| \{f < q_1 \}^\circ \big|
                    \\
            &\leq 
                \inf_{q \in S_+(- \lambda)} \rho^{-1}(|q_2|)^n \ \big| (\{ f < q_1 \}^*)^\circ \big |
                & (\mbox{by Blaschke-Santalo's inequality } \eqref{thm: Blaschke-Santalo rearrangement bis})
                    \\
            &= 
                \inf_{q \in S_+(- \lambda)}  \big| \rho^{-1}(|q_2|) \ (\{ f < q_1 \}^*)^\circ \big |
                & (\mbox{by homogeneity of the measure})
                    \\
            &=
                \left|\bigcap_{q \in S_+(- \lambda)} \rho^{-1}(|q_2|)  \ (\{ f < q_1 \}^*)^\circ \right|
                    \\
            &=
                \left|\bigcap_{q \in S_+(- \lambda)} \rho^{-1}(|q_2|) \ (\{ f_* < q_1 \})^\circ \right|
                    \\
            &=
                | \{ \mathcal{T} f_* \leq \lambda \} |
    \end{align*}
    where the first inequality is $\mu( \cap_n A_n) \leq \inf_n \mu(A_n)$ combined with an application of the homogeneity of the Lebesgue measure.  
    %The second inequality is Blaschke-Santalo's inequality \eqref{thm: Blaschke-Santalo rearrangement bis}. 
    %The following equality is again the homogeneity of the Lebesgue measure. 
    The antepenultimate equality is due to the fact that we are considering an intersection of centered balls which are totally ordered. The penultimate equality is the characterization of the spherically symmetric increasing rearrangement and the final one is \eqref{eq:sub-level-description} applied to the spherically symmetric increasing rearrangement $f_*$.
\end{proof}
We will now provide two direct applications of the above theorem.  
\begin{definition} \label{def:legendre}
    For a Borel measurable $f : \mathbb{R}^n \to \mathbb{R}$ define the Legendre transform of $f$,  $\mathcal{L}f:\mathbb{R}^n \to \mathbb{R}$
    \[
        \mathcal{L}f(x) = \sup_{y \in \mathbb{R}^n} \bigg[ \langle x , y  \rangle -  f(y) \bigg].
    \]
\end{definition}
 
\begin{theorem} \label{thm: majorization rearrangement Legendre}
        If $f: \mathbb{R}^n \to [0,\infty)$ is convex even and $f(0) = 0$, then
        \[
            | \{ \mathcal{L}f \leq \lambda \} | \leq  | \{ \mathcal{L}f_* \leq \lambda \} | , \qquad \lambda \in \mathbb{R}.
        \]
\end{theorem}
\begin{proof}
    It suffices to apply the above theorem to $\rho(x)=x$.
\end{proof}

\begin{definition}\label{def:polar}
    For a Borel measurable $f:\mathbb{R}^{n}\rightarrow[0,\infty)$ that is %even, 
    convex, and takes the value zero at the origin, define the Polar transform of f to be $$
    f^\circ(x)=\sup_{y \in \mathbb{R}^n}\frac{\langle x,y \rangle-1}{f(y)}
    $$
    when $f$ is not identically zero, and $0^\circ=0$.
\end{definition}
Here, following \cite{artstein-milman}\footnote{see  more precisely \cite[Section 2.1]{artstein-milman}.},  we convey that $\frac{0}{0}=0$ and $\frac{+}{0}=+\infty$. Observe that necessarily $f^\circ(0)=0$.
As is standard in the literature, we denote the set of non-negative 
even 
convex functions on $\mathbb{R}^{n}$ taking the value zero at the origin\footnote{often called geometric convex functions.} by $\mathrm{Cvx}_0(\mathbb{R}^n)$.
%$\Tilde{C}(\mathbb{R}^{n})$.
\begin{theorem} \label{thm: majorization rearrangement Polar}
        For all $f\in\mathrm{Cvx}_0(\mathbb{R}^{n})$ it holds
        \[
            | \{f^\circ \leq \lambda \} | \leq  | \{ (f_*)^\circ \leq \lambda \} | ,
            \qquad \qquad \lambda \in \mathbb{R}.
        \] 
\end{theorem}
\begin{proof}
    The proof is based on the following two observations and an immediate application of Theorem \ref{thm: majorization rearrangement Legendre}. 
    One may observe the following for $f\in\mathrm{Cvx}_0(\mathbb{R}^{n})$
    \begin{enumerate}
        \item $\forall \lambda>0, \{f^\circ>\lambda\}=\lambda\{\mathcal{L}f>\frac{1}{\lambda}\}$;
        \item $\forall t\geq 0, (tA)^{c}=tA^{c}$.
    \end{enumerate} The argument is then as follows (using the homogeneity of the Lebesgue measure at the forth and fifth  line)
    \begin{equation*}
        \begin{split}
            \left|\{(f_{*})^\circ\leq\lambda\}\right|&=\left|\{(f_{*})^{\circ}>\lambda\}^{c}\right|\\
            &\stackrel{1.}{=}\left|\left(\lambda\{Lf_{*}>\frac{1}{\lambda}\}\right)^{c}\right|\\
            &\stackrel{2.}{=}\left|\lambda\{Lf_{*}>\frac{1}{\lambda}\}^{c}\right|\\
            &=\lambda^{n}\left|\{Lf_{*}\leq\frac{1}{\lambda}\}\right|\\
            &\stackrel{Th.\ \ref{thm: majorization rearrangement Legendre}}{\geq}\left|\lambda\{Lf\leq\frac{1}{\lambda}\}\right|\\
            &\stackrel{1.\& 2.}{=}\left|\{f^\circ\leq\lambda\}\right| .
        \end{split}
    \end{equation*}
\end{proof}







\section{Applications} \label{sec:applications}

In this section we collect some applications of our rearrangement results.
The first is about comparison of the solutions of a Hamilton-Jacobi equation with initial data $f$ versus $f_*$. 
In the second subsection, we provide a simple proof of a result by Ben Li \cite{li} about extremizers in functional Blaschke-Santalo type inequalities. Finally, we reduce the very recent semi-group proof by Cordero et al.
\cite{cordero} on functional Blaschke-Santalo by heat-flow, to dimension 1, therefore, simplifying their final argument 
from the use of the variance Brascamp-Lieb inequality to the trivial observation that a square is non-negative. 


\subsection{Hamilton-Jacobi equation and the Hopf-Lax formula}\label{sec:HJ}

In this section, we prove a comparison result for the solution of the following Hamilton-Jacobi equation
\begin{equation} \label{eq:hj}
\begin{cases}
u_t(x) + A(\nabla u_t) = 0 & \mbox{in } \mathbf{R}^n \times (0 ,\infty) \\
u_0=f & \mbox{on } \mathbf{R}^n \times \{t =0\} 
\end{cases}
\end{equation}
with $A \colon \mathbb{R}^n \to \mathbb{R}$ a potential defined as $A(x)=H(\|x\|)$, with $H$ convex, smooth and satisfying $\lim_{x \to \infty} H(x)/x=+\infty$. This is a first order parabolic equation with initial data (at $t=0$) $f$.
Above $u(t,x)$ (for short $u_t(x)$) is the unknown,  function of the two variables $t \in [0,\infty)$ and $x \in \mathbb{R}^n$ and the gradient is acting on the space variable. Here, $\|\cdot\|$ denotes the usual Euclidean norm.

Such PDEs are well understood and we refer to the book by Evans \cite{evans} for an introduction, a historical presentation, and references and to the book by Villani \cite{villani} for connections with the transport of mass phenomenon.
In particular, if $f$ is assumed to be Lipschitz continuous (\textit{i.e.}\ such that $\sup_{\genfrac{}{}{0pt}{}{x,y\in \mathbb{R}^n}{x\neq y}} \frac{|f(x)-f(y)|}{\|x-y\|} < \infty$), the following function 
$$
Q_tf(x) \coloneqq \inf_{y \in \mathbb{R}^n} \left\{ f(y) + t\mathcal{L}H \left(\frac{|x-y|}{t} \right)  \right\}, \qquad t \geq 0
$$
solves the Hamilton-Jacobi equation in the sense that
$Q_tf$ is Lipschitz continuous, differentiable almost everywhere in $\mathbb{R}^n \times(0,\infty)$, it satisfies $u_t(x) + A(\nabla u_t) = 0$ almost everywhere in $\mathbf{R}^n \times (0 ,\infty)$
and $Q_0f=f$ (such a solution is called a generalized solution). Furthermore, it is the unique solution of \eqref{eq:hj} in the viscosity sense. Above $\mathcal{L}H$ is the Legendre transform of $H$ in the sense of Definition \ref{def:legendre}.
The formula that defines $Q_tf$ is known as the Hopf-Lax formula.

Our aim is to compare the solution of \eqref{eq:hj} to the solution of the rearranged problem
\begin{equation} \label{eq:hj-rearranged}
\begin{cases}
v_t(x) + A(\nabla v_t) = 0 & \mbox{in } \mathbf{R}^n \times (0 ,\infty) \\
v_0=f_* & \mbox{on } \mathbf{R}^n \times \{t =0\} 
\end{cases}
\end{equation}
where $f_*$ is the spherically symmetric increasing rearrangement of $f$.

The first occurrence of a comparison problem for Hamilton Jacobi equations apparently appeared in a paper by Ferrone, Posteraro and Volpicelli \cite{FPV} where the potential $H$ is essentially $H(x)=\|x\|$ (see also \cite{giarrusso-nunziante,giarrusso,mercaldo} for related result), followed by 
Alvino, Ferrone, Trombetti and Lions \cite{AFTL} who extended the comparison results of \cite{FPV} to a larger class of $H$ that are assumed to be $1$ homogeneous ($H(tx)=|t|H(x), t \in \mathbb{R}$, $x \in \mathbb{R}^n$) and satisfying $a\|x\| \leq H(x) \leq b\|x\|$ for some constants $a,b >0$. To be really precise, these authors deal with 
Problem \eqref{eq:hj} with a term $A(t,\nabla u_t)$ that satisfies   $A(t,x) \geq H(x)$ with $H(x)=\|x\|$ in \cite{FPV} and $H$ $1$-homogeneous and satisfying $a\|x\| \leq H(x) \leq b\|x\|$ in \cite{AFTL}. More importantly the initial data for the rearranged problem considered in this series of papers is the symmetric decreasing rearrangement $f^*$ rather than $f_*$, and the equation is considered on a general domain $\Omega \subset \mathbb{R}^n$.  Their technique relies on a careful use of the co-area formula together with the classical Euclidean isoperimetric inequality (or some of its variant through more general Brunn-Minkowski's inequality \cite{busemann}, see \cite[Sections 2 and 4.2]{AFTL}).

Here instead, we will exploit the fact the the solution of \eqref{eq:hj} is given by the explicit Hopf-Lax formula to prove a comparison theorem.
In fact, we don't need $H$ to be $1$-homogeneous but rather that it is convex and super linear at infinity
(to guarantee that $Q_tf$ is a generalized solution of Problem \eqref{eq:hj}). Therefore, our results go in a different direction than the above cited references and deal with a different class of potentials $H$.

Before moving to our comparison result for the first order  Hamilton-Jacobi Problem \eqref{eq:hj}, let us briefly mention that the literature is huge on elliptic operators. After the pioneering work of Talenti \cite{talenti}, there have been a continuous and extended production of papers dealing with rearrangement and comparison theorems for solution of elliptic operators (the Laplacian being the model example). Talenti's argument relies, again, as main arguments, on the co-area formula together with the classical Isoperimetric inequality in the Euclidean space. One key observation by Talenti is also that, starting from a rearranged function leads to an equality in the successive inequalities used in the proof of the comparison result, therefore leading to explicit constants in various functional inequalities, with possibly a full description of the extremizers. It is impossible to cite and describe all papers that built upon Talenti's approach as the reader can convince him/herself by consulting the list of more than 200 references published before 2000 contained in \cite{trombetti}\footnote{It is also very hard to cite just a few.}. Let us just mention that people were interested in comparison results for more general elliptic PDE, linear and non linear, degenerate,etc.


\medskip


To make use of our general argument, we consider 
$\mathcal{X}=\mathcal{Y}=\mathbb{R}^n$ equipped with the Lebesgue measure. The Hopf-Lax formula above corresponds to the infimum convolution
$Q_\phi$ given in Definition \ref{def:inf-convolution}
with $\phi (x,y)= t  \mathcal{L}H \left( \frac{\|x-y\|}{t} \right)$, $x,y \in \mathbb{R}^n$, $t>0$ being a fixed parameter. Rearrangement are understood in the usual sense (symmetric decreasing with respect to the Lebesgue measure). Then the usual isoperimetric inequality guarantees that the triple $*$, the Lebesgue measure and $\phi$ is an isoperimetric rearrangement in the sense of Definition  \ref{def:isoperimetric-rearrangement}.
%{\color{red}Cyril: I'm not totally sure of this...}

One of our main theorem is the following.

\begin{theorem} \label{thm:HJ}
Let $H \colon \mathbb{R}_+ \to \mathbb{R}$ be smooth, convex and super-linear ($H(x)/x \to \infty$ when $x \to \infty$). Let $f$ be a Lipschitz continuous function on $\mathbb{R}^n$ and $f_*$ be its symmetric increasing rearrangement.

Denote by $u_t$ and $v_t$ the solution of Problem 
\eqref{eq:hj} and \eqref{eq:hj-rearranged} respectively. Then, for all $\lambda \in \mathbb{R}$, it holds
$$
|\{ v_t < \lambda \}| \leq |\{u_t < \lambda\}| .
$$
%In particular
\end{theorem}

%{\color{red}Cyril: is the conclusion implies the $u_t$ majorizes $v_t$ in the convex order?}

%{\color{red}Cyril: in the paper by Ferone et al, Alvino et al., the comparison is between $u_t^*$ and $v_t$. So to say between $(Q_\phi f)^*$ and $Q_\phi f^*$. Are we missing something here? }


\begin{proof}
From Theorem \ref{th:isoperimetric-rearrangeement} (see \eqref{eq:inf-conv}) we know that
$$
|\{ Q_\phi f_* < \lambda \}| \leq |\{Q_\phi f < \lambda\}| .
$$
The expected result follows at once since by construction $u_t=Q_\phi f$ and $v_t=Q_\phi f_*$ thanks to Lemma \ref{lem:Lipschitz} below (that guarantees that $f_*$ is Lipschitz continuous and therefore that $Q_\phi f_*$ is indeed solution of the rearranged Hamilton-Jacobi equation.
\end{proof}

The proof of Theorem \ref{thm:HJ} relied on the fact that Lipschitz continuity is preserved by rearrangement, a property that is of independent interest and that is most probaly well-known. We give a proof for completeness. Recall that a function on $\mathbb{R}^n$ is $L$-Lipschitz (or Lipschitz continuous with Lipschitz constant $L$) if $\sup_{x \neq y}\frac{|f(x)-f(y)|}{\|x-y\|} \leq L$.

\begin{lemma} \label{lem:Lipschitz}
Let $f \colon \mathbb{R}^n \to \mathbb{R}$ be $L$-Lipschitz with $L \in (0,\infty)$. Then its symmetric increasing rearrangement $f_*$ is also $L$-Lipschitz.
\end{lemma}

\begin{proof}
   We follow some ideas from \cite{MDA25}.
   We claim that a function $f$ is $L$-Lipschitz if and only if for any $\lambda \in \mathbb{R}$, any $\varepsilon>0$, 
   $$
   \{f < \lambda\}_\varepsilon \subset \{f < \lambda +L\varepsilon\},
   $$
   where we recall that $A_\varepsilon=A+\varepsilon B$ with $B$ the Euclidean unit ball. Strict inequality is needed. Indeed assume that $f$ is $L$-Lipschitz. Then if $x \in \{f < \lambda\}_\varepsilon$, by definition, there exists
   $y$ with $f(y)<\lambda$ such that $\|x-y\| < \varepsilon$. In particular,
   $$
   f(x) 
   = 
   f(y) + f(x) - f(y)  
   <
   \lambda + 
   L\|x-y\| 
   <
   \lambda + L \varepsilon  
   $$
which guarantees that $x \in \{f < \lambda +L\varepsilon\}$.

Assume reciprocally that $\{f < \lambda\}_\varepsilon \subset \{f < \lambda +L\varepsilon\}$ and take $x \neq y$. Fix $\delta,\eta>0$. By definition of the enlargement, $x \in \{ f < f(y)+\eta\}_{\|x-y\|+\delta}$ (because $f(y)<f(y)+\eta$). Therefore 
$f(x) < f(y)+\eta + L(\|x-y\|+\delta)$. By symmetry in $x,y$, we conclude that, for any $\delta,\eta>0$ it holds 
$|f(x)-f(y)| < \eta + L(\|x-y\|+\delta)$ from which we get that $f$ is $L$-Lipschitz and therefore the claim.

Now the key observation is given in \eqref{eq: increasing rearrangement characterization} as $\{f_* < \lambda\}=\{f<\lambda\}^*$. Fix $\lambda \in \mathbb{R}$, we will prove that 
$\{f_* < \lambda\}_\varepsilon \subset \{f_* < \lambda +L\varepsilon\}=\{f < \lambda +L\varepsilon\}^*$.
Observe that the set 
$$
\{f_* < \lambda\}_\varepsilon =
\{f_* < \lambda\} + \varepsilon B
=
\{f < \lambda\}^* + \varepsilon B
$$ 
is a centered ball. Since  $\{f < \lambda +L\varepsilon\}^*$ is also a centered ball, the desired inclusion 
$\{f_* < \lambda\}_\varepsilon \subset \{f < \lambda +L\varepsilon\}^*$ will follow by comparison of the volume of these two balls. Now, by the equality case in the Brunn-Minkowski's inequality and then Brunn-Minkowski's inequality (recall that $|\cdot|$ denotes the volume with respect to the Lebesgue measure), it holds
\begin{align*}
|\{ f_* < \lambda\}_\varepsilon \}|^\frac{1}{n} 
& = |\{ f < \lambda\}^* + \varepsilon B |^\frac{1}{n} \\\
& = 
|\{ f < \lambda\}^*|^\frac{1}{n} + \varepsilon |B|^\frac{1}{n} & \mbox{(equality case in Brunn-Minkoswki)} \\
& 
=
|\{ f < \lambda\}|^\frac{1}{n} + \varepsilon |B|^\frac{1}{n}  & \mbox{(by construction of the rearrangement)} \\
& \leq 
|\{ f < \lambda\}+ \varepsilon B|^\frac{1}{n} & \mbox{(Brunn-Minkoswki)} \\
&= 
|\{ f < \lambda\}_\varepsilon|^\frac{1}{n} \\
& \leq 
|\{ f < \lambda+L\varepsilon\}|^\frac{1}{n}
& \mbox{(Since } f \mbox{ is } L-\mbox{Lipschitz)} \\
& =
|\{ f < \lambda+L\varepsilon\}^*|^\frac{1}{n} & \mbox{(by construction of the rearrangement)} \\
& =
|\{ f_* < \lambda+L\varepsilon\}|^\frac{1}{n}
& \mbox{(by \eqref{eq: increasing rearrangement characterization})}.
\end{align*}
The volume are therefore comparable and hence $\{f_* < \lambda\}_\varepsilon \subset \{f_* < \lambda +L\varepsilon\}$ as announced, leading to the desired conclusion of the lemma. 
% {\color{red}Devraj: Actually I really didn't have anything interesting to say but I was solely keen on asking a question. I know that initially we had spoken on talking about a comparison on (12) and (13) where the space variable x belongs to the sphere $S^{n-1}$. Am i correct in assuming that we have to restrict our discussion to $\mathbb{R}^n$ as lemma 5.2 may not necessarily hold for f being L-Lipschitz on the sphere? It seems like the proof requires some Brunn-Minkowski inequality which I dont think we have on the sphere..}
\end{proof}




\subsection{Extremal Functions in the Functional Santalo Inequality}

In all the following we will use $\mathcal{A}$ to denote either the Legendre transform ($\mathcal{L}$, Definition \ref{def:legendre}) or the Polar transform ($^\circ$, Definition \ref{def:polar}). Our aim is to recover, in a shorter way, the following result of Ben Li, using rearrangement. Recall that we denote $\mathrm{Cvx}_0(\mathbb{R}^{n})$ the set of non-negative even convex functions on $\mathbb{R}^{n}$ taking the value zero at at the origin. 
Let $\omega_{n}=\pi^{n/2}/\Gamma(\frac{n}{2}+1)$ denote the volume of the $n$ dimensional Euclidean unit ball.

\begin{theorem}[\cite{li}] \label{thm: existence of extremizer}
    There exists a radially symmetric extremizer of the following variational problem 
    $$
    \sup_{f\in\mathrm{Cvx}_0(\mathbb{R}^{n})}\int e^{-f}\int e^{-\mathcal{A}f}
    $$ 
    where the integrals are implicitly assumed to be taken with respect to the Lebesgue measure on $\mathbb{R}^{n}$. 
% \begin{enumerate}
%     \item For $\mathcal{A}=\mathcal{L}$, $\Lambda$ is the class of non-negative convex functions with f(0)=0.
%     \item For $\mathcal{A}=\circ$,  $\Lambda = \mathrm{Cvx}_0(\mathbb{R}^{n})$.
% \end{enumerate}
In particular
$$
    \sup_{f\in \mathrm{Cvx}_0(\mathbb{R}^{n})} 
    \int e^{-f}\int e^{-\mathcal{L}f}
    =
    (n\omega_{n})^2 
    \sup_{\Psi} 
    \int_0^\infty e^{-\Psi(r)}r^{n-1}dr 
    \int_0^\infty e^{-\mathcal{L}\Psi}r^{n-1}dr 
$$
and
$$
    \sup_{f\in \mathrm{Cvx}_0(\mathbb{R}^{n})} 
    \int e^{-f}\int e^{-f^\circ}
    =
    (n\omega_{n})^2 
    \sup_{\Psi} 
    \int_0^\infty e^{-\Psi(r)}r^{n-1}dr 
    \int_0^\infty e^{-\Psi^\circ(r)}r^{n-1}dr
$$
where the supremum runs over all convex non-decreasing function $\Psi \colon [0,\infty) \to [0,\infty)$ with $\Psi(0)=0$, and 
$\mathcal{L}\Psi(r)=\sup_{s>0}\{rs-\Psi(s)\}$, 
$\Psi^\circ(r) \coloneqq \sup_{s>0} \frac{rs-1}{\Psi(s)}$. 
\end{theorem} 

The theorem is a direct consequence of a series of elementary lemmas.

Let $\mathcal{M}$ be the set of absolutely continuous measures on $[0,\infty)$ (absolutely continuous with respect to the Lebesgue measure) whose density is log-concave. Furthermore, we define the class $\mathcal{S}$ in the following manner 
$$
\mathcal{S} :=\{\mu \in \mathcal{M} :  \frac{d\mu}{dx}=e^{-\Psi(x)}
\mbox{ with } \Psi(0)=0, \Psi \mbox{ is non-decreasing  and } \int_{0}^{\infty}\!e^{-\Psi(r)}r^{n-1}dr=1\} .
$$
Note that log-concave just says that $\Psi \colon [0,\infty) \to \mathbb{R}$ is convex. We need some preparation. 
Since $\Psi$ is non-decreasing, it holds  
    $$
    1
    =
    \int_{0}^{\infty}e^{-\Psi(r)}r^{n-1}dr\geq\int_{0}^{2}e^{-\Psi(r)}r^{n-1}dr
    \geq 
    e^{-\Psi(2)}\frac{2^{n}}{n} .
    $$ 
    In particular $\Psi(2) \geq \log\left(\frac{2^{n}}{n}\right)$. Set 
    \begin{equation} \label{eq:def_a}
      a \coloneqq 
      \frac{1}{2}\log\left(\frac{2^{n}}{n}\right)    
    \end{equation}
    and notice that $a>0$ does not depend on $\Psi$. Now, observe that since $\Psi$ is convex and $\Psi(0)=0$, the map $(0,\infty) \ni r \mapsto \Psi(r)/r$ is non-decreasing so that, for any $r \geq 2$, (again, uniformly in $\Psi$)
    \begin{equation} \label{eq:psi}
    \Psi(r) \geq \frac{\Psi(2)}{2}r \geq ar .
    \end{equation}


\begin{lemma} \label{lem:tight}
    The set $\mathcal{S}$ is tight.
\end{lemma}
\begin{proof}
    To prove tightness we need to show that, for all 
    $\varepsilon>0$, there exists a compact set $K_{\varepsilon} \subset [0,\infty)$ such that $\mu([0,\infty)\setminus K_{\varepsilon})<\varepsilon$ for all $\mu\in\mathcal{S}$.  
    
    Fix $\varepsilon >0$ and take $K_{\varepsilon}=[0,b]$ with $b \geq 2$ chosen so that $\frac{e^{-ab}}{a} < \varepsilon$ (and $a$ defined in \eqref{eq:def_a}). Then, 
    %for $\varepsilon \in (0,1/2]$ 
    it holds
    $$
    \mu([0,\infty)\setminus K_{\varepsilon})
    =
    \int_{b}^{\infty}e^{-\Psi(r)}dr
    \leq
    \int_{b}^{\infty}e^{-ar}dr
    =
    \frac{e^{-ab}}{a} 
    <
    \epsilon 
    $$
    as expected.
\end{proof} 

\begin{lemma}\label{lem:absolutely-continuous}
If $\mu_{n}\in\mathcal{S}$ converges weakly to a measure $\mu$,  then $\mu$ is absolutely continuous with respect to the Lebesgue measure. Moreover, if $e^{-\Psi_n}$ denotes the density of $\mu_n$ and $e^{-\Psi}$ the density of $\mu$, then $\Psi_n$ converges point-wise to $\Psi$ on the positive axis.
\end{lemma}

\begin{proof}
Denote by $e^{-\Psi_n}$ the density of $\mu_n$ with respect to the Lebesgue measure,
on $[0,\infty)$. Since $\mu_n \in \mathcal{S}$ we have $\mu_n(U) = \int_U e^{-\Psi_n(x)} dx \leq |U| e^{- \Psi_n(0)} = |U|$ for any Borel set $U$.  Thus for an open set $U$ such that $A \subseteq U$, by the weak convergence,
\begin{align*}
    \mu(A) \leq \mu(U) \leq \liminf_n \mu_n(U) \leq \liminf_n |U| = |U|.
\end{align*}
By the regularity of the Lebesgue measure, taking the infimum over all such $U$ gives,
\[
    \mu(A) \leq |A|.
\]
Thus $\mu$ has a density with respect to the Lebesgue measure, as expected. 

We will now show that $\lim_{n\rightarrow\infty}\Psi_{n}(x)=\Psi(x)$ for all $x \in [0,\infty)$.
By Helly's selection theorem we can extract from the sequence $(e^{-\Psi_n})$ a subsequence $(e^{-\Psi_{\varphi(n)}})$ converging point-wise to $e^{-\Psi}$.
In particular $\Psi$ is convex and therefore continuous on $(0,\infty)$.  
%Now every subsequence of $(\Psi_n(x))$ must converge to $\Psi(x)$, so that the whole sequence converge to $\Psi(x)$ and the proof is complete.



Remark that for all $a,b>0$, weak convergence implies that $\lim_{n\rightarrow\infty}\mu_{n}(a,b)=\mu(a,b)$. By contradiction, assume that $\limsup_{n\rightarrow\infty}e^{-\Psi_{n}(r)}>e^{-\Psi(r)}$ for some $r > 0$. 
The strict inequality guarantees  the existence of $\epsilon>0$ such that $\limsup_{n\rightarrow\infty}e^{-\Psi_{n}(r)} > \varepsilon + e^{-\Psi(r)}$. By monotonicity of $\Psi_n$, for any $0<a<r$ it holds
     \begin{align*}
      (r-a)(\varepsilon + e^{-\Psi(r)})
      < (r-a) \limsup_{n\rightarrow\infty} e^{-\Psi_n(r)} 
      \leq 
      \limsup_{n\rightarrow\infty}\mu_{n}(a,r)
      =\mu(a,r) 
      \leq (r-a)e^{-\Psi(a)} .
      \end{align*} 
      Therefore $\varepsilon + e^{-\Psi(r)} < e^{-\Psi(a)}$ that leads, by continuity of $\Psi$, in the limit $a \to r$, to a contradiction. Hence,  for all $r >0$, $\limsup_{n\rightarrow\infty}e^{-\Psi_{n}(r)} \leq e^{-\Psi(r)}$.
     
     Assume now that  $\liminf_{n\rightarrow\infty}e^{-\Psi_{n}(r)}<e^{-\Psi(r)}$ for some $r > 0$. Similarly, for some $\varepsilon>0$ and any $a>r$ it holds
  \begin{align*}
      (a-r)(-\varepsilon + e^{-\Psi(r)})
      > (a-r) \liminf_{n\rightarrow\infty} e^{-\Psi_n(r)} 
      \geq 
      \limsup_{n\rightarrow\infty}\mu_{n}(r,a)
      =\mu(r,a) 
      \geq (a-r)e^{-\Psi(a)} .
      \end{align*}    
This leads in the limit $a \to r$ to a contradiction and therefore to the desired conclusion: $\Psi_n(x) \to \Psi$ point-wise on $(0,\infty)$ and hence on $[0,\infty)$ since $\Psi_n(0)=\Psi(0)=0$.



% The strict inequality together with the continuity of $\Psi$ guarantee  the existence of $\epsilon,\delta>0$ such that $\limsup_{n\rightarrow\infty}\Psi_{n}(r)-\varepsilon>\Psi(r+\delta)$. By monotonicity of $\Psi_n$
%      \begin{align*}
%               \limsup_{n\rightarrow\infty}\mu_{n}(r,r+\delta)&\geq\delta\limsup_{n\rightarrow\infty}\Psi_{n}(r)\\
%              &>\delta\varepsilon+\delta\Psi(r+\delta)\\
%              &\geq\mu(r,r+\delta)
%       \end{align*} 
%       {\color{red}Cyril: the above computation is not correct.}
%      The case $\liminf_{n\rightarrow\infty}\Psi_{n}(r)<\Psi(r)$ may be treated analogously. 


%(but this density can be taken to be )bounded by $1$. 
%Curiously, we didn't use log-concavity like I suggested originally in our discussions!


% This is actually a general fact, if $\alpha$ is a fixed Radon measure, the operator $M_\alpha(\mu) = \left\| \frac{d\mu}{d\alpha} \right\|_\infty$ where the essential suprema is taken with respect to $\alpha$ and defined to be infinity when $\mu$ is not absolutely continuous with respect to $\alpha$, is lower semi-continuous by the same argument.  Even more generally, if  I recall correctly and the space of action is Polish, this can be taken as a statement about the lower semi-continuity of the $p$-R\'enyi divergence for $p > 0$ (this is the $p = \infty$ case above), or more generally still, any $f$-divergence being jointly lower semicontinuous (in both arguments simultaneously).

% Now denote $\frac{d\mu}{dx}=e^{-\Psi(x)}$. We will now show that $\lim_{n\rightarrow\infty}\Psi_{n}(x)=\Psi(x)$ for all $x \in [0,\infty)$.
% First remark that for $\forall a,b>0$, weak convergence implies that $\lim_{n\rightarrow\infty}\mu_{n}(a,b)=\mu(a,b)$. By contradiction, assume that $\limsup_{n\rightarrow\infty}\Psi_{n}(r)>\Psi(r)$ for some $r\geq 0$. The strict inequality guarantees  the existence of $\epsilon,\delta>0$ such that $\limsup_{n\rightarrow\infty}\Psi_{n}(r)-\varepsilon>\Psi(r+\delta)$. Observe the following
%     \begin{equation*}
%         \begin{split}
%             \limsup_{n\rightarrow\infty}\mu_{n}(r,r+\delta)&\geq\delta\limsup_{n\rightarrow\infty}\Psi_{n}(r)\\
%             &>\delta\varepsilon+\delta\Psi(r+\delta)\\
%             &\geq\mu(r,r+\delta)\\
%         \end{split}
%     \end{equation*} The case $\liminf_{n\rightarrow\infty}\Psi_{n}(r)<\Psi(r)$ may be treated analogously. 
\end{proof}


\begin{lemma} \label{lem:compact}
    The set $\mathcal{S}$ is compact  for the topology of the weak convergence.
\end{lemma}


\begin{proof}
    Consider a sequence $(\mu_n)$ of $\mathcal{S}$. We will prove that we can extract a converging subsequence (for the weak convergence) to a limit $\mu \in \mathcal{S}$.
     Denote by $e^{-\Psi_n}$ the density of $\mu_n$ with respect to the Lebesgue measure,
on $[0,\infty)$. Since $\mu_n \in \mathcal{S}$, $e^{-\Psi_n} \in [0,1]$ is non-increasing. Applying Helly's selection theorem we can extract a subsequence
$e^{-\Psi_{\varphi(n)}}$ converging pointwise to some limit $e^{-\Psi}$. Now by Lemma \ref{lem:tight} and Prokhorov's theorem we can further extract from the sub-sequence $(\mu_{\varphi(n)})$ a sub-subsequence weakly converging to some measure $\mu$ (with density $e^{-\Psi}$ by Lemma \ref{lem:absolutely-continuous}). Let's call for simplicity again $\mu_n$ this converging sub-subsequence.

To prove that $\mathcal{S}$ is compact, it only remains to prove that the limit $\mu$ belongs to $\mathcal{S}$. 

By the point-wise convergence we infer that $\Psi(0)=0$ and that $\Psi$ is non-decreasing and convex. We therefore only need to check that
$$
\int_{0}^{\infty}e^{-\Psi(r)}r^{n-1}dr=1.
$$
Recall the definition of $a$ defined in \eqref{eq:def_a}
and set $g(r)=r^{n-1} \left(\mathds{1}_{[0,2]} + e^{-ar}\mathds{1}_{(2,\infty)}\right)$. From \eqref{eq:psi} we see that $\Psi_m \leq g$ for all $m$, and since $g$ is integrable on the positive axis, by the dominated convergence theorem, we finally obtain
$$
1 = \lim_{m \to \infty} \int_0^\infty e^{-\Psi_m(r)}r^{n-1}dr = \int_0^\infty e^{-\Psi(r)}r^{n-1}dr
$$
as expected. This ends the proof of the Lemma.
\end{proof}
    
We are now in position to proove Theorem \ref{thm: existence of extremizer}.

\begin{proof} [Proof of Theorem \ref{thm: existence of extremizer}]
In the variational problem $\sup_{f\in\mathrm{Cvx}_0(\mathbb{R}^{n})}\int e^{-f(x)}dx\int e^{-\mathcal{A}f(x)}dx$ under consideration, $e^{-f}$ and $e^{-\mathcal{A}f}$ are non-negative functions. Therefore, we use the layer cake decomposition to get
\begin{align*}
\int e^{-f(x)}dx\int e^{-\mathcal{A}f(x)}dx
& =
\int_0^\infty |\{e^{-f} > t\}|dt \int_0^\infty |\{e^{-\mathcal{A}f} > t\}|dt \\
& =
\int_0^\infty |\{f < -\log t\}|dt \int_0^\infty |\{\mathcal{A}f < -\log t\}|dt .
\end{align*}
Now by \eqref{eq: increasing rearrangement characterization} and the definition of the rearrangement, it holds
$$
|\{f < -\log t\}|=|\{f < -\log t\}^*|=|\{f_* < -\log t\}|
$$ 
for any $t>0$. Therefore, using Theorem \ref{thm: majorization rearrangement Legendre}/Theorem \ref{thm: majorization rearrangement Polar} (according to the transform $\mathcal{A}$ under consideration), we get
$$
\int e^{-f(x)}dx\int e^{-\mathcal{A}f(x)}dx
\leq 
\int_0^\infty |\{f_* < -\log t\}|dt \int_0^\infty |\{\mathcal{A}f_* < -\log t\}|dt .
$$
Using the layer cake decomposition backward it follows that
$$
\int e^{-f(x)}dx\int e^{-\mathcal{A}f(x)}dx
\leq 
\int e^{-f_*(x)}dx\int e^{-\mathcal{A}f_*(x)}dx .
$$
Hence, as a first intermediate conclusion, since $f_* \in \mathrm{Cvx}_0(\mathbb{R}^{n})$,
%since $\Lambda$ is stable under the action of the Legendre/polar transform, 
it holds
$$
\sup_{f \in \mathrm{Cvx}_0(\mathbb{R}^{n})} \int e^{-f(x)}dx\int e^{-\mathcal{A}f(x)}dx
=
\sup_{f \in \mathrm{Cvx}_0(\mathbb{R}^{n})} \int e^{-f_*(x)}dx\int e^{-\mathcal{A}f_*(x)}dx .
$$
Now we observe that for any $f \in \mathrm{Cvx}_0(\mathbb{R}^{n})$, there exists a convex non-decreasing function $\Psi \colon [0,\infty) \to [0,\infty)$, with $\Psi(0)=0$ such that $f_*(x)=\Psi(\|x\|)$, $x \in \mathbb{R}^d$. Moreover, it is not difficult to check that 
$\mathcal{A}f_*(x)=\mathcal{A}\Psi(\|x\|)$ for any $x \in \mathbb{R}^d$. Therefore, passing to polar coordinates, we end up with
$$
\sup_{f \in \mathrm{Cvx}_0(\mathbb{R}^{n})} \int e^{-f(x)}dx\int e^{-\mathcal{A}f(x)}dx
= (n\omega_n)^2
\sup_{\Psi} \int_0^\infty e^{-\Psi(r)}r^{n-1}dx
\int_0^\infty e^{-\mathcal{A}\Psi(r)}r^{n-1}dr 
$$
where the supremum is running over all $\Psi$ convex, non-decreasing and satisfying $\Psi(0)=0$, and where $\omega_n\coloneqq\pi^{n/2}/\Gamma(\frac{n}{1}+1)$ denotes the volume of the $n$ dimensional Euclidean unit ball. 

We also observe that the left and right hand side are invariant under scaling, \textit{i.e.}\ changing $f$ into $x \mapsto f(\lambda x)$, and similarly for $\Psi$ does not affect the products $\int e^{-f(x)}dx\int e^{-\mathcal{A}f(x)}dx$ and $\int_0^\infty e^{-\Psi(r)}r^{n-1}dx\int_0^\infty e^{-\mathcal{A}\Psi(r)}r^{n-1}dr$. Therefore, we can assume $\int e^{-\Psi(r)}r^{n-1}dr=1$ and 
$$
\sup_{f \in \mathrm{Cvx}_0(\mathbb{R}^{n})} \int e^{-f(x)}dx\int e^{-\mathcal{A}f(x)}dx
= (n\omega_n)^2 \sup_{\mu \in \mathcal{S}} \int_0^\infty e^{-\mathcal{A}\Psi(r)}r^{n-1}dr .
$$
To conclude the proof of the lemma we claim that the map $\mathcal{S} \ni \mu \mapsto \int_0^\infty e^{-\mathcal{A}\Psi(r)}r^{n-1}dr$  is upper semi continuous. Indeed, consider a weakly converging sequence $(\mu_m)$ of $\mathcal{S}$. By compactness of $\mathcal{S}$ (Lemma \ref{lem:compact}), $\mu \in \mathcal{S}$, and by Lemma \ref{lem:absolutely-continuous},  if $e^{-\Psi_m}$ denote the density of $\mu_m$ and $e^{-\Psi}$ the density of $\mu$, $\Psi_m$ converge point-wise to $\Psi$.
%$\Psi_{n}\rightarrow\Psi$ in a pointwise fashion. 
By Fatou's Lemma, it follows
$$
\limsup_{m\rightarrow\infty}\int_{0}^{\infty}e^{-\mathcal{A}\Psi_{m}(r)}r^{n-1}dr\leq\int_{0}^{\infty}\limsup_{m\rightarrow\infty}e^{-\mathcal{A}\Psi_{m}(r)}r^{n-1}dr
\leq 
\int_{0}^{\infty}e^{-\mathcal{A}\Psi(r)}r^{n-1}dr,
$$ 
where in the last inequality we used that $\liminf \mathcal{A}\Psi_m \geq \mathcal{A}\Psi$, a property that can easily be checked from the definition of the Legendre/polar transform, proving the claim. 
Since upper semi-continuous functions attains their maximum on compact sets, the proof of the theorem is complete.

% Consider a sequence $(\mu_n)$ of $\mathcal{S}$ approximating the supremum in the right hand side. since $\mathcal{S}$ is compact we can extract a subsequence converging to some $\mu \in \mathcal{S}$, with density $\Psi$, realizing the supremum. This achieves the proof of the theorem.
% ith the fact that the transform $\mathcal{A}$ preserves rotational invariance, we see that $$
% \sup_{f_\in\Lambda}\int e^{-f}\int e^{-\mathcal{A}f}\leq\sup_{f_*\in\Lambda}\int e^{-f_*}\int e^{-\mathcal{A}f_*}
% $$ Observe that $$
% \sup_{f_*\in\Lambda}\int e^{-f_*}\int e^{-\mathcal{A}f_*}=\sup_{f_*\in\Lambda,\int e^{-f_{*}}=\frac{1}{|S_{n-1}|}}\int e^{-f_*}\int e^{-\mathcal{A}f_*}
% $$ 
% It suffices to consider $f_{\lambda}(x)=f(\lambda x)$ for $\lambda=\left(\frac{1}{|S_{n-1}|\int e^{-f}}\right)^{\frac{1}{n}}$. \\\\The above deduction justifies the following reduction 
% $$
% \sup_{f\in\Lambda}\int e^{-f}\int e^{-\mathcal{A}f}\leq|S_{n-1}|\sup_{\mu\in\mathcal{S}}\int_{0}^{\infty}e^{-\mathcal{A}\Psi(r)}r^{d-1}dr
% $$ 
% To conclude the proof of the lemma we claim that the map $\mathcal{S} \ni \mu \mapsto \int_0^\infty e^{-\mathcal{A}\Psi(r)}r^{d-1}dr$  is upper semi continuous. Indeed, consider $\mu_{n}\rightarrow\mu$ in the weak topology. By an above lemma, $\Psi_{n}\rightarrow\Psi$ in a pointwise fashion. Let us now enjoy a simple application of Fatou's Lemma, $$
% \limsup_{n\rightarrow\infty}\int_{0}^{\infty}e^{-\mathcal{A}\Psi_{n}(r)}r^{d-1}dr\leq\int_{0}^{\infty}\limsup_{n\rightarrow\infty}e^{-\mathcal{A}\Psi_{n}(r)}r^{d-1}dr=\int_{0}^{\infty}e^{-\mathcal{A}\Psi(r)}r^{d-1}dr
% $$ We conclude the proof by virtue of the elementary fact that upper semi continuous functions obtain attain their maximum defined on a compact set. 
\end{proof}




% \subsection{Legendre Transform}
% \begin{definition} \label{def:legendre}
%     For a Borel measurable $f : \mathbb{R}^n \to \mathbb{R}$ define the Legendre transform of $f$,  $\mathcal{L}f:\mathbb{R}^n \to \mathbb{R}$
%     \[
%         \mathcal{L}f(x) = \sup_{y \in \mathbb{R}^n} \bigg[ \langle x , y  \rangle -  f(y) \bigg].
%     \]
% \end{definition}
 
% \begin{theorem} \label{thm: majorization rearrangement Legendre}
%         If $f: \mathbb{R}^n \to [0,\infty)$ and $f(0) = 0$, then
%         \[
%             | \{ \mathcal{L}f \leq \lambda \} | \leq  | \{ \mathcal{L}f_* \leq \lambda \} |
%         \]
% \end{theorem}

% The proof of the theorem will be based on the following fundamental result from convex geometry, for which we will need the following definition,
% \begin{definition}
%     For a Borel set $K$, define $h_K(x) = \sup_{y \in K} \langle x, y \rangle$, and
%     \[
%         K^\circ \coloneqq \{ x: h_K(x) \leq 1 \}
%     \]
%     when $K$ is non-empty, and $K^\circ = \mathbb{R}^n$ otherwise.
% \end{definition}

% \begin{theorem}[Blaschke-Santalo]
%         For a non-empty Borel set $K$,
%         \[
%             |K||K^\circ| \leq |B_n |^2
%         \]
% \end{theorem}
% Note that since, $(\lambda K)^\circ = \lambda^{-1} K^\circ$, $| \lambda K | |(\lambda K)^\circ | = |K||K^\circ|$, and the Blaschke-Santalo inequality can be reformulated as a rearrangement inequality.


% \begin{theorem}[Rearrangement Blaschke-Santalo]  \label{thm: Blaschke-Santalo rearrangement}
% For a Borel set $K \subseteq \mathbb{R}^n$,
% \begin{align*}
%     |K^\circ| \leq | (K^*)^\circ|.
% \end{align*}
% \end{theorem}

% \begin{proof}[Proof of Theorem \ref{thm: majorization rearrangement Legendre}]
%     Let us first note that since $f$ takes its minimum $0$ at $0$, 
%         \[
%             \mathcal{L}f(x) = \sup_y \big[ \langle x,y\rangle - f(y) \big] \geq \langle x ,0 \rangle - f(0) = 0.
%         \]
%     Thus for $\lambda < 0$, $\{ \mathcal{L} f \leq \lambda \}$ is empty and there is nothing to prove.
%    Thus we assume that $\lambda \geq 0$. Then with the cost function $\phi(x,y) = - \langle x,y \rangle$,
%     \begin{align*}
%         \{ \mathcal{L}f \leq \lambda \} 
%             &=
%                 \{ Q_\phi f \geq - \lambda \}
%                     \\
%             &=
%                 \{ Q_\phi f < - \lambda \}^c
%                     \\
%             &=
%                \left( \bigcup_{q \in S(-\lambda)} \{ f < q_1 \}_{\phi(q_2)} \right)^c
%                     \\
%             &=
%                 \bigcap_{q \in S(-\lambda)} \{ f < q_1 \}_{\phi(q_2)}^c
%     \end{align*}
%     where $S(-\lambda) = \{q = (q_1,q_2) \in \mathbb{Q}^2 :  q_1 + q_2 < - \lambda \}$.  Note that since $\{ f < q_1 \}$, is empty unless $q_1 > 0$, it suffices to consider the union over $q \in S_+(-\lambda)$, where $S_+(-\lambda) = S(-\lambda) \bigcap \{x: x_1 > 0 \}$.  Summarizing, we have
%     \begin{align} \label{eq: exp Legendre transform sublevel set}
%         \{ \mathcal{L} f \leq \lambda \} = \bigcap_{q \in S_+(-\lambda)} \{ f < q_1 \}_{\phi(q_2)}^c
%     \end{align}
%     Note that for $q \in S_+(-\lambda)$, since $q_1 \geq 0$ and $q_1 + q_2 < -\lambda$, $q_2$ is necessarily negative.
%     We claim that for $\varepsilon >0$, and a Borel set $K$,
%      \begin{align} \label{eq: enlargement with dot product cost}
%         (K_{\phi(-\varepsilon)})^c = \varepsilon \ K^\circ.
%     \end{align}
%     Indeed,
%     \begin{align*}
%         (K_{\phi(-\varepsilon)})^c
%             &=
%                 \{ x : \exists y \in K, \phi(x,y) < - \varepsilon \}^c
%                     \\
%             &=
%                 \left\{ x: \inf_{y \in K} - \langle x,y \rangle < - \varepsilon \right\}^c
%                     \\
%             &=
%                 \left\{ x : \sup_{y \in K } \langle x,y \rangle > \varepsilon \right\}^c
%                     \\
%             &=
%                 \left\{ x : \sup_{y \in K} \langle x,y \rangle \leq \varepsilon \right\}
%                     \\
%             &=
%                 \varepsilon \left\{ x : \sup_{y \in K} \langle x,y \rangle \leq 1 \right\}
%                     \\
%             &=
%                 \varepsilon \ K^\circ.
%     \end{align*}
%     Applying \eqref{eq: enlargement with dot product cost} to \eqref{eq: exp Legendre transform sublevel set}, we have 
%     \begin{align*}
%         \{ \mathcal{L} f \leq \lambda \} = \bigcap_{q \in S_+(-\lambda)} |q_2| \{ f < q_1 \}^\circ,
%     \end{align*}
%     Thus,
%     \begin{align*}
%         |\{ \mathcal{L} f \leq \lambda \}| 
%             &=
%                 \left|  \bigcap_{q \in S_+(-\lambda)} |q_2| \{ f < q_1 \}^\circ \right|
%                     \\
%             &\leq 
%                 \inf_{q \in S_+(-\lambda)} |q_2|^n \ \big| \{f < q_1 \}^\circ \big|
%                     \\
%             &\leq 
%                 \inf_{q \in S_+(- \lambda)} |q_2|^n \ \big| (\{ f < q_1 \}^*)^\circ \big |
%                     \\
%             &= 
%                 \inf_{q \in S_+(- \lambda)}  \big| |q_2| \ (\{ f < q_1 \}^*)^\circ \big |
%                     \\
%             &=
%                 \left|\bigcap_{q \in S_+(- \lambda)} |q_2|  \ (\{ f < q_1 \}^*)^\circ \right|
%                     \\
%             &=
%                 \left|\bigcap_{q \in S_+(- \lambda)} |q_2| \ (\{ f_* < q_1 \})^\circ \right|
%                     \\
%             &=
%                 | \{ \mathcal{L} f_* \leq \lambda \} |
%     \end{align*}
%     where the first inequality is $\mu( \cap_n A_n) \leq \inf_n \mu(A_n)$ combined with an application of the homogeneity of the Lebesgue measure.  The second inequality is Theorem \ref{thm: Blaschke-Santalo rearrangement}, the following equality is again the homogeneity of the Lebesgue measure, while the following equality is due to the fact that $|q_2| (\{ f > q_1 \}^*)^\circ$ is always an origin symmetric ball (potentially of infinite radius) so that the we are considering an intersection of nested sets, the 2nd to last equality is the characterization of the spherically symmetric increasing rearrangement, an the final inequality is the description of the sub-level sets of the Legendre transform applied to the spherically symmetric increasing rearrangement.
% \end{proof}

%Theorem \ref{thm: existence of extremizer} admits the following corollary.

%\begin{corollary} \label{cor:extremizers}
% It holds
% $$
%     \sup_{f\in \tilde C(\mathbb{R}^n)} 
%     \int e^{-f}\int e^{-f^\circ}
%     =
%     (n\omega_{n})^2 
%     \sup_{\Psi} 
%     \int_0^\infty e^{-\Psi(r)}r^{n-1}dr 
%     \int_0^\infty e^{-\Psi^\circ(r)}r^{n-1}dr
% $$
% where the supremum runs over all convex non-decreasing function $\Psi \colon [0,\infty) \to [0,\infty)$ with $\Psi(0)=0$, and 
% $\Psi^\circ(r) \coloneqq \sup_{s>0} \frac{rs-1}{\Psi(s)}$. Here $\omega_{n}=\pi^{n/2}/\Gamma(\frac{n}{2}+1)$ denotes the volume of the $n$ dimensional Euclidean unit ball.
% \end{corollary}

% {\color{red}Cyril: observe that the same corollary holds for the Legendre transform in the sense that
% $$
%     \sup_{f\in \tilde C(\mathbb{R}^n)} 
%     \int e^{-f}\int e^{-\mathcal{L}f}
%     =
%     (n\omega_{n})^2 
%     \sup_{\Psi} 
%     \int_0^\infty e^{-\Psi(r)}r^{n-1}dr 
%     \int_0^\infty e^{-\mathcal{L}\Psi}r^{n-1}dr 
% $$
% with the same family of $\Psi$ (convex, non-decreasing with $\Psi(0)=0$). With this formulation one would expect to be able to use the semi-group approach of Cordero et al (that should be simpler in dimension 1). 
% %But I don't see how to handle the $r^{n-1}$ extra term coming from the change of variables in polar coordinate.
% }

% \begin{proof}
% Any radially symmetric function $f$ of $\tilde C(\mathbb{R}^n)$ can be written as $f(x)=\Psi(|x|)$ for some convex non-decreasing function $\Psi \colon ([0,\infty) \to [0,\infty)$ with $\Psi(0)=0$.
% Then it is easy to see that $f^\circ(x)=\Psi^\circ(|x|)$, $x \in \mathbb{R}^n$, so that the expected result follows from Theorem \ref{thm: existence of extremizer} passing to polar coordinates.
% \end{proof}



% Given a Lagrangian $f: \mathbb{R} \to [0,\infty)$ such that $f(0) = 0$, then the partition function $Z: (0,\infty) \to (0,\infty)$ define by
% \[
%     Z(\beta) \coloneqq  \int_{\mathbb{R}^n} e^{-\beta \mathcal{L} f(x)} dx 
% \]
% If we define $f^{*}$ as the symmetric increasing rearrangement of $f$, and $$Z_*(\beta)  \coloneqq  \int_{\mathbb{R}^n} e^{-\beta \mathcal{L} f_*(x)} dx $$ we have the following corollary.

% \begin{corollary}
%     For a Lagrangian $f$, and $f_*$ its symmetric increasing rearrangement, then we have a pointwise comparison of their respective partition functions $Z$ and $Z_*$,
%     \[
%         Z(\beta) \leq Z_*(\beta)
%     \]
% \end{corollary}

% \begin{proof}
    
% \end{proof}
% \section{Distinct Spaces using (Approximate) Isoperimetry}
% \subsection{Strictly Log-Concave Measures in $\mathbb{R}^{n}$}
% Let us consider the class of measures in $\mathbb{R}^{n}$ admitting the following density $$\frac{d\mu}{dx}=\frac{1}{\mathcal{Z}}e^{-W(x)}, V''(x)\geq cId>0$$ Denote the following quantities as such $$\Phi(a)=\int_{-\infty}^{a}\frac{1}{\sqrt{2\pi}}e^{-\frac{x^{2}}{2}}dx, \varphi=\Phi', \gamma(A)=\int_{A}\varphi(x)dx$$ It is well known that $\mu$ satisfies the following isoperimetric inequality \begin{equation}
%     \mu^+(A)\geq\sqrt{c}\varphi\circ\Phi^{-1}(\mu(A)), A\in\mathcal{B}(\mathbf{R}^{n})
% \end{equation} For $A\in\mathcal{B}(\mathbb{R}^{n})$, define $
% A_{\varepsilon}:=\{x\in\mathbb{R}^{n},\exists y\in\mathbb{R}^{n}, |x-y|<\varepsilon\}
% $
% To consider rearrangements, it will be most helpful to employ the "integrated" form of (6). Assume by simple approximation that $A$ is the finite union of open balls. \\\\{\color{blue}Devraj: This is where I am a bit confused, for this argument to work we need the following $$
% \frac{d}{d\varepsilon}\mu\left(A_{\varepsilon}\right)=\lim_{h\rightarrow 0}\frac{\mu\left(A_{\varepsilon+h}\right)-\mu\left(A_{\varepsilon}\right)}{h}=\mu^+(A_{\epsilon})
% $$ Can we ensure this is true for a generic enlargement? I.e. for a generic cost function enlargement? I don't really see how this is true if the order isn't of h (meaning, it is a linear function of the euclidean norm). Perhaps someone can enlighten me?}
% {\color{red}Cyril: above the left hand side is a set. I guess you wanted to talk about measure of sets? In any case I suggest we specify to $R^n$ here and with enlargement with usual euclidean balls.}


% Define $v(\varepsilon)=\Phi^{-1}(\mu(A_{\varepsilon}))$. Observe that $$
% v'(\varepsilon)=\frac{\mu^+(A_{\varepsilon})}{\varphi\circ\Phi^{-1}(\mu(A_{\epsilon}))}\geq\sqrt{c}
% $$ Consequently, $$  
% v(\varepsilon)=v(0)+\int_{0}^{\epsilon}v'(t)dt\geq v(0)+\varepsilon\sqrt{c}
% $$ $$\iff$$ $$
% \Phi^{-1}(\mu(A_{\varepsilon}))\geq\Phi^{-1}(\mu(A))+\varepsilon\sqrt{c}
% $$ Thus, associating $A\subseteq\mathbb{R}^{n}$ to $A^*\subseteq\mathbb{R}^{n}$ where $$
% A^*=\{x\in\mathbb{R}^{n}:<x,u><a\}, a=\Phi^{-1}(\mu(A)), |u|=1
% $$ yields the conclusion $$
% \mu(A_{\varepsilon})\geq\Phi(\Phi^{-1}(\mu(A))+\varepsilon\sqrt{c})=\Phi(a+\varepsilon\sqrt{c})
% $$ Note that $$\Phi(a+\varepsilon\sqrt{c})=\mu(S), S:=\{x\in\mathbb{R}^{n}:<x,u><a+\varepsilon\sqrt{c}\}$$

% We claim that $A^{*}_{\varepsilon\sqrt{c}}\subseteq S$. Indeed, if $y\in A^{*}_{\varepsilon\sqrt{c}}$, there exists $x\in A^*$ such that $|x-y|<\varepsilon\sqrt{c}$, thus $$
% <y,u>=<x,u>+<y-x,u><a+|y-x||u|<a+\varepsilon\sqrt{c}
% $$
% This gives us the following conclusion 
% $$
% \mu(A_{\varepsilon})\geq\mu(A^*_{\varepsilon\sqrt{c}})
% $$This gives the following for $$
% Q_{c}f(x)=\sup_{y\in\mathbb{R}^{n}}f(y)-|x-y|
% $$\begin{equation}
%     \begin{split}
%         \mu\left(\{Q_{c}f>\lambda\}\right)&=\mu\left(\bigcup_{q_{1}-q_{2}>\lambda}\{f>q_{1}\}_{q_{2}}\right)\\
%         &\geq\sup_{q_{1}-q_{2}>\lambda}\mu\left(\{f>q_{1}\}_{q_{2}}\right)\\
%         &\geq\sup_{q_{1}-q_{2}>\lambda}\mu\left(\{f>q_{1}\}^*_{q_{2}\sqrt{c}}\right)\\
%         &=\mu\left(\bigcup_{q_{1}-q_{2}>\lambda}\{f>q_{1}\}^*_{q_{2}\sqrt{c}}\right)\\
%         &=\mu\left(\bigcup_{q_{1}-q_{2}>\lambda}\{f^*>q_{1}\}_{q_{2}\sqrt{c}}\right)\\
%         &=\mu\left(\{Q_{d}f_{*}>\lambda\}\right)
%     \end{split}
% \end{equation}where $$
% Q_{d}f_{*}(x)=\sup_{y\in\mathbb{R}^{n}}f_{*}(y)-\frac{|x-y|}{\sqrt{c}}
% $$

% \subsection{Product of Log-Concave measures in $\mathbf{R}^{n}$}
% In this section we consider measures of the form $
% \frac{d\mu_{n}}{dx}=\prod_{i\geq 1}^{n}\frac{1}{Z_{\Phi}}e^{-\Phi(|x_{i}|)}=\prod_{i\geq 1}^{n}\rho(x_{i})
% $ where $\Phi:\mathbb{R}_{+}\rightarrow\mathbb{R}_+$ is an increasing convex function where $\Phi(0)=0$ and $\sqrt{\Phi}$ is concave. Furthermore, assume that $\Phi$ is $\mathcal{C}^{2}(\Phi^{-1}(1),\infty)$. \\\\It was shown that there exists $K>0$ such that $$
% \mathcal{I}_{\mu_{n}}(t)\geq K\mathcal{I}_{\mu_{1}}(t)=K \rho\circ H^{-1}(t)
% $$ Recall that the isoperimetric profile decreases in dimension, this yields $$
% \mu_{n}^+(A)\geq K\rho\circ H^{-1}(\mu_{n}(A)), \forall A\in\mathcal{B}(R^{n})
% $$ Thus, associating $A\subseteq\mathbb{R}^{n}$ to $A^*\subseteq\mathbb{R}^{n}$ where $$
% A^*=\{x\in\mathbb{R}^{n}:<x,u><a\}, a=H^{-1}(\mu_{n}(A)), |u|=1
% $$ yields the conclusion $$
% \mu_{n}(A_{\varepsilon})\geq H(H^{-1}(\mu_{n}(A))+\varepsilon K)=H(a+\varepsilon K)
% $$ Note that $$H(a+\varepsilon K)=\mu_{n}(S), S:=\{x\in\mathbb{R}^{n}:<x,u><a+\varepsilon K\}$$ Thus, reasoning as above, 
% This gives us the following conclusion 
% $$
% \mu_{n}(A_{\varepsilon})\geq\mu_{n}(A^*_{\varepsilon\sqrt{c}})
% $$
% This would furnish the following 
% $$
% \mu_{n}\left(\{Q_{c}f>\lambda\}\right)\geq\mu_{n}\left(\{Q_{d}f_{*}>\lambda\}\right)
% $$ where $$
% Q_{d}f_{*}(x)=\sup_{y\in\mathbb{R}^{n}}f_{*}(y)-\frac{|x-y|}{K}
% $$

% \subsection{Euclidean Sphere}
% Consider the Euclidean sphere $S^{n-1}\subseteq\mathbb{R}^{n}$. Let d denote the geodesic distance in the section. In an expected manner, for some set $A\in\mathcal{B}(S^{n-1})$, let $$
% A_{\epsilon}=\{x\in S^{n-1}:d(x,A)<\epsilon\}
% $$ One may speak about balls on the sphere that we center at the north pole. Let us associate $A\subseteq S^{n-1}$ to $A^*\subseteq S^{n-1}$ where $A^*$ is a ball such that $\sigma_{n-1}(A)=\sigma_{n-1}(A^*)$ where $\sigma_{n-1}$ is the normalized uniform measure on $S^{n-1}$ One then has the following $$
% \sigma_{n-1}(A_\epsilon)\geq\sigma_{n-1}(A^*_{\epsilon})
% $$ This is the classical Isoperimetry theorem due to Lévy. From above, we see that for $$
% Q_{c}f(x)=\sup_{y\in S^{n-1}}f(y)-d(x,y)
% $$the following inequality holds 
% $$
% \sigma_{n-1}\left(\{Q_{c}f>\lambda\}\right)\geq\sigma_{n-1}\left(\{Q_{d}f_{*}>\lambda\}\right)
% $$ Let us now briefly comment on an application to the Hamilton-Jacobi equation. In this section, we will restrict our study to the sphere. Let us consider $\alpha:\mathbf{R}^+\rightarrow\mathbf{R}^+$ an increasing convex function of class $\mathcal{C}^1$ such that $\alpha(0)=0$. If $f:S^{n-1}\rightarrow\mathbf{R}$ is bounded, define 
% $$
% P_{t}f(x)=\sup_{y\in S^{n-1}}\left\{f(y)-t\alpha\left(\frac{d(x,y)}{t}\right)\right\}
% $$ For $\alpha^*(u)=\sup_{h\geq 0}\left\{hu-\alpha(h)\right\}$, the following equation holds $$
% \frac{d}{dt}P_{t}f(x)=\alpha^*\left(|\nabla^+P_{t}f|(x)\right)
% $$
% {\color{blue}Devraj:Can we even apply our rearrangement to this? Above we would take $\alpha(x)=x$ and then everything is trivial as $P_{t}f(x)$ no longer depends on t...am I neglecting something? Or is this structure not useful for a cost equal to the distance?}

% {\color{red}Cyril: if one considers 
% $\alpha(x)=|x|$ then $\alpha^*(x)=0$ on $[-1,1]$ and $\infty$ elsewhere. This is kind of pathological. One should check if the HJ equation has still a meaning (I'm not sure).
% So probably better not to speak about HJ here I guess, as a first approximation.
% }








\subsection{Functional Blaschke-Santalo's inequality via Semi-group revisited}

In this section we will prove the following well-known functional Blaschke-Santalo's inequality \cite{ball86}: for all measurable super linear function $f \in \mathrm{Cvx}_0(\mathbb{R}^n)$ 
%on $\mathbb{R}^n$ with 
%$\int e^{-f(x})dx < \infty$ 
it holds
\begin{equation} \label{eq:functionalBS}
    \int e^{-f(x)}dx \int e^{- \mathcal{L}f(x)}dx \leq (2\pi)^n .
\end{equation}
Recall that $\mathrm{Cvx}_0(\mathbb{R}^{n})$ denotes the set of non-negative even convex functions on $\mathbb{R}^{n}$ taking the value zero at at the origin and a function is called super-linear if for all $x \in \mathbb{R}^n$, $f(x) \geq a|x|+b$ for some $a>0$, $b \in \mathbb{R}$. Super-linearity  will allow us to avoid some non-crucial  technicalities. We follow \cite{cordero} here, with this choice.

We stress that there is nothing  new in this section: our aim is to show how our rearrangement results developed earlier can simplify some aspect of known proofs. More specifically, following the semi-group argument of 
\cite{cordero}, thanks to our reduction to dimension 1, their final Brascamp-Lieb argument will reduce to the trivial fact that a square is non-negative. We may, at different step, refer to \cite{artstein-klartag-milman,cordero} for some technical details.

% To achieves this we will proceed in two steps. In the first step we use some of our rearrangement's result in order to reduce the above statement to a one dimensional one. In the second step we will follow Cordero et al. and proceed through a semi-group argument. Thanks to the first step, it appears that the semi-group argument becomes technically much simpler than in Cordero et al \cite{cordero}. In particular, we won't need to  deal with convex superlinear functions as in \cite{cordero}.

% \subsection{Step 1}
% Step 1 is collected in the following lemma. Recall that $\omega_n$ denotes the volume of the $n$-dimensional Euclidean unit ball.

% \begin{lemma} \label{lem:step1}
%    For any $f$ on $\mathbb{R}^n$ with 
% $\int e^{-f(x})dx < \infty$ it holds
% \begin{equation*} 
%     \int e^{-f(x)}dx \int e^{- \mathcal{L}f(x)}dx 
%     \leq 
%     (n\omega_n)^2
%     \sup_{\Psi} 
% \int_0^\infty e^{-\Psi(r)}r^{n-1}dr \int_0^\infty e^{- \mathcal{L}\Psi(r)}r^{n-1}dr
% \end{equation*}  
% where the supremum runs over all non-decreasing function $\Psi$ defined on the positive axis, with $\Psi(0) = 0$ and $\mathcal{L}\Psi(r)\coloneqq \sup_{s >0} rs-\Psi(s)$. In particular 
% $$
% \sup_{f : \int e^{-f}<\infty} 
% \int e^{-f(x)}dx \int e^{- \mathcal{L}f(x)}dx 
%     = 
%     (n\omega_n)^2
%     \sup_{\Psi} 
% \int_0^\infty e^{-\Psi(r)}r^{n-1}dr \int_0^\infty e^{- \mathcal{L}\Psi(r)}r^{n-1}dr .
% $$
% \end{lemma}


% \begin{proof}
% Set 
% $$
% S \coloneqq (n\omega_n)^2
%     \sup_{\Psi} 
% \int_0^\infty e^{-\Psi(r)}r^{n-1}dr \int_0^\infty e^{- \mathcal{L}\Psi(r)}r^{n-1}dr
% $$
% where the supremum runs over all non-decreasing function $\Psi$ defined on the positive axis, with $\Psi(0) \geq 0$ (and observe that, since changing $\Psi$ into $\Psi+t$, $t \in \mathbb{R}$ does not affect the product of the integrals, the supremum defining $S$ is actually the same as the supremum appearing in the statement of the lemma (with the condition $\Psi(0)=0$).

% Assume first that $f$ is bounded below. Again, since the left hand side of the functional Blashke-Santalo's inequality is invariant by changing $f$ into $f+t$ for any $t \in \mathbb{R}$, we can assume that $f$ is non-negative. We then proceed as in the proof of Theorem \ref{thm: existence of extremizer}. Using the layer cake decomposition to get
% \begin{align*}
% \int e^{-f(x)}dx\int e^{-\mathcal{L}f(x)}dx
% & =
% \int_0^\infty |\{e^{-f} > t\}|dt \int_0^\infty |\{e^{-\mathcal{L}f} > t\}|dt \\
% & =
% \int_0^\infty |\{f < -\log t\}|dt \int_0^\infty |\{\mathcal{L}f < -\log t\}|dt .
% \end{align*}
% By \eqref{eq: increasing rearrangement characterization}, it holds
% $
% |\{f < -\log t\}|=|\{f < -\log t\}^*|=|\{f_* < -\log t\}|$ 
% for any $t>0$. Therefore, using Theorem \ref{thm: majorization rearrangement Legendre}, and the layer cake decomposition again  
% \begin{align*}
% \int e^{-f(x)}dx\int e^{-\mathcal{L}f(x)}dx
% & \leq 
% \int_0^\infty |\{f_* < -\log t\}|dt \int_0^\infty |\{\mathcal{L}f_* < -\log t\}|dt \\
% & =
% \int e^{-f_*(x)}dx\int e^{-\mathcal{L}f_*(x)}dx .
% \end{align*}
% Now by construction of the rearrangement, for any $f \geq 0$, there exists a function $\Psi \colon [0,\infty) \to [0,\infty)$, non-decreasing, such that
% $f_*(x)=\Psi(\|x\|)$, $x \in \mathbb{R}^n$. Passing to polar coordinates gives $\int e^{-f(x)}dx\int e^{-\mathcal{L}f(x)}dx \leq S$ when $f$ is bounded below. 

% Now consider a measurable function $f$ on $\mathbb{R}^n$ with $\int e^{-f}<\infty$ and set $f_N=\max(f,-N)$, for $N$ integer. Observe that $\mathcal{L}f_N \geq \mathcal{L}f$ so that, since $f_N$ is bounded below
% $$
% \int e^{-f_N} \int e^{-\mathcal{L}f} \leq 
% \int e^{-f_N} \int e^{-\mathcal{L}f_N} 
% \leq 
% S.
% $$
% By the dominated convergence theorem we conclude that the inequality in the lemma holds for any $f$ with $\int e^{-f}<\infty$ as expected. 

% The second part of the lemma is trivial, ending the proof of the lemma.
% \end{proof}

By Theorem \ref{thm: existence of extremizer}
we know that
$$
    \sup_{f\in \mathrm{Cvx}_0(\mathbb{R}^{n})} 
    \int e^{-f}\int e^{-\mathcal{L}f}
    =
    (n\omega_{n})^2 
    \sup_{\Psi} 
    \int_0^\infty e^{-\Psi(r)}r^{n-1}dr 
    \int_0^\infty e^{-\mathcal{L}\Psi(r)}r^{n-1}dr
$$
where the second supremum runs now over all super-linear\footnote{It is straight forward to adapt the proof of  Theorem \ref{thm: existence of extremizer} to take into account the extra super-linearity assumption considered here.}  convex non-decreasing function $\Psi \colon [0,\infty) \to [0,\infty)$ with $\Psi(0)=0$, and 
$\mathcal{L}\Psi(r)=\sup_{s>0}\{rs-\Psi(s)\}$. 
Super-linearity ensures in particular that the Laplace transform of $\Psi$ is finite on the whole $[0,\infty)$.




\subsubsection*{Technical preparation: the Bessel semi-group}

On $(0,\infty)$, define the following generator
$$
L^bf(r)=f''(r) +\frac{n-1}{r}f'(r), \qquad r>0
$$
where $f \colon (0,\infty) \to \mathbb{R}$ is smooth enough. We denote $(P_t^b)_{t \geq 0}$
its associated semi-group and observe that the reversible measure is $\nu(dr)=r^{n-1}dr$, $r \in (0,\infty)$. 
This is the projection of the heat semi-group $(P_t^h)_{t \geq 0}$ on radial functions, called the Bessel semi-group. More precisely, for $f \colon (0,\infty) \to \mathbb{R}$, define $\tilde{f}(x)=f(\|x\|)$, $x \in \mathbb{R}^n$ so that
$$
P_t^bf(\|x\|) = P_t^h \tilde f(x) = \tilde f* g_t(x), \qquad x \in \mathbb{R}^n ,
$$ 
with $g_t(x)\coloneqq e^{-\|x\|^2/(4t)}/(4\pi t)^{n/2}$.

%  The operator $L$ is the projection on radial functions of the Ornstein-Uhlenbeck generator $L^{ou}=\Delta -x \cdot \nabla$, whose semi-group we denote $R_t$.
% For $\tilde{f}(x)=f(\|x\|)$, $x \in \mathbb{R}^n$ it holds $R_t \tilde{f}(x)=P_tf(\|x\|)$.
% As is classical, from Melher's formula, for any function $h$  on $\mathbb{R}^n$, $R_th(x)=(h*g_{1-\rho})(\sqrt{\rho}x)$, $x \in \mathbb{R}^n$, with $g_s(x) \coloneqq e^{-\|x\|^2/(2s)}/(2\pi s)^{n/2}$  the $n$-dimensional density of a centered Gaussian of covariance $s\mathrm{Id}_{\mathbb{R}^n}$ and $\rho \coloneqq e^{-2t}$. It follows that  $P_tf(\|x\|)=(\tilde{f}*g_{1-\rho})(\sqrt{\rho}x)$ for any $x \in \mathbb{R}^n$.

% Next, we introduce the "Polar Focker-Franck" operator, \textit{i.e.}\ the dual of $L$ in the space 
% $\mathbb{L}^2(r^{n-1}dr,(0,\infty)])$, namely the operator $M$ satisfying for all $u,v  \colon (0,\infty) \to \mathbb{R}$ smooth enough,
% $$
% \int_0^\infty v(r)Lu(r) r^{n-1}dr 
% = 
% \int_0^\infty u(r)Mv(r) r^{n-1}dr .
% $$
% Integrating by parts one easily checks that 
% $$
% Mv(r)=v''(r) + (\frac{n-1}{r} + r)v'(r) + nv(r), \qquad r >0 .
% $$
% The associated semi-group we denote $(Q_t)_{t \geq 0}$.
% Given a function $v$ on the positive axis, define  $\tilde{v}(x)=v(\|x\|)$, $x \in \mathbb{R}^n$. 
% Let $R_t^* f=\frac{1}{\rho^{n/2}}f_\rho*g_{1-\rho}$ be the dual of the Ornstein-Uhlenbeck semi-group in $\mathbb{L}^2(dx,\mathbb{R}^n)$, with $f_\rho(x) \coloneqq f(x/\sqrt{\rho})$, $x \in \mathbb{R}^n$. 
% Then, the duality between $R_t$ and $R_t^*$, transferred to polar coordinates ensures that 
% $$
% Q_tv(r) = 
% %\frac{1}{n\omega_n} 
% \int_{\mathbb{S}^{n-1}} 
% R_t^*\tilde{v}(r\theta)d\sigma(\theta) ,
% %=\frac{1}{n\omega_n \rho^{n/2}} \int_{\mathbb{S}^{n-1}} (\tilde{v}*g_{1-\rho})(r\theta/\sqrt{\rho })d\theta 
% \qquad r >0
% $$
% where $d\sigma$ is the uniform measure on $\mathbb{S}^{n-1}$.

%In the limit $t \to \infty$ it holds %$R_t^*\tilde{v}(x) \to \frac{e^{-\|x\|^2/2}}{(2\pi)^{n/2}}\int_{\mathbb{R}^n} \tilde{v}(y)dy$ so that 
%$$
%\lim_{t \to \infty} Q_t v(r)=\frac{e^{-r^2/2}}{(2\pi)^{n/2}}\int_0^\infty v(r)r^{n-1}dr, \qquad r > 0 .
%$$

Consider a super-linear, convex, non-decreasing function $\Psi \colon (0,\infty) \to \mathbb{R}$, with $\Psi(0)=0$, and set 
$$
\Psi_t(r)
\coloneqq
-\log P_t^b(e^{-\Psi})(r), \qquad r \geq 0 .
$$
for $\Psi$ defined on the positive axis, super-linear non-decreasing, with $\Psi(0) = 0$. 
Being a convolution with a Gaussian kernel $P_t^be^{-\Psi}$ is positive, $\mathcal{C}^\infty$ smooth. 
Moreover, super-linearity ensures that $\Psi_t$ is also super-linear (see \cite{cordero} for the details), and together with the convexity property, $\Psi_t$ is strictly convex \cite{cordero}. 
In particular, $\Psi_t'$ has an reciprocal function we denote ${\Psi_t'}^{-1}$ and, as is easy to check from the definition,  
$(\mathcal{L}\Psi_t)'={\Psi_t'}^{-1}$
and $\mathcal{L}\Psi_t(r) = r\Psi_t(r)- \Psi_t((\mathcal{L}\Psi_t)'(r)))$
on the positive axis. Taking the derivative in $t$, it follows that 
$$
\frac{d}{dt}\mathcal{L}\Psi_t(r) 
=
r \frac{d}{dt}(\mathcal{L}\Psi_t)'(r)
-\Psi_t'((\mathcal{L}\Psi_t)'(r)) \frac{d}{dt}(\mathcal{L}\Psi_t)'(r)
- \frac{d}{dt}\Psi_t ((\mathcal{L}\Psi_t)'(r)) 
=
- \frac{d}{dt}\Psi_t ((\mathcal{L}\Psi_t)'(r)) 
$$
where the last equality follows from the identity $\Psi_t'((\mathcal{L}\Psi_t)'(r)) =r$.
Since $\Psi_t'=- \frac{P_t^b(e^{-\Psi})'}{P_t^b(e^{-\Psi})}$ and
 $\Psi_t''=- \frac{P_t^b(e^{-\Psi})''}{P_t^b(e^{-\Psi})}+ \left(\frac{P_t^b(e^{-\Psi})'}{P_t^b(e^{-\Psi})} \right)^2$, it holds
 \begin{align*}
     \frac{d}{dt}\Psi_t (r)
     & = 
     -\frac{L^b P_t^b(e^{-\Psi})(r)}{P_t^b(e^{-\Psi})(r)} \\
     & =
     \frac{-1}{P_t^b(e^{-\Psi})(r)} 
     \left( P_t^b(e^{-\Psi})''(r) + \frac{n-1}{r} P_t^b(e^{-\Psi})'(r)\right) \\
     & =
     -\Psi_t'(r)^2 + \Psi_t''(r) + \frac{n-1}{r} \Psi_t'(r) .
 \end{align*}
 Therefore
 \begin{align*}
 \frac{d}{dt} \mathcal{L}\Psi_t(r) 
 & = - \frac{d}{dt}\Psi_t((\mathcal{L}\Psi_t)'(r)) 
  = 
 - \frac{d}{dt}\Psi_t({\Psi_t'}^{-1}(r)) \\
 & =
 r^2 - \Psi_t''({\Psi_t'}^{-1}(r)) - \frac{n-1}{{\Psi_t'}^{-1}(r)}r , \qquad r >0 .
 \end{align*}
 Now $\Psi_t'( ({\mathcal{L}\Psi_t)'}(r))=r$ guarantees that
 $\Psi_t''((\mathcal{L}\Psi_t)'(r))(\mathcal{L}\Psi_t)''(r)=1$, so that, since $(\mathcal{L}\Psi_t)'={\Psi_t'}^{-1}$,
 $\Psi_t''({\Psi_t'}^{-1}(r))=\frac{1}{(\mathcal{L}\Psi_t)''(r)}$.
 It follows that
 \begin{equation} \label{eq:cordero}
 \partial_t \mathcal{L}\Psi_t(r) 
 =
 -\frac{1}{(\mathcal{L}\Psi_t)''(r)} + r^2 - \frac{n-1}{(\mathcal{L}\Psi_t)'(r)}r , \qquad r >0 .
 \end{equation}





% %By Prekopa-Leindler's inequality, Q_t
% .Therefore $\Psi_t$ is $\mathcal{C}^\infty$ smooth as well. Super-linearity guarantees that 
% $\mathcal{L}\Psi_t$ has support the whole $[0,\infty)$ and is strictly convex.

%   Then, as is standard, it holds for $r >0$
%   $$
%   \frac{\partial}{\partial_t}\Psi_t(r)=\frac{MQ_t(e^{-\Psi})}{Q_t(e^{-\Psi})},
%   $$
%   $$
%   (\mathcal{L}\Psi_t)'(r) = {\Psi_t'}^{-1}(r),
%   $$ 
%   and
%   $$
%   \mathcal{L}\Psi_t(r) = r(\mathcal{L}\Psi_t)'(r) - \Psi_t((\mathcal{L}\Psi_t)'(r))
%   $$
%   where ${\Psi_t'}^{-1}$ denotes the inverse function of $\Psi_t'$ that does exists since $\Psi_t$ is strictly convex. 
  % Let $(f_t)_{t \geq 0}$ be a family of $\mathcal{C}^1$ strictly convex functions on the line taking finite values. Assume that, for any $x \in \mathbb{R}$, $t \mapsto f_t(x)$ is differentiable. Then for all $t \geq 0$, $f_t'$ is invertible with inverse $(\mathcal{L}f_t)'$. Moreover,
  % $\mathcal{L}f_t(x)=x(\mathcal{L}f_t)'(x)-f_t((\mathcal{L}f_t)'(x))$ and $\partial_t \mathcal{L}f_t(x) = - \partial_tf_t((\mathcal{L}f)'(x))$. Here $\partial_t$ is a shorthand notation for the partial derivative with respect to the variable $t$ while
  % prime denotes the derivative with respect to $x$. 
  % In particular, for $f_t=-\log P_tg$ for some semi-group $(P_t)_{t \geq0}$ with infinitesimal generator $L$, and some function $g$, $\partial_t \mathcal{L}f_t(x) = \frac{LP_tg(y)}{P_tg(y)}|_{y=(\mathcal{L}f_t)'(x)}$, $x \in \mathbb{R}$. 
 %\end{lemma}

% \begin{proof}
%   That  $f_t'$ is invertible with inverse $(\mathcal{L}f_t)'$ and
%   $\mathcal{L}f_t(x)=x(\mathcal{L}f_t)'(x)-f((\mathcal{L}f)'(x))$
%   is clear. Taking the derivative with respect to $t$, we obtain
% $$
% \partial_t \mathcal{L}f_t (x) = 
% x \partial_t (\mathcal{L}f_t)'(x)
% -\partial_t f_t((\mathcal{L}f_t)'(x))
% - f_t'((\mathcal{L}f_t)'(x))\partial_t (\mathcal{L}f_t)'(x) .
% $$
% The identity ${f_t'}((\mathcal{L}f_t)'(x))=x$ leads to
% $\partial_t \mathcal{L}f_t (x) = 
% -\partial_t f_t((\mathcal{L}f_t)'(x))$
%   as expected.
% \end{proof}




\subsubsection*{Proof of  the functional Blaschke-Santalo's inequality \eqref{eq:functionalBS}}

Following \cite{cordero}, our aim is to prove that
$
t \mapsto 
\int_0^\infty e^{-\Psi_t(r)}r^{n-1}dr  
\int_0^\infty e^{-\mathcal{L}\Psi_t(r)}r^{n-1}dr  
$
 is non-decreasing. 
This would imply that 
$$
\int_0^\infty \!e^{-\Psi(r)}r^{n-1}dr  
\int_0^\infty \!e^{-\mathcal{L}\Psi(r)}r^{n-1}dr
\leq 
\lim_{t \to \infty}
\int_0^\infty \!e^{-\Psi_t(r)}r^{n-1}dr  
\int_0^\infty \!e^{-\mathcal{L}\Psi_t(r)}r^{n-1}dr
= \frac{(2\pi)^n}{(n \omega_n)^2}
$$
 where the limit can easily be computed thanks to the previous technical preparation (see \cite{artstein-klartag-milman,cordero} for  details). 
 
Observe first that, by invariance,
\begin{equation} \label{eq:0}
\int_0^\infty e^{-\Psi_t(r)}r^{n-1}dr = 
\int_0^\infty P_t^b(e^{-\Psi})r^{n-1}dr 
=
\int_0^\infty e^{-\Psi(r)} r^{n-1}dr .
\end{equation}
Therefore the first factor in $\int_0^\infty e^{-\Psi_t(r)}r^{n-1}dr  
\int_0^\infty e^{-\mathcal{L}\Psi_t(r)}r^{n-1}dr$  does not depend on $t$ and we are reduced to proving that
$$
t \mapsto \alpha(t) \coloneqq \int_0^\infty e^{-\mathcal{L}\Psi_t(r)}r^{n-1}dr 
$$
is non-decreasing on the positive axis.
Now, thanks to \eqref{eq:cordero}, it holds
\begin{align*}
\alpha'(t) 
& = 
-   
    \int_0^\infty  \partial_t \mathcal{L}\Psi_t(r) e^{-\mathcal{L}\Psi_t(r)}r^{n-1}dr  \\
&    =
    \int_0^\infty   \left( \frac{1}{(\mathcal{L}\Psi_t)''(r)} - r^2 + \frac{n-1}{(\mathcal{L}\Psi_t)'(r)}r \right) e^{-\mathcal{L}\Psi_t(r)}r^{n-1}dr .     
\end{align*}
% Since $\Psi_t'=- \frac{Q_t(e^{-\Psi})'}{Q_t(e^{-\Psi})}$ and
% $\Psi_t''=-\frac{Q_t(e^{-\Psi})''}{Q_t(e^{-\Psi})}+\left(\frac{Q_t(e^{-\Psi})'}{Q_t(e^{-\Psi})}\right)^2$, for $s:=(\mathcal{L}\Psi_t)'(r)$, from the technical preparation we infer that 
% \begin{align*}
%    \partial_t \mathcal{L}\Psi_t(r) 
%    & = 
%     \frac{M Q_t(e^{-\Psi})(s)}{Q_t(e^{-\Psi})(s)} \\
%     & =
%     \frac{Q_t(e^{-\Psi})''(s)}{Q_t(e^{-\Psi})(s)}
%     + \left( \frac{n-1}{s} + s \right) \frac{Q_t(e^{-\Psi})'(s)}{Q_t(e^{-\Psi})(s)}
%     + n \\
%     & =
%     \Psi_t'(s)^2 - \Psi_t''(r) -\left( \frac{n-1}{s}+s \right) \Psi_t'(s) +n  .
% \end{align*}
% In turn, since $(\mathcal{L}\Psi_t)'(r)={\Psi_t'}^{-1}(r)$, 
% \begin{align*}
% \partial_t \mathcal{L}\Psi_t(r) 
% & =
% r^2 - \Psi_t''({\Psi_t'}^{-1}(r)) -\left( \frac{n-1}{(\mathcal{L}\Psi_t)'(r)}+(\mathcal{L}\Psi_t)'(r) \right) r + n , \qquad r >0 .
% \end{align*}
% Now $\Psi_t'( {\mathcal{L}\Psi_t'}(r))=r$ guarantees that
% $\Psi_t''({\mathcal{L}\Psi_t'}(r)){\mathcal{L}\Psi_t''}(r)=1$, so that
% $\Phi_t''({\Phi_t'}^{-1}(r))=\frac{1}{(\mathcal{L}\Phi_t)''(r)}$.
% It follows that
% $$
% \partial_t \mathcal{L}\Phi_t(r) 
% =
% r^2 - \frac{1}{(\mathcal{L}\Phi_t)''(r)} -\left( \frac{n-1}{(\mathcal{L}\Phi_t)'(r)}+(\mathcal{L}\Phi_t)'(r) \right)r + n  , \qquad r >0 .
% $$
% Therefore $\alpha'(t) \geq 0$ is equivalent to proving  that
% $$
% \int_0^\infty \left(\frac{1}{(\mathcal{L}\Psi_t)''(r)} - r^2 + \left( \frac{n-1}{(\mathcal{L}\Psi_t)'(r)}+(\mathcal{L}\Psi_t)'(r) \right) - n  \right) e^{-(\mathcal{L}\Psi_t)(r)}r^{n-1}dr \geq 0  %$$
To ease notation, we set   $V=\mathcal{L}\Psi_t$ and observe that, $V',V''>0$. In order to handle the border case in the forthcoming integration by part, we introduce for $\varepsilon >0$,
$V_\varepsilon(r) \coloneqq V(r) + \varepsilon r$, $r \geq 0$. We claim that
\begin{equation} \label{eq:claim}
\int_0^\infty  \left( \frac{1}{V_\varepsilon''(r)} - r^2 + \frac{n-1}{V_\varepsilon'(r)}r \right) e^{-V_\varepsilon(r)}r^{n-1}dr     \geq 0 .
\end{equation}
This would imply $\alpha'(t) \geq 0$ by the monotone convergence theorem, in the limit $\varepsilon \downarrow 0$. Therefore to prove \eqref{eq:functionalBS} it only remains to establish \eqref{eq:claim}. To that aim, we use an integration by part to get 
\begin{align*}
-\int e^{-V_\varepsilon(r)}r^{n+1}dr
& =
\int -V_\varepsilon'(r)e^{-V_\varepsilon(r)} \frac{r^{n+1}}{V_\varepsilon'(r)}dr \\
& =
\left[ e^{-V_\varepsilon(r)}r^{n+1}/V_\varepsilon'(r)\right]_0^\infty 
- \int_0^\infty \left(
\frac{(n+1)r^n}{V_\varepsilon'(r)} - \frac{V''_\varepsilon(r)r^{n+1}}{V'_\varepsilon(r)^2} \right)
e^{-V_\varepsilon(r)} dr\\
& =
- \int_0^\infty 
\frac{(n-1)r}{V_\varepsilon'(r)} e^{-V_\varepsilon(r)} r^{n-1}dr
- \int_0^\infty \left(
\frac{2r}{V_\varepsilon'(r)} - \frac{V''_\varepsilon(r)r^2}{V_\varepsilon'(r)^2} \right)
e^{-V_\varepsilon(r)} r^{n-1}dr 
\end{align*}
where in the last equality we used that $V_\varepsilon$ is strictly convex and that $V_\varepsilon'(0) \geq \varepsilon$ so that the bracket term vanishes.
%{(\color{red}Cyril: I used that the first term (bracket term) vanishes. To be checked.)}
%Similarly
%\begin{align*}
%\int_0^\infty e^{-V(r)} V'(r) r^{n}dr
%& = 
%\left[ -e^{-V(r)} r^n \right]_0^\infty
%- \int_0^\infty n e^{-V(r)} r^{n-1}dr     
% =
%- \int_0^\infty n e^{-V(r)} r^{n-1}dr .
%\end{align*}
Therefore,
\begin{align*}
 & \int_0^\infty \left(\frac{1}{V_\varepsilon''(r)} -r^2 + \frac{n-1}{V_\varepsilon'(r)}r \right)  e^{-V_\varepsilon(r)}r^{n-1}dr \\
  & \qquad \qquad \qquad\qquad=
  \int_0^\infty \left(\frac{1}{V_\varepsilon''(r)} -\frac{2r}{V_\varepsilon'(r)} + \frac{V_\varepsilon''(r)r^2}{V_\varepsilon'(r)^2} \right) e^{-V_\varepsilon(r)}r^{n-1}dr \\
  &  \qquad \qquad \qquad\qquad=
  \int_0^\infty \frac{(rV_\varepsilon''(r)-V_\varepsilon'(r))^2}{V''_\varepsilon(r)V'_\varepsilon(r)^2} e^{-V_\varepsilon(r)}r^{n-1}dr \\
  & \qquad \qquad \qquad\qquad 
  \geq 0
\end{align*}
as expected. 

% Hence we proved that $\alpha$ is non-decreasing which, as mentioned at the begining of step 2, proove the funcitonal Blashke-Santalo's inequality of \cite{artstein-klartag-milman}.












% \subsubsection{Reproducing Cordero et al after rearrangement}

% As an illustration of the rearrangement techniques, let us reproduce the proof of Cordero et al of the functional Blaschke Santalo.  

% We start with an obvious lemma. {(\color{red}Cyril: which however should be stated carrefully, there are some technical issues here.)}

% \begin{lemma} \label{lem:Cyril}
%   Let $(f_t)_{t \geq 0}$ be a family of $\mathcal{C}^1$ strictly convex functions on the line taking finite values. Assume that, for any $x \in \mathbb{R}$, $t \mapsto f_t(x)$ is differentiable. Then for all $t \geq 0$, $f_t'$ is invertible with inverse $(\mathcal{L}f_t)'$. Moreover,
%   $\mathcal{L}f_t(x)=x(\mathcal{L}f_t)'(x)-f_t((\mathcal{L}f_t)'(x))$ and $\partial_t \mathcal{L}f_t(x) = - \partial_tf_t((\mathcal{L}f)'(x))$. Here $\partial_t$ is a shorthand notation for the partial derivative with respect to the variable $t$ while
%   prime denotes the derivative with respect to $x$. 
%   In particular, for $f_t=-\log P_tg$ for some semi-group $(P_t)_{t \geq0}$ with infinitesimal generator $L$, and some function $g$, $\partial_t \mathcal{L}f_t(x) = \frac{LP_tg(y)}{P_tg(y)}|_{y=(\mathcal{L}f_t)'(x)}$, $x \in \mathbb{R}$. 
%  \end{lemma}

% \begin{proof}
%   That  $f_t'$ is invertible with inverse $(\mathcal{L}f_t)'$ and
%   $\mathcal{L}f_t(x)=x(\mathcal{L}f_t)'(x)-f((\mathcal{L}f)'(x))$
%   is clear. Taking the derivative with respect to $t$, we obtain
% $$
% \partial_t \mathcal{L}f_t (x) = 
% x \partial_t (\mathcal{L}f_t)'(x)
% -\partial_t f_t((\mathcal{L}f_t)'(x))
% - f_t'((\mathcal{L}f_t)'(x))\partial_t (\mathcal{L}f_t)'(x) .
% $$
% The identity ${f_t'}((\mathcal{L}f_t)'(x))=x$ leads to
% $\partial_t \mathcal{L}f_t (x) = 
% -\partial_t f_t((\mathcal{L}f_t)'(x))$
%   as expected.
% \end{proof}

%  Recall that we "proved" earlier that
%  $$
%      \sup_{f\in \tilde C(\mathbb{R}^n)} 
%      \int e^{-f}\int e^{-\mathcal{L}f}
%      =
%      \sup_{\Psi} 
%     \int e^{-\Psi(|x|)}dx 
%      \int_0^\infty e^{-\mathcal{L}\Psi(|x|}dx  
%  $$
%  with $\Psi$ convex, increasing with $\Psi(0)=0$.
%  Changing $\Psi$ into $\Psi+c$ for any constant $c$ doesn't affect the volume product, therefore the constraint $\Psi(0)=0$ might be removed if needed.

%  In Cordero et al they consider the heat semi-group $P_t^h$, with generator the Laplacian $L^h=\Delta$ and put $f=-\log P_t^h(e^{-g})$. In view of the rearrangement procedure, it is natural to consider the projection of the heat-semi-group on its radial part. This is known as the Bessel semi-group we denote $P_t^b$, and whose infinitesimal generator is acting on smooth functions $\Psi \colon (0,\infty) \to \mathbb{R}$ as

%  $$
%  L^b \Psi (r) = \Psi''(r) + \frac{n-1}{r}\Psi'(r), \qquad r \in (0,\infty) .
%  $$
%  In fact it is not difficult to see that, for $f(x)=\Psi(|x|)$,
%  $P_t^hf(x)=P_t^b \Psi(|x|)$. The invariant measure for the Bessel semi-group is the measure with density $r^{n-1}$ with respect to the Lebesgue measure on $(0,\infty)$.

%  Therefore, the natural counterpart of the choice $f_t=-\log P_t^h(e^{-g})$ is 
%  $$
%  \Phi_t(x)=-\log P_t^b(e^{-\Psi})(|x|), 
%  $$
%  for $\Psi$ defined on the positive axis, convex increasing, with $\Psi(0)=0$. 

%  Set, passing to polar coordinates and by invariance,
%  \begin{align*}
%  \alpha(t)
%  & \coloneqq
%  \int_{\mathbb{R}^n} e^{-\Phi_t(x)}dx 
%      \int_{\mathbb{R}^n} e^{-\mathcal{L}\Phi_t(x)}dx \\
%  & =
%  \int_{\mathbb{R}^n} P_t^b(e^{-\Psi})(|x|) dx 
%      \int_{\mathbb{R}^n} e^{-\mathcal{L}\Phi_t(|x|)}dx \\
%      & =
%  (n\omega_n)^2    \int_0^\infty P_t^b(e^{-\Psi})(r)r^{n-1}dr  
%      \int_0^\infty e^{-\mathcal{L}\Phi_t(r)}r^{n-1}dr \\
%      & =
%     (n\omega_n)^2    \int_0^\infty e^{-\Psi(r)}r^{n-1}dr  
%      \int_0^\infty e^{-\mathcal{L}\Phi_t(r)}r^{n-1}dr  
%  \end{align*}
%  Our aim is to prove that $\alpha$ is non-decreasing. To see this, take the derivative 
%  $$
%  \alpha'(t) = 
%  - (n\omega_n)^2 \int_0^\infty e^{-\Psi(r)}r^{n-1}dr  
%      \int_0^\infty \partial_t \mathcal{L}\Phi_t(r) e^{-\mathcal{L}\Phi_t(r)}r^{n-1}dr .
%  $$
%  {\color{red}Cyril: In the language of Cordero et al the latter is written for a generic $f_t=-\log P_t^h(e^{-g})$ in $\mathbb{R}^n$ (not necessarily spherically symmetric) and reads
%  $$
%  \alpha'(t)
%  =
%  - \int_{\mathbb{R}^n} e^{-g(x)}dx 
%     \int_{\mathbb{R}^n} \partial_t \mathcal{L}f_t(x) e^{-\mathcal{L}f_t(x)}dx
%  $$
%  In fact both expressions are equal and one can pass from one to the other by changing variables.}

%  By Lemma \ref{lem:Cyril} it holds
%  $\partial_t \mathcal{L}\Phi_t(r) = - \partial_t\Phi_t((\mathcal{L}\Phi_t)'(r))$. Now for the Bessel semi-group, (since $\Phi_t'=-P_t^b(e^{-\Psi})'/P_t^b(e^{-\Psi})$ and
%  $\Phi_t''=-P_t^b(e^{-\Psi})''/P_t^b(e^{-\Psi})+(P_t^b(e^{-\Psi})'/P_t^b(e^{-\Psi}))^2$)
%  \begin{align*}
%      \partial_t\Phi_t (r)
%      & = 
%      -\frac{L^b P_t^b(e^{-\Psi})(r)}{P_t^b(e^{-\Psi})(r)} \\
%      & =
%      \frac{-1}{P_t^b(e^{-\Psi})(r)} 
%      \left( P_t^b(e^{-\Psi})''(r) + \frac{n-1}{r} P_t^b(e^{-\Psi})'(r)\right) \\
%      & =
%      -\Phi_t'(r)^2 + \Phi_t''(r) + \frac{n-1}{r} \Phi_t'(r) .
%  \end{align*}
%  In turn
%  \begin{align*}
%  \partial_t \mathcal{L}\Phi_t(r) 
%  & = - \partial_t\Phi_t((\mathcal{L}\Phi_t)'(r)) \\
%  & = 
%  - \partial_t\Phi_t({\Phi_t'}^{-1}(r)) \\
%  & =
%  r^2 - \Phi_t''({\Phi_t'}^{-1}(r)) - \frac{n-1}{{\Phi_t'}^{-1}(r)}r , \qquad r >0 .
%  \end{align*}
%  Now $\Phi_t'( {\mathcal{L}\Phi_t'}(r))=r$ guarantees that
%  $\Phi_t''({\mathcal{L}\Phi_t'}(r)){\mathcal{L}\Phi_t''}(r)=1$, so that, since ${\mathcal{L}\Phi_t'}={\Phi_t'}^{-1}$,
%  $\Phi_t''({\Phi_t'}^{-1}(r))=\frac{1}{(\mathcal{L}\Phi_t)''(r)}$.
%  It follows that
%  $$
%  \partial_t \mathcal{L}\Phi_t(r) 
%  =
%  -\frac{1}{(\mathcal{L}\Phi_t)''(r)} + r^2 - \frac{n-1}{{\Phi_t'}^{-1}(r)}r , \qquad r >0 .
%  $$
%  {\color{red}Cyril: The same computation is done in Cordero et al, leading, for the heat flow, to
%  $$
%  \partial_t \mathcal{L}\Phi_t(x) 
%  =
%  - Tr(D^2_x \mathcal{L}\Phi_t)^{-1} + |x|^2 ,
%  \qquad x \in \mathbb{R}^n .
%  $$
%  The extra term in front of $n-1$ is precisely coming from the change in polar coordinates.

%  This suggests the following fact that looks non trivial to me:
%  for $f(x)=g(|x|)$ (recall that ${g'}^{-1}=(\mathcal{L}g)'$)
%  $$
%  Tr(D^2_x \mathcal{L}f)^{-1} = 
%  \frac{1}{(\mathcal{L}g)''(|x|)} + \frac{n-1}{(\mathcal{L}g)'(|x|)}|x|, \qquad x \in \mathbb{R}^n .
%  $$
%  }

%  It follows that $\alpha'(t) \geq 0$ if we can prove that
%  $$
%  \int_0^\infty \left(\frac{1}{(\mathcal{L}\Phi_t)''(r)} - r^2 + \frac{n-1}{{\Phi_t'}^{-1}(r)}r \right) e^{-(\mathcal{L}\Phi_t)(r)}r^{n-1}dr \geq 0  
%  $$
%  {\color{red}Cyril: while Cordero et al needed to prove that
%  $$
%  \int_{\mathbb{R}^n} \left( Tr(D^2_x \mathcal{L}\Phi_t)^{-1} - |x|^2\right) e^{-\mathcal{L}\Phi_t(x)} dx 
%  \geq 0
%  $$
%  which is proven by Brascamp-Lieb applied to the function $x \mapsto x_i$ and summing up over $i$.}

%  In order to reproduce the argument by Cordero et al, we would need to prove that, since ${\Phi'_t}^{-1}=(\mathcal{L}\Phi)'$,
%  $$
%  \int_0^\infty \left(\frac{1}{(\mathcal{L}\Phi_t)''(r)} -r^2 + \frac{n-1}{(\mathcal{L}\Phi_t)'(r)}r \right) e^{-(\mathcal{L}\Phi_t)(r)}r^{n-1}dr \geq 0 .
%  $$
%  {\color{red}Cyril: whats is about to come should be made rigorous. In particular, we need some properties on $V:=\mathcal{L}\Phi_t$.}
%  Set $V=\mathcal{L}\Phi_t$. Then the latter is recast
%  $$
%  \int_0^\infty \left(\frac{1}{V''(r)} -r^2 + \frac{n-1}{V'(r)}r \right) e^{-V(r)}r^{n-1}dr \geq 0 .
%  $$
%  Integrating by part, we have {(\color{red}Cyril: we need $V'>0$ here)}
%  \begin{align*}
%  -\int e^{-V(r)}r^{n+1}dr
%  & =
%  \int -V'(r)e^{-V(r)} \frac{r^{n+1}}{V'(r)}dr \\
%  & =
%  \left[ e^{-V(r)}r^{n+1}/V'(r)\right]_0^\infty 
%  - \int_0^\infty \left(
%  \frac{(n+1)r^n}{V'(r)} - \frac{V''(r)r^{n+1}}{V'(r)^2} \right)
%  e^{-V(r)} dr\\
%  & =
%  - \int_0^\infty 
%  \frac{(n-1)r}{V'(r)} e^{-V(r)} r^{n-1}dr
%  - \int_0^\infty \left(
%  \frac{2r}{V'(r)} - \frac{V''(r)r^2}{V'(r)^2} \right)
%  e^{-V(r)} r^{n-1}dr .
%  \end{align*}
%  {(\color{red}Cyril: I used that the first term (bracket term) vanishes. To be checked.)}
%  Therefore,
%  \begin{align*}
%    \int_0^\infty \left(\frac{1}{V''(r)} -r^2 + \frac{n-1}{V'(r)}r \right) e^{-V(r)}r^{n-1}dr
%    & =
%    \int_0^\infty \left(\frac{1}{V''(r)} -\frac{2r}{V'(r)} + \frac{V''(r)r^2}{V'(r)^2} \right) e^{-V(r)}r^{n-1}dr \\
%    & =
%    \int_0^\infty \frac{(rV''(r)-V'(r))^2}{V''(r)V'(r)^2} e^{-V(r)}r^{n-1}dr \\
%    & \geq 0
%  \end{align*}
%  as expected {(\color{red}Cyril: we can assume $V$ strictly convex for example here. i.e. we need to check that $\Psi$ convex increasing implies $V=\mathcal{L}\Phi_t$ is strictly convex increasing.)}
%  This reproduces Cordero et al argument, after rearrangement. {\color{red}Cyril: There is no real gain with respect to Cordero et al, just that we are dealing with convex functions on the positive axis only, which probably makes life easier for regularity/technicalities.}


%  \bigskip 

%  \begin{center}
%  \underline{\hspace{5cm}}
% \end{center}

%  {\color{red}Cyril: I don't really like the fact that, in the proof of Cordero et al, we need, in the end, to rescale, which, at the level of the semi-groups means to go from the heat flow to the Focker Planck flow. 

%  In this section I would like to use and introduce only one semi-group in order to avoid rescalings etc. The natural candidate is the Focker Planck semi-group, after passing to polar coordinate. I now reproduce the proof above without details (there are some technical issues I guess one would need to fix if one wants to make the argument rigorous. The idea is to see what could be implemented for the polar version of Baschke Santalo functional inequality: we could try to mimick the proof below with the Laguerre semi-group).
%  }

% On $(0,\infty)$, define the following generator (the projection of the Ornstein-Uhlenbeck semi-group on radial functions),
% $$
% Lf(r)=f''(r) +(\frac{n-1}{r} - r)f'(r), \qquad r>0
% $$
% where $f \colon (0,\infty) \to \mathbb{R}$ is smooth enough. We denote $(P_t)_{t \geq 0}$
% its associated semi-group and observe that the reversible measure is $\nu(dr)=n\omega_n r^{n-1}e^{-r^2/2}/(2 \pi)^{n/2}dr$, $r \in (0,\infty)$.

% Next, we introduce the "Polar Focker-Franck" operator, that is the dual of $L$ in 
% $\mathbb{L}^2(r^{n-1}dr)$. Namely the operator $M$ satisfying for all $f,g  \colon (0,\infty) \to \mathbb{R}$ smooth enough,
% $$
% \int_0^\infty f(r)Lg(r) r^{n-1}dr 
% = 
% \int_0^\infty g(r)Mf(r) r^{n-1}dr .
% $$
% Integrating by parts one easily gets
% $$
% Mf(r)=f''(r) + (\frac{n-1}{r} + r)f'(r) + nf(r), \qquad r >0 .
% $$
% The associated semi-group we denote $(Q_t)_{t \geq 0}$.
% Consider
% $$
% \Phi_t(x)=-\log Q_t(e^{-\Psi})(|x|), 
% $$
% for $\Psi$ defined on the positive axis, convex increasing, with $\Psi(0)=0$. 
% Set, passing to polar coordinates,
% \begin{align*}
% \alpha(t)
% & \coloneqq
% \int_{\mathbb{R}^n} e^{-\Phi_t(x)}dx 
%     \int_{\mathbb{R}^n} e^{-\mathcal{L}\Phi_t(x)}dx \\
% & =
% \int_{\mathbb{R}^n} Q_t(e^{-\Psi})(|x|) dx 
%     \int_{\mathbb{R}^n} e^{-\mathcal{L}\Phi_t(|x|)}dx \\
%     & =
% (n\omega_n)^2    \int_0^\infty Q_t(e^{-\Psi})(r)r^{n-1}dr  
%     \int_0^\infty e^{-\mathcal{L}\Phi_t(r)}r^{n-1}dr 
%     %\\
%     %& =
%    %(n\omega_n)^2    \int_0^\infty e^{-\Psi(r)}r^{n-1}dr  
%    % \int_0^\infty e^{-\mathcal{L}\Phi_t(r)}r^{n-1}dr  
% \end{align*}
% Our aim is to prove that $\alpha$ is non-decreasing. To see this, observe first that
% (by duality)
% \begin{equation} \label{eq:0}
% \frac{d}{dt} \int_0^\infty Q_t(e^{-\Psi})(r)r^{n-1}dr 
% =
% \int_0^\infty M Q_t(e^{-\Psi})(r)r^{n-1}dr
% =
% \int_0^\infty Q_t(e^{-\Psi})(r) L \mathds{1} (r)r^{n-1}dr
% = 0 .
% \end{equation}
% Therefore
% $$
% \alpha'(t) = 
% - (n\omega_n)^2 \int_0^\infty e^{-\Psi(r)}r^{n-1}dr  
%     \int_0^\infty \partial_t \mathcal{L}\Phi_t(r) e^{-\mathcal{L}\Phi_t(r)}r^{n-1}dr .
% $$
% By Lemma \ref{lem:Cyril} it holds
% (since $\Phi_t'=-Q_t(e^{-\Psi})'/Q_t(e^{-\Psi})$ and
% $\Phi_t''=-Q_t(e^{-\Psi})''/Q_t(e^{-\Psi})+(Q_t(e^{-\Psi})'/Q_t(e^{-\Psi}))^2$), for $s:=(\mathcal{L}\Phi_t)'(r)$,
% \begin{align*}
%    \partial_t \mathcal{L}\Phi_t(r) 
%    & = 
%     \frac{M Q_t(e^{-\Psi})(s)}{Q_t(e^{-\Psi})(s)} \\
%     & =
%     \frac{Q_t(e^{-\Psi})''(s)}{Q_t(e^{-\Psi})(s)}
%     + \left( \frac{n-1}{s} + s \right) \frac{Q_t(e^{-\Psi})'(s)}{Q_t(e^{-\Psi})(s)}
%     + n \\
%     & =
%     \Phi_t'(s)^2 - \Phi_t''(r) -\left( \frac{n-1}{s}+s \right) \Phi_t'(s) +n  .
% \end{align*}
% In turn, since $(\mathcal{L}\Phi_t)'(r)={\Phi_t'}^{-1}(r)$, 
% \begin{align*}
% \partial_t \mathcal{L}\Phi_t(r) 
% & =
% r^2 - \Phi_t''({\Phi_t'}^{-1}(r)) -\left( \frac{n-1}{(\mathcal{L}\Phi_t)'(r)}+(\mathcal{L}\Phi_t)'(r) \right) r + n , \qquad r >0 .
% \end{align*}
% Now $\Phi_t'( {\mathcal{L}\Phi_t'}(r))=r$ guarantees that
% $\Phi_t''({\mathcal{L}\Phi_t'}(r)){\mathcal{L}\Phi_t''}(r)=1$, so that
% $\Phi_t''({\Phi_t'}^{-1}(r))=\frac{1}{(\mathcal{L}\Phi_t)''(r)}$.
% It follows that
% $$
% \partial_t \mathcal{L}\Phi_t(r) 
% =
% r^2 - \frac{1}{(\mathcal{L}\Phi_t)''(r)} -\left( \frac{n-1}{(\mathcal{L}\Phi_t)'(r)}+(\mathcal{L}\Phi_t)'(r) \right)r + n  , \qquad r >0 .
% $$
% It follows that $\alpha'(t) \geq 0$ if we can prove that
% $$
% \int_0^\infty \left(\frac{1}{(\mathcal{L}\Phi_t)''(r)} - r^2 + \left( \frac{n-1}{(\mathcal{L}\Phi_t)'(r)}+(\mathcal{L}\Phi_t)'(r) \right) - n  \right) e^{-(\mathcal{L}\Phi_t)(r)}r^{n-1}dr \geq 0  
% $$

% Set $V=\mathcal{L}\Phi_t$ for simplicity. Then the latter is recast
% $$
% \int_0^\infty \left(\frac{1}{V''(r)} -r^2 + \left( \frac{n-1}{V'(r)}+V'(r) \right)r  - n\right) e^{-V(r)}r^{n-1}dr \geq 0 .
% $$
% Integrating by part, we have {(\color{red}Cyril: we need $V'>0$ here)}
% \begin{align*}
% -\int e^{-V(r)}r^{n+1}dr
% & =
% \int -V'(r)e^{-V(r)} \frac{r^{n+1}}{V'(r)}dr \\
% & =
% \left[ e^{-V(r)}r^{n+1}/V'(r)\right]_0^\infty 
% - \int_0^\infty \left(
% \frac{(n+1)r^n}{V'(r)} - \frac{V''(r)r^{n+1}}{V'(r)^2} \right)
% e^{-V(r)} dr\\
% & =
% - \int_0^\infty 
% \frac{(n-1)r}{V'(r)} e^{-V(r)} r^{n-1}dr
% - \int_0^\infty \left(
% \frac{2r}{V'(r)} - \frac{V''(r)r^2}{V'(r)^2} \right)
% e^{-V(r)} r^{n-1}dr .
% \end{align*}
% {(\color{red}Cyril: I used that the first term (bracket term) vanishes. To be checked.)}
% Similarly
% \begin{align*}
% \int_0^\infty e^{-V(r)} V'(r) r^{n}dr
% & = 
% \left[ -e^{-V(r)} r^n \right]_0^\infty
% - \int_0^\infty n e^{-V(r)} r^{n-1}dr     \\
% & =
% - \int_0^\infty n e^{-V(r)} r^{n-1}dr .
% \end{align*}
% Therefore,
% \begin{align*}
%  & \int_0^\infty \left(\frac{1}{V''(r)} -r^2 + \left( \frac{n-1}{V'(r)}+V'(r) \right)r + n \right)  e^{-V(r)}r^{n-1}dr \\
%   & \qquad \qquad \qquad\qquad=
%   \int_0^\infty \left(\frac{1}{V''(r)} -\frac{2r}{V'(r)} + \frac{V''(r)r^2}{V'(r)^2} \right) e^{-V(r)}r^{n-1}dr \\
%   &  \qquad \qquad \qquad\qquad=
%   \int_0^\infty \frac{(rV''(r)-V'(r))^2}{V''(r)V'(r)^2} e^{-V(r)}r^{n-1}dr \\
%   & \qquad \qquad \qquad\qquad 
%   \geq 0
% \end{align*}
% as expected. As a consequence $t \mapsto \alpha(t)$
% is non-decreasing so that
% $$
% \alpha(0) = \int_{\mathbb{R}^n} e^{-\Psi} \int e^{\mathcal{L}\Psi}
% \leq 
% \lim_{t \to \infty} \alpha(t) 
% $$
% and we are left with computing 
% \begin{align*}
% \lim_{t \to \infty} \alpha(t) 
% & =
% \lim_{t \to \infty} (n\omega_n)^2    \int_0^\infty Q_t(e^{-\Psi})(r)r^{n-1}dr  
%     \int_0^\infty e^{-\mathcal{L}\Phi_t(r)}r^{n-1}dr \\
% &    =
% \int_{\mathbb{R}^n} e^{-\Psi}(x)dx  
% \lim_{t \to \infty}  \int_{\mathbb{R}^n} e^{-\mathcal{L}\Phi_t(x)}dx 
% \end{align*}
% where in the last line we used \eqref{eq:0}.

% By construction, when $t$ tends to infinity, $Q_tg(r) \to n\omega_n \frac{e^{-r^2/2}}{(2\pi)^{n/2}} \int_0^\infty g(r)r^{n-1}dr$ (to see this one can use the relation $\int_0^\infty f(r)Q_t g(r) r^{n-1}dr = \int_0^\infty g(r)P_tf(r)r^{n-1}$ with $f=\delta_r$,
% the Dirac mass at some fixed point $r \in (0,\infty)$, and the fact that $P_tf \to n\omega_n \int f(s)e^{-s^2/2}s^{n-1}dr /(2\pi)^{n/2} = n\omega_n e^{-r^2/2}r^{n-1}/(2\pi)^{n/2}$ in the limit $t \to \infty$). 

% Hence,
% $$
% \lim_{t \to \infty} \Phi_t (x)
% =
% - \log \left( n\omega_n \frac{e^{-|x|^2/2}}{(2\pi)^{n/2}} \int_0^\infty e^{-\Psi(r)}r^{n-1}dr \right)
% =
% - \log \left( \frac{1}{(2\pi)^{n/2}} \int_{\mathbb{R}^n} e^{-\Psi(x)}dx \right)
% + \frac{|x|^2}{2}
% $$
% so that {(\color{red}Cyril: again there are some technical issues here)}
% $$
% \lim_{t \to \infty}  \mathcal{L}\Phi_t(x)
% =
% \log \left( \frac{1}{(2\pi)^{n/2}} \int_{\mathbb{R}^n} e^{-\Psi(x)}dx  \right) 
% +
% \frac{|x|^2}{2}  
% .
% $$
% In turn 
% $$
% \lim_{t \to \infty} 
%  \int_{\mathbb{R}^n} e^{-\mathcal{L}\Phi_t(x)}dx 
%  =
%  (2\pi)^{n/2} \int_{\mathbb{R}^n} e^{-|x|^2/2}dx \left( \int_{\mathbb{R}^n} e^{-\Psi(x)}dx  \right)^{-1} 
%  =
% (2\pi)^{n}  \left( \int_{\mathbb{R}^n} e^{-\Psi(x)}dx  \right)^{-1}  .
% $$
% This leads to
% $$
% \lim_{t \to \infty} \alpha(t) 
% =  (2\pi)^{n}  
% $$
% therefore proving the functional form of Blaschke Santalo inequality.

% This reproduces Cordero et al argument, after rearrangement. {\color{red}Cyril: There is no real gain with respect to Cordero et al, just that we are dealing with convex functions on the positive axis only, which probably makes life easier for regularity/technicalities. 

% Also, I wonder if our approach could say anything about robustness (quantitative version of BS)?}


% \subsection{Polar Bascke Santalo}

% Let us present a tentative of proving Artstein-Avidan-Shlomka's result related to  polar Baschke Santalo. Namely, we want to maximize the following quantity
% $$
% \int e^{-f} \int e^{-f^\circ}
% =
% \int e^{-f} 
% \int \Gamma(n+1, \frac{1}{\mathcal{L}f}) .
% $$
% which, after rearrangement and change of variable, reduces to maximizing
% $$
% \int_0^\infty e^{-f(r)}r^{n-1}dr 
% \int_0^\infty \Gamma(n+1, \frac{1}{\mathcal{L}f(r)})r^{n-1}dr 
% $$
% over all convex increasing function $f \colon [0,\infty) \to [0,\infty)$ satisfying $f(0)=0$.

% We want to use the semi-group approach developped by Cordero et al. To that aim, consider the Laguerre operator
% $$
% Lf(r)=rf''(r)+(n-r)f'(r), \qquad r >0
% $$
% symmetric in $\mathbb{L}^2(\nu)$, with $\nu$ the Gamma distribution with density
% $$
% \varphi(r) \coloneqq \frac{r^{n-1}e^{-r} }{\Gamma(n)}
% $$
% on the positive axis. This means that, for any functions $f,g$ smooth enough
% $$
% \int_0^\infty fLg d\nu = \int g Lf d\nu .
% $$
% Then define the dual $M$ of $L$ in $\mathbb{L}^2(r^{n-1}dr)$. Namely, the operator acting on smooth functions as
% $$
% Mf(r) = rf''(r)+(n+r)f'(r)+nf(r), \qquad r >0
% $$
% satisfies
% $$
% \int_0^\infty f(r)Lg(r)r^{n-1}dr = \int_0^\infty g(r)Mf(r)r^{n-1}dr 
% $$
% as one can easily check using various integration by parts.

% Denote by $P_t$, $t \geq 0$ the associated Laguerre semi-group and by $Q_t$ the semi-group assocaited to $M$, and set 
% $$
% f_t = -\log Q_t(e^{-g}),
% $$
% for some fixed convex increasing function $g \colon [0,\infty) \to [0,\infty)$ satisfying $g(0)=0$. For simplicity we set
% $F_t \coloneqq \mathcal{L}f_t$ that is also convex increasing and $F_t(0)=0$. In order to avoid unnecessary technicalities we assume that $g$ is strictly convex so that $F_t$ does not take infinite values.

% Our aim is to prove that
% \begin{align*}
% \alpha(t) 
% &\coloneqq
% \int_0^\infty e^{-f_t(r)}r^{n-1}dr 
% \int_0^\infty \Gamma(n+1, \frac{1}{F_t})r^{n-1}dr \\
% & = 
% \int_0^\infty Q_t(e^{-g})(r)r^{n-1}dr 
% \int_0^\infty \Gamma(n+1, \frac{1}{F_t})r^{n-1}dr, \qquad t >0
% \end{align*}
% is non decreasing.

% Observe that
% \begin{align*}
% \frac{d}{dt} 
% \int_0^\infty Q_t(e^{-g})(r)r^{n-1}dr
% & =
% \int_0^\infty MP_t(e^{-g})(r)r^{n-1}dr \\
% & =
% \int_0^\infty P_t(e^{-g})(r) L \mathds{1}(r) e^r d\nu(r) & \mbox{(by duality)}\\
% & =
% 0 
% \end{align*}
% so that $t \mapsto \int_0^\infty Q_t(e^{-g})(r)r^{n-1}dr$ is a constant function. 
% Differentiating in $t$, we therefore obtain that
% \begin{align*}
% \alpha'(t)
% & = 
% - \int_0^\infty e^{-g(r)}r^{n-1}dr 
% \int_0^\infty \Gamma'(n+1, \frac{1}{F_t}) \frac{\frac{d}{dt}F_t(r)}{F_t(r)^2}r^{n-1}dr
% \end{align*}
% where $\Gamma'(a,x)=\frac{d}{dx}\Gamma(a,x)=-x^{a-1}e^{-x}$ is the derivative of the incomplete gamma function.
% % Collecting the two term, we arrive at
% % $$
% % \alpha'(t)
% %  = 
% % \int_0^\infty P_t(e^{-g})(r)r^{n-1}dr 
% % \int_0^\infty \left( n \Gamma(n+1, \frac{1}{F_t})  -\Gamma'(n+1, \frac{1}{F_t}) \frac{\frac{d}{dt}F_t(r)}{F_t(r)^2} \right) r^{n-1}dr .
% % $$
% Hence the sign of $\alpha'$ is governed by 
% $$
% I \coloneqq 
% \int_0^\infty \left(   -\Gamma'(n+1, \frac{1}{F_t}) \frac{\frac{d}{dt}F_t(r)}{F_t(r)^2} \right) r^{n-1}dr
% $$
% we now analyze.
% By Lemma \ref{lem:Cyril} we observe that
% \begin{align*}
% \frac{d}{dt}F_t 
% & = - \frac{d}{dt} f_t(F_t') \\
% & = 
% -\frac{MQ_t(e^{-g})(F_t')}{Q_t(e^{-g})(F_t')} \\
% & =
% -F_t'\frac{Q_t(e^{-g})''(F_t')}{Q_t(e^{-g})(F_t')} -\left(n+F_t'\right) \frac{Q_t(e^{-g})'(F_t')}{Q_t(e^{-g})(F_t')} - n .
% \end{align*}
% \\{\color{blue}Devraj: The second line should be $$\frac{MQ_t(e^{-g})(F_t')}{Q_t(e^{-g})(F_t')}$$
% }
% However, 
% $$
% f_t'= - \frac{Q_t(e^{-g})'}{Q_t(e^{-g})}
% \qquad \mbox{and} \qquad 
% f_t''=-\frac{Q_t(e^{-g})''}{Q_t(e^{-g})} + 
% \frac{{Q_t(e^{-g})'}^2}{Q_t(e^{-g})^2} .
% $$
% Hence
% \begin{align*}
% \frac{d}{dt}F_t 
% & =
% -F_t'\left( {f_t'}^2(F_t') - f_t''(F_t') \right)
% +\left(n+F_t'\right) f_t'(F_t') -n \\
% & =
% - F_t'(r)\left( r^2 - \frac{1}{F_t''(r)} \right)
% +\left(n+F_t'(r)\right) r -n 
% \end{align*}
% since $f_t'(F_t'(r))=r$, which also leads to
% $f_t''(F_t')F_t''(r)=1$ and therefore to 
% $f_t''(F_t')=\frac{1}{F_t''}$.
% Changing notation ($V=F_t$), we obtain
% $$
% I = \int_0^\infty \frac{e^{-1/V(r)}}{V^{n+2}(r)}
% \left(  \frac{V'(r)}{V''(r)}-r^2V'(r) + nr+V'(r) r -n \right)r^{n-1}dr .
% $$
% {\color{red}Cyril: the scope would be to prove that $I \geq 0$, maybe like in the previous section showing that, after some appropriate change of variables, the bracket is nothing but a square (or anything trivially positive)...}
% \\{\color{blue}Devraj: I then end up with $$I = -\int_0^\infty \frac{e^{-1/V(r)}}{V^{n+2}(r)}
% \left(  \frac{V'(r)}{V''(r)}-r^2V'(r) + nr+V'(r) r -n \right)r^{n-1}dr$$, so with the opposite sign, evidently, the difficult tasks remains}

% Later on we will use the following trivial fact
% $$
% L e^r=ne^r .
% $$
% {\color{red}Cyril: There might be a problem here since $r \mapsto e^r$ does not belong to $\mathbb{L}^2(\nu)$.}

% Denote by $P_t$, $t \geq 0$ the associated Laguerre semi-group and set 
% $$
% f_t = -\log P_t(e^{-g}),
% $$
% for some fixed convex increasing function $g \colon [0,\infty) \to [0,\infty)$ satisfying $g(0)=0$. For simplicity we set
% $F_t \coloneqq \mathcal{L}f_t$ that is also convex increasing and $F_t(0)=0$. In order to avoid unnecessary technicalities we assume that $g$ is strictly convex so that $F_t$ does not take infinite values.

% Our aim is to prove that
% \begin{align*}
% \alpha(t) 
% &\coloneqq
% \int_0^\infty e^{-f_t(r)}r^{n-1}dr 
% \int_0^\infty \Gamma(n+1, \frac{1}{F_t})r^{n-1}dr \\
% & = 
% \int_0^\infty P_t(e^{-g})(r)r^{n-1}dr 
% \int_0^\infty \Gamma(n+1, \frac{1}{F_t})r^{n-1}dr, \qquad t >0
% \end{align*}
% is non decreasing.

% Observe that
% \begin{align*}
% \frac{d}{dt} 
% \int_0^\infty P_t(e^{-g})(r)r^{n-1}dr
% & =
% \int_0^\infty LP_t(e^{-g})(r)r^{n-1}dr \\
% & =
% \Gamma(n) \int_0^\infty LP_t(e^{-g})(r) e^r d\nu(r) \\
% & =
% \Gamma(n) \int_0^\infty P_t(e^{-g})(r) L e^r d\nu(r) & \mbox{(by symmetry)}\\
% & =
% n \int_0^\infty P_t(e^{-g})(r) r^{n-1}dr &  \mbox{(by the trivial fact above)} .
% \end{align*}
% Differentiating in $t$, we therefore obtain that
% \begin{align*}
% \alpha'(t)
% & = 
% n \alpha(t)
% - 
% \int_0^\infty P_t(e^{-g})(r)r^{n-1}dr 
% \int_0^\infty \Gamma'(n+1, \frac{1}{F_t}) \frac{\frac{d}{dt}F_t(r)}{F_t(r)^2}r^{n-1}dr
% \end{align*}
% where $\Gamma'(a,x)=\frac{d}{dx}\Gamma(a,x)=-x^{a-1}e^{-x}$ is the derivative of the incomplete gamma function.
% Collecting the two term, we arrive at
% $$
% \alpha'(t)
%  = 
% \int_0^\infty P_t(e^{-g})(r)r^{n-1}dr 
% \int_0^\infty \left( n \Gamma(n+1, \frac{1}{F_t})  -\Gamma'(n+1, \frac{1}{F_t}) \frac{\frac{d}{dt}F_t(r)}{F_t(r)^2} \right) r^{n-1}dr .
% $$
% Hence the sign of $\alpha'$ is governed by the sign of the second integral
% $$
% I \coloneqq 
% \int_0^\infty \left( n \Gamma(n+1, \frac{1}{F_t})  -\Gamma'(n+1, \frac{1}{F_t}) \frac{\frac{d}{dt}F_t(r)}{F_t(r)^2} \right) r^{n-1}dr
% $$
% we now analyze.

% By Lemma \ref{lem:Cyril} we observe that
% \begin{align*}
% \frac{d}{dt}F_t 
% & = - \frac{d}{dt} f_t(F_t') \\
% & = 
% \frac{LP_t(e^{-g})(F_t')}{P_t(e^{-g})(F_t')} \\
% & =
% F_t'\frac{P_t(e^{-g})''(F_t')}{P_t(e^{-g})(F_t')} +\left(n-F_t'\right) \frac{P_t(e^{-g})'(F_t')}{P_t(e^{-g})(F_t')} .
% \end{align*}






% \newpage

% Recall that we "proved" earlier that
% $$
%     \sup_{f\in \tilde C(\mathbb{R}^n)} 
%     \int e^{-f}\int e^{-\mathcal{L}f}
%     =
%     (n\omega_{n})^2 
%     \sup_{\Psi} 
%     \int_0^\infty e^{-\Psi(r)}r^{n-1}dr 
%     \int_0^\infty e^{-\mathcal{L}\Psi}r^{n-1}dr 
% $$
% with $\Psi$ convex, non-decreasing with $\Psi(0)=0$.
% Changing $\Psi$ into $\Psi+c$ for any constant $c$ doesn't affect the volume product, therefore the constraint $\Psi(0)=0$ might be removed if needed.

% Since we will use the heat flow on the line, we extend $\Psi$ on the negative axis by parity. Moreover by an approximation argument
% we can assume that $\Psi$ is strictly convex. Our aim is therefore to compute
% $$
% M = \sup_{\Psi} 
%     \int_{-\infty}^\infty e^{-\Psi(r)}|r|^{n-1}dr 
%     \int_{-\infty}^\infty e^{-\mathcal{L}\Psi}|r|^{n-1}dr 
% $$
% for $\Psi \colon \mathbb{R} \to [0,\infty)$  even, non-decreasing  and strictly convex on $(0,\infty)$ (whatever the value of $\Psi(0)$). 
% Then, we would get
% $$
%     \sup_{f\in \tilde C(\mathbb{R}^n)} 
%     \int e^{-f}\int e^{-\mathcal{L}f}
%     =
% \frac{1}{4}(n\omega_{n})^2 M 
% $$
% and we are left with an estimation of $M$.
% To that aim, and following Cordero et al., consider $\Psi_t \coloneqq - \log P_tf$
% where $P_t$ is the one dimensional heat flow and $f(r)=e^{-\Psi(r)}$, $r \in \mathbb{R}$ for a given $\Psi$ as above. 
% %Since we can change the value of $\Psi(0)$, we assume that $\int e^{-\Psi(r)}dr=1$, turning $f$ into a probability density on the line.
% Now, for $x \in \mathbb{R}$ {\color{red}Cyril: I'm not sure about the normalization here, the $4t$ and $\sqrt{2t}$}
% $$
% P_tf(x) \coloneqq \mathbb{E}(x+\sqrt{2t}Z) = \frac{1}{\sqrt{4\pi t}} \int_{-\infty}^\infty f(y)e^{-|x-y|^2/4t}dy
% =
% \frac{1}{\sqrt{2t+1}}
% Q_{\log\sqrt{2t+1}}f(\frac{x}{\sqrt{2t+1}})
% $$
% with $Z \sim \mathcal{N}(0,1)$ a one dimensional standard normal. It is generated by the Laplacien $L$ which, in dimension 1, takes the form $Lf=f''$ for $f$ smooth enough. Here $Q$ is the Focker-Planck semi-group, that is, the adjoint of the Ornstein-Uhlenbeck semi-group in $\mathbb{L}^2(dx)$ (for the Lebesgue measure), see \eqref{eq:Focker-Planck}. One important fact is that $\lim_{t \to \infty} Q_tf = \frac{1}{\sqrt{2\pi}}e^{-r^2/2} \int_{-\infty}^\infty f(r)dr$.
% For $t \geq 0$, consider the following quantity
% $$
% \alpha(t) \coloneqq 
%     \int_{-\infty}^\infty e^{-\Psi_t(r)}|r|^{n-1}dr 
%     \int_{-\infty}^\infty e^{-\mathcal{L}\Psi_t}|r|^{n-1}dr 
%     =
%     \int_{-\infty}^\infty P_tf(r)|r|^{n-1}dr 
%     \int_{-\infty}^\infty e^{-\mathcal{L}\Psi_t}|r|^{n-1}dr
% $$
% and observe that 
% $$
% \alpha(0)= \int_{-\infty}^\infty e^{-\Psi(r)}|r|^{n-1}dr 
%     \int_{-\infty}^\infty e^{-\mathcal{L}\Psi}|r|^{n-1}dr
% $$
% while, (after change of variable $s=r/\sqrt{2t+1}$), {\color{red}Cyril: and some convergence theorem that I'm skipping but that should be controlled properly}
% \begin{align*}
% \lim_{t \to \infty} \alpha(t)
% & = \int_{-\infty}^\infty Q_\infty f(s) |s|^{n-1}ds
%    \lim_{t \to \infty}  \int_{-\infty}^\infty e^{-\mathcal{L}\Psi_t}|r|^{n-1}dr \\
%    & =
%   \frac{1}{\sqrt{2\pi}}\int_{-\infty}^\infty e^{-s^2/2} |s|^{n-1}ds \int_{-\infty}^\infty f(r)dr   \int_{-\infty}^\infty \sqrt{2\pi}e^{-s^2/2} |s|^{n-1}ds \left( \int_{-\infty}^\infty f(r)dr \right)^{-1} \\
%   & = 
%   4 (2^{\frac{n}{2}-1} \Gamma(n/2))^2 
% \end{align*}
% where we used, in the second line, similar computations than in \eqref{eq:Devraj1}, \eqref{eq:Devraj2}, \eqref{eq:Devraj3}.

% Assume one can prove that $\alpha$ is non-decreasing, one would get that 
% $M \leq 4 (2^{\frac{n}{2}-1} \Gamma(n/2))^2$, and therefore, after some algebra, the functional Blashke Santalo inequality
% $$
%  \sup_{f\in \tilde C(\mathbb{R}^n)} 
%     \int e^{-f}\int e^{-\mathcal{L}f}
%     = (2 \pi)^n .
% $$
% Therefore, to fully recover the proof by Cordero et al. it remains to prove that $\alpha$ is non decreasing.


% %Using Brascamp-Lieb Cordero et al prove that $\Phi_t''>0$
% %{\color{red}Cyril: computations are in the appendix of their paper, would be better to have a self-contained easy argument of this fact}. 
 
% {\color{red}Cyril: what is about to come is essentially wrong. $r^{n-1}$ is the bad guy. I don't know how to deal with it so far.}
 
% % Therefore $\Phi_t'$ has an inverse function we denote ${\Phi_t'}^{-1}$ and it is easy to check that
% % ${\Phi_t'}^{-1}=(\mathcal{L}\Phi_t)'$ and
% % \begin{align*}
% % \mathcal{L}\Phi_t (x)
% % & = x {\Phi_t'}^{-1}(x) - \Phi_t({\Phi_t'}^{-1}(x)) \\
% % & = x (\mathcal{L}\Phi_t)'(x) - \Phi_t((\mathcal{L}\Phi_t)'(x))
% % , \qquad x \in \mathbb{R}.
% %\end{align*}
% % Taking the derivative in $t$, it holds
% % $$
% % \partial_t \mathcal{L}\Phi_t (x) = 
% % x \partial_t (\mathcal{L}\Phi_t)'(x)
% % -\partial_t \Phi_t((\mathcal{L}\Phi_t)'(x))
% % - \Phi_t'((\mathcal{L}\Phi_t)'(x))\partial_t \mathcal{L}\Phi_t)'(x) .
% % $$
% % Now the identity ${\Phi_t'}((\mathcal{L}\Phi_t)'(x))=x$ leads to
% % $$
% % \partial_t \mathcal{L}\Phi_t (x) = 
% % -\partial_t \Phi_t((\mathcal{L}\Phi_t)'(x)) ,
% % \qquad x \in \mathbb{R}, \quad t >0  .
% %  $$
% For the heat flow,
% \begin{align*}
%     \partial_t \Phi_t
%     =
%     - \frac{LP_t(f)}{P_t(f)}
%     =
%     - \frac{P_t(f)''}{P_t(f)}
%     = 
%     \Phi_t'' - {\Phi_t'}^2 .
% \end{align*}
% We conclude that
% \begin{align*}
% \partial_t \mathcal{L}\Phi_t (x) 
% & = 
% -\Phi_t''((\mathcal{L}\Phi_t)'(x)) + {\Phi_t'}^2 ((\mathcal{L}\Phi_t)'(x)) \\
% & =
% -\Phi_t''({\Phi_t'}^{-1}(x)) + x^2 \\
% & =
% -\frac{1}{(\mathcal{L}\Phi_t)''(x)} + x^2 .
% \end{align*}
% Now, given $\Psi$ set
% $$
% \alpha(t) = \int_{-\infty}^\infty e^{-\Phi_t(r)}|r|^{n-1}dr 
%     \int_{-\infty}^\infty e^{-\mathcal{L}\Phi_t}|r|^{n-1}dr
% $$
% By construction of $\Phi_t$, and since the heat semi-group is invariant for the Lebesgue measure, in fact
% $$
% \alpha(t) = 
% \int_{-\infty}^\infty e^{-\Psi(r)}|r|^{n-1}dr 
%     \int_{-\infty}^\infty e^{-\mathcal{L}\Phi_t}|r|^{n-1}dr
% $$
% Therefore
% \begin{align*}
% \alpha'(t)
% & =
% - \int_{-\infty}^\infty e^{-\Psi(r)}|r|^{n-1}dr 
%     \int_{-\infty}^\infty \partial_t \mathcal{L}\Phi_t (x)  e^{-\mathcal{L}\Phi_t}|r|^{n-1}dr \\
%     & =
% \int_{-\infty}^\infty e^{-\Psi(r)}|r|^{n-1}dr 
%     \int_{-\infty}^\infty \left( \frac{1}{(\mathcal{L}\Phi_t)''(x)} - x^2 \right)  e^{-\mathcal{L}\Phi_t}|r|^{n-1}dr .
% \end{align*}


% \newpage

% {\color{red}Cyril: what comes next is Devraj's atempt to proving the polar functional Blashle-Santalo.}

% In this section we would like to show that $$
% \int_{\mathbb{{R}}^{n}}e^{-\phi(x)}dx\int_{\mathbb{{R}}^{n}}e^{-\phi(x)^\circ}dx\leq\int_{\mathbb{R}^{n}}e^{-|x|}dx\int_{\mathbb{{R}}^{n}}e^{-|x|^\circ}dx=\left[|S_{n-1}|(n-1)!\right]^{2}
% $$ where $\phi\in\Tilde{C}(\mathbb{R}^{n})$
% Using arguments shown above we have the following 
% \begin{equation}
%     \begin{split}
%         &\sup_{\phi\in\Tilde{C}(\mathbb{R}^{n})}\int_{\mathbb{{R}}^{n}}e^{-\phi(x)}dx\int_{\mathbb{{R}}^{n}}e^{-\phi(x)^\circ}dx\\
%         \leq&\sup_{\phi\in\Tilde{C}(\mathbb{R}^{n})}\int_{\mathbb{{R}}^{n}}e^{-\phi(|x|)}dx\int_{\mathbb{{R}}^{n}}e^{-\phi(|x|)^\circ}dx\\
%         \leq&\sup_{\phi\in\Tilde{C}(\mathbb{R}^{n})}|S_{n-1}|^2\int_{0}^{\infty}e^{-\phi(r)}r^{n-1}dr\int_{0}^{\infty}e^{-\phi(r)^\circ}r^{n-1}dr\\
%         =&\sup_{\phi\in\Tilde{C}(\mathbb{R}^{n}),\int_{0}^{\infty}e^{-\phi(r)}dr=1}|S_{n-1}|^2\int_{0}^{\infty}e^{-\phi(r)}r^{n-1}dr\int_{0}^{\infty}e^{-\phi(r)^\circ}r^{n-1}dr\\
%     \end{split}
% \end{equation}
% Note that if $\phi\in\Tilde{C}(\mathbb{R}^{n})$ and $\int_{0}^{\infty}e^{-\phi(r)}dr=1$ we must have that $$
% \int_{0}^{\infty}e^{-\phi(r)}r^{n-1}dr\leq\int_{0}^{\infty}e^{-r}r^{n-1}dr\leq(n-1)!
% $$ This gives the following $$
% \sup_{\phi\in\Tilde{C}(\mathbb{R}^{n})}\int_{\mathbb{{R}}^{n}}e^{-\phi(x)}dx\int_{\mathbb{{R}}^{n}}e^{-\phi(x)^\circ}dx\leq\sup_{\phi\in\Tilde{C}(\mathbb{R}^{n}),\int_{0}^{\infty}e^{-\phi(r)}dr=1}|S_{n-1}|^{2}(n-1)!\int_{0}^{\infty}e^{-\phi(r)^\circ}r^{n-1}dr
% $$ It is rather evident that one much show that $$
% \sup_{\phi\in\Tilde{C}(\mathbb{R}^{n}),\int_{0}^{\infty}e^{-\phi(r)}dr=1}\int_{0}^{\infty}e^{-\phi(r)^\circ}r^{n-1}dr\leq (n-1)!
% $$ To simplify notation, for $f=e^{-\phi}$, define $$
% M(f):=\int_{0}^{\infty}exp\left[-\left(-log f\right)^\circ\right]r^{n-1}dr
% $$ To this aim, we introduce the Laguerre semigroup $P_{t}$ that is associated to the one sided exponential measure $\lambda,\frac{d\lambda}{dx}=e^{-x}, x>0$. We need the dual semigroup (with respect to the Lebesgue measure on the half line) $\Tilde{P}_{t}$ of this semigroup which satisfies the following $$
% \int_{0}^{\infty}g(x)P_{t}f(x)dx=\int_{0}^{\infty}f(x)\Tilde{P}_{t}g(x)dx
% $$ From the above, we see that $$
% \Tilde{P}_{0}f=f, \Tilde{P}_{\infty}f=e^{-r}\int_{0}^{\infty}f(r)dr
% $$ Note that $\Tilde{P}_{t}f(x)$ is the solution to the following differential equation $$
% \partial_{t}q(t,x)=q(t,x)+(x+1)\frac{d}{dx}q(t,x)+x\frac{d^2}{dx^2}q(t,x),q(0,x)=f(x)
% $$ We need to show that $\alpha(t)=log M(\Tilde{P}_{t}f)$ is non-decreasing in time for $-log(f)\in\Tilde{C}(R^{n}),\int_{0}^{\infty}f(r)dr=1$. This would yield our result. $$
% \alpha'(t)=\int_{0}^{\infty}-\frac{d}{dt}[-log \Tilde{P}_{t}(f)(r)]^\circ\frac{\mu_{t}}{\int_{0}^{\infty}\mu_{t}(r)dr}dr
% $$ where $$
% \mu_{t}(r)=exp\left[-(-log \Tilde{P}_{t}f)^\circ\right]r^{n-1}
% $$ 
% {\color{blue}Devraj:The above is precisely what was done in the article by Cordero et Al. Regrettably, we've to work with this strange semigroup as the scaling argument won't work well with the polar transform (meaning we don't have $(h+c)^\circ=h^\circ-c$ in stark contrast to the Legendre transform. To follow their article, we would have to show \begin{enumerate}
%     \item $\mu_{t}$ is strictly log concave (so we can use some brascamp-lieb inequality). For this it may be better to use the relationship with the Orstein-Uhlenbeck semigroup as things are much more tractable. 
%     \item Compute $\frac{d}{dt}[-log \Tilde{P}_{t}(f)(r)]^\circ$ to see if for some choice of function in the brascamp lieb we can bound this by something nice.
% \end{enumerate}
% But to be quite honest...the computations are proving to be rather atrocious :( 
% }
















\bibliographystyle{plain}
\bibliography{rearrangement}






\end{document}





As a consequence $t \mapsto \alpha(t)$
is non-decreasing so that
$$
\alpha(0) = \int_{\mathbb{R}^n} e^{-\Psi} \int e^{\mathcal{L}\Psi}
\leq 
\lim_{t \to \infty} \alpha(t) 
$$
and we are left with computing 
\begin{align*}
\lim_{t \to \infty} \alpha(t) 
& =
\lim_{t \to \infty} (n\omega_n)^2    \int_0^\infty Q_t(e^{-\Psi})(r)r^{n-1}dr  
    \int_0^\infty e^{-\mathcal{L}\Phi_t(r)}r^{n-1}dr \\
&    =
\int_{\mathbb{R}^n} e^{-\Psi}(x)dx  
\lim_{t \to \infty}  \int_{\mathbb{R}^n} e^{-\mathcal{L}\Phi_t(x)}dx 
\end{align*}
where in the last line we used \eqref{eq:0}.

By construction, when $t$ tends to infinity, $Q_tg(r) \to n\omega_n \frac{e^{-r^2/2}}{(2\pi)^{n/2}} \int_0^\infty g(r)r^{n-1}dr$ (to see this one can use the relation $\int_0^\infty f(r)Q_t g(r) r^{n-1}dr = \int_0^\infty g(r)P_tf(r)r^{n-1}dr$ with $f=\delta_{r_o}$,
the Dirac mass at some fixed point $r_o \in (0,\infty)$, and the fact that $P_tf \to n\omega_n \int f(s)e^{-s^2/2}s^{n-1}dr /(2\pi)^{n/2} = n\omega_n e^{-r_o^2/2}r_o^{n-1}/(2\pi)^{n/2}$ in the limit $t \to \infty$). 

Hence,
$$
\lim_{t \to \infty} \Phi_t (x)
=
- \log \left( n\omega_n \frac{e^{-|x|^2/2}}{(2\pi)^{n/2}} \int_0^\infty e^{-\Psi(r)}r^{n-1}dr \right)
=
- \log \left( \frac{1}{(2\pi)^{n/2}} \int_{\mathbb{R}^n} e^{-\Psi(x)}dx \right)
+ \frac{|x|^2}{2}
$$
so that {(\color{red}Cyril: again there are some technical issues here)}
$$
\lim_{t \to \infty}  \mathcal{L}\Phi_t(x)
=
\log \left( \frac{1}{(2\pi)^{n/2}} \int_{\mathbb{R}^n} e^{-\Psi(x)}dx  \right) 
+
\frac{|x|^2}{2}  
.
$$
In turn 
$$
\lim_{t \to \infty} 
 \int_{\mathbb{R}^n} e^{-\mathcal{L}\Phi_t(x)}dx 
 =
 (2\pi)^{n/2} \int_{\mathbb{R}^n} e^{-|x|^2/2}dx \left( \int_{\mathbb{R}^n} e^{-\Psi(x)}dx  \right)^{-1} 
 =
(2\pi)^{n}  \left( \int_{\mathbb{R}^n} e^{-\Psi(x)}dx  \right)^{-1}  .
$$
This leads to
$$
\lim_{t \to \infty} \alpha(t) 
=  (2\pi)^{n}  
$$
therefore proving the functional form of Blaschke Santalo inequality.

This reproduces Cordero et al argument, after rearrangement. {\color{red}Cyril: There is no real gain with respect to Cordero et al, just that we are dealing with convex functions on the positive axis only, which probably makes life easier for regularity/technicalities. 

Also, I wonder if our approach could say anything about robustness (quantitative version of BS)?}





\subsubsection{Reproducing Cordero et al after rearrangement}

As an illustration of the rearrangement techniques, let us reproduce the proof of Cordero et al of the functional Blaschke Santalo.  

We start with an obvious lemma. {(\color{red}Cyril: which however should be stated carrefully, there are some technical issues here.)}

\begin{lemma} \label{lem:Cyril}
  Let $(f_t)_{t \geq 0}$ be a family of $\mathcal{C}^1$ strictly convex functions on the line taking finite values. Assume that, for any $x \in \mathbb{R}$, $t \mapsto f_t(x)$ is differentiable. Then for all $t \geq 0$, $f_t'$ is invertible with inverse $(\mathcal{L}f_t)'$. Moreover,
  $\mathcal{L}f_t(x)=x(\mathcal{L}f_t)'(x)-f_t((\mathcal{L}f_t)'(x))$ and $\partial_t \mathcal{L}f_t(x) = - \partial_tf_t((\mathcal{L}f)'(x))$. Here $\partial_t$ is a shorthand notation for the partial derivative with respect to the variable $t$ while
  prime denotes the derivative with respect to $x$. 
  In particular, for $f_t=-\log P_tg$ for some semi-group $(P_t)_{t \geq0}$ with infinitesimal generator $L$, and some function $g$, $\partial_t \mathcal{L}f_t(x) = \frac{LP_tg(y)}{P_tg(y)}|_{y=(\mathcal{L}f_t)'(x)}$, $x \in \mathbb{R}$. 
 \end{lemma}

\begin{proof}
  That  $f_t'$ is invertible with inverse $(\mathcal{L}f_t)'$ and
  $\mathcal{L}f_t(x)=x(\mathcal{L}f_t)'(x)-f((\mathcal{L}f)'(x))$
  is clear. Taking the derivative with respect to $t$, we obtain
$$
\partial_t \mathcal{L}f_t (x) = 
x \partial_t (\mathcal{L}f_t)'(x)
-\partial_t f_t((\mathcal{L}f_t)'(x))
- f_t'((\mathcal{L}f_t)'(x))\partial_t (\mathcal{L}f_t)'(x) .
$$
The identity ${f_t'}((\mathcal{L}f_t)'(x))=x$ leads to
$\partial_t \mathcal{L}f_t (x) = 
-\partial_t f_t((\mathcal{L}f_t)'(x))$
  as expected.
\end{proof}

% Recall that we "proved" earlier that
% $$
%     \sup_{f\in \tilde C(\mathbb{R}^n)} 
%     \int e^{-f}\int e^{-\mathcal{L}f}
%     =
%     \sup_{\Psi} 
%     \int e^{-\Psi(|x|)}dx 
%     \int_0^\infty e^{-\mathcal{L}\Psi(|x|}dx  
% $$
% with $\Psi$ convex, increasing with $\Psi(0)=0$.
% %Changing $\Psi$ into $\Psi+c$ for any constant $c$ doesn't affect the volume product, therefore the constraint $\Psi(0)=0$ might be removed if needed.

% In Cordero et al they consider the heat semi-group $P_t^h$, with generator the Laplacian $L^h=\Delta$ and put $f=-\log P_t^h(e^{-g})$. In view of the rearrangement procedure, it is natural to consider the projection of the heat-semi-group on its radial part. This is known as the Bessel semi-group we denote $P_t^b$, and whose infinitesimal generator is acting on smooth functions $\Psi \colon (0,\infty) \to \mathbb{R}$ as

% $$
% L^b \Psi (r) = \Psi''(r) + \frac{n-1}{r}\Psi'(r), \qquad r \in (0,\infty) .
% $$
% In fact it is not difficult to see that, for $f(x)=\Psi(|x|)$,
% $P_t^hf(x)=P_t^b \Psi(|x|)$. The invariant measure for the Bessel semi-group is the measure with density $r^{n-1}$ with respect to the Lebesgue measure on $(0,\infty)$.

% Therefore, the natural counterpart of the choice $f_t=-\log P_t^h(e^{-g})$ is 
% $$
% \Phi_t(x)=-\log P_t^b(e^{-\Psi})(|x|), 
% $$
% for $\Psi$ defined on the positive axis, convex increasing, with $\Psi(0)=0$. 

% Set, passing to polar coordinates and by invariance,
% \begin{align*}
% \alpha(t)
% & \coloneqq
% \int_{\mathbb{R}^n} e^{-\Phi_t(x)}dx 
%     \int_{\mathbb{R}^n} e^{-\mathcal{L}\Phi_t(x)}dx \\
% & =
% \int_{\mathbb{R}^n} P_t^b(e^{-\Psi})(|x|) dx 
%     \int_{\mathbb{R}^n} e^{-\mathcal{L}\Phi_t(|x|)}dx \\
%     & =
% (n\omega_n)^2    \int_0^\infty P_t^b(e^{-\Psi})(r)r^{n-1}dr  
%     \int_0^\infty e^{-\mathcal{L}\Phi_t(r)}r^{n-1}dr \\
%     & =
%    (n\omega_n)^2    \int_0^\infty e^{-\Psi(r)}r^{n-1}dr  
%     \int_0^\infty e^{-\mathcal{L}\Phi_t(r)}r^{n-1}dr  
% \end{align*}
% Our aim is to prove that $\alpha$ is non-decreasing. To see this, take the derivative 
% $$
% \alpha'(t) = 
% - (n\omega_n)^2 \int_0^\infty e^{-\Psi(r)}r^{n-1}dr  
%     \int_0^\infty \partial_t \mathcal{L}\Phi_t(r) e^{-\mathcal{L}\Phi_t(r)}r^{n-1}dr .
% $$
% {\color{red}Cyril: In the language of Cordero et al the latter is written for a generic $f_t=-\log P_t^h(e^{-g})$ in $\mathbb{R}^n$ (not necessarily spherically symmetric) and reads
% $$
% \alpha'(t)
% =
% - \int_{\mathbb{R}^n} e^{-g(x)}dx 
%    \int_{\mathbb{R}^n} \partial_t \mathcal{L}f_t(x) e^{-\mathcal{L}f_t(x)}dx
% $$
% In fact both expressions are equal and one can pass from one to the other by changing variables.}

% By Lemma \ref{lem:Cyril} it holds
% $\partial_t \mathcal{L}\Phi_t(r) = - \partial_t\Phi_t((\mathcal{L}\Phi_t)'(r))$. Now for the Bessel semi-group, (since $\Phi_t'=-P_t^b(e^{-\Psi})'/P_t^b(e^{-\Psi})$ and
% $\Phi_t''=-P_t^b(e^{-\Psi})''/P_t^b(e^{-\Psi})+(P_t^b(e^{-\Psi})'/P_t^b(e^{-\Psi}))^2$)
% \begin{align*}
%     \partial_t\Phi_t (r)
%     & = 
%     -\frac{L^b P_t^b(e^{-\Psi})(r)}{P_t^b(e^{-\Psi})(r)} \\
%     & =
%     \frac{-1}{P_t^b(e^{-\Psi})(r)} 
%     \left( P_t^b(e^{-\Psi})''(r) + \frac{n-1}{r} P_t^b(e^{-\Psi})'(r)\right) \\
%     & =
%     -\Phi_t'(r)^2 + \Phi_t''(r) + \frac{n-1}{r} \Phi_t'(r) .
% \end{align*}
% In turn
% \begin{align*}
% \partial_t \mathcal{L}\Phi_t(r) 
% & = - \partial_t\Phi_t((\mathcal{L}\Phi_t)'(r)) \\
% & = 
% - \partial_t\Phi_t({\Phi_t'}^{-1}(r)) \\
% & =
% r^2 - \Phi_t''({\Phi_t'}^{-1}(r)) - \frac{n-1}{{\Phi_t'}^{-1}(r)}r , \qquad r >0 .
% \end{align*}
% Now $\Phi_t'( {\mathcal{L}\Phi_t'}(r))=r$ guarantees that
% $\Phi_t''({\mathcal{L}\Phi_t'}(r)){\mathcal{L}\Phi_t''}(r)=1$, so that, since ${\mathcal{L}\Phi_t'}={\Phi_t'}^{-1}$,
% $\Phi_t''({\Phi_t'}^{-1}(r))=\frac{1}{(\mathcal{L}\Phi_t)''(r)}$.
% It follows that
% $$
% \partial_t \mathcal{L}\Phi_t(r) 
% =
% -\frac{1}{(\mathcal{L}\Phi_t)''(r)} + r^2 - \frac{n-1}{{\Phi_t'}^{-1}(r)}r , \qquad r >0 .
% $$
% {\color{red}Cyril: The same computation is done in Cordero et al, leading, for the heat flow, to
% $$
% \partial_t \mathcal{L}\Phi_t(x) 
% =
% - Tr(D^2_x \mathcal{L}\Phi_t)^{-1} + |x|^2 ,
% \qquad x \in \mathbb{R}^n .
% $$
% The extra term in front of $n-1$ is precisely coming from the change in polar coordinates.

% This suggests the following fact that looks non trivial to me:
% for $f(x)=g(|x|)$ (recall that ${g'}^{-1}=(\mathcal{L}g)'$)
% $$
% Tr(D^2_x \mathcal{L}f)^{-1} = 
% \frac{1}{(\mathcal{L}g)''(|x|)} + \frac{n-1}{(\mathcal{L}g)'(|x|)}|x|, \qquad x \in \mathbb{R}^n .
% $$
% }

% It follows that $\alpha'(t) \geq 0$ if we can prove that
% $$
% \int_0^\infty \left(\frac{1}{(\mathcal{L}\Phi_t)''(r)} - r^2 + \frac{n-1}{{\Phi_t'}^{-1}(r)}r \right) e^{-(\mathcal{L}\Phi_t)(r)}r^{n-1}dr \geq 0  
% $$
% {\color{red}Cyril: while Cordero et al needed to prove that
% $$
% \int_{\mathbb{R}^n} \left( Tr(D^2_x \mathcal{L}\Phi_t)^{-1} - |x|^2\right) e^{-\mathcal{L}\Phi_t(x)} dx 
% \geq 0
% $$
% which is proven by Brascamp-Lieb applied to the function $x \mapsto x_i$ and summing up over $i$.}

% In order to reproduce the argument by Cordero et al, we would need to prove that, since ${\Phi'_t}^{-1}=(\mathcal{L}\Phi)'$,
% $$
% \int_0^\infty \left(\frac{1}{(\mathcal{L}\Phi_t)''(r)} -r^2 + \frac{n-1}{(\mathcal{L}\Phi_t)'(r)}r \right) e^{-(\mathcal{L}\Phi_t)(r)}r^{n-1}dr \geq 0 .
% $$
% {\color{red}Cyril: whats is about to come should be made rigorous. In particular, we need some properties on $V:=\mathcal{L}\Phi_t$.}
% Set $V=\mathcal{L}\Phi_t$. Then the latter is recast
% $$
% \int_0^\infty \left(\frac{1}{V''(r)} -r^2 + \frac{n-1}{V'(r)}r \right) e^{-V(r)}r^{n-1}dr \geq 0 .
% $$
% Integrating by part, we have {(\color{red}Cyril: we need $V'>0$ here)}
% \begin{align*}
% -\int e^{-V(r)}r^{n+1}dr
% & =
% \int -V'(r)e^{-V(r)} \frac{r^{n+1}}{V'(r)}dr \\
% & =
% \left[ e^{-V(r)}r^{n+1}/V'(r)\right]_0^\infty 
% - \int_0^\infty \left(
% \frac{(n+1)r^n}{V'(r)} - \frac{V''(r)r^{n+1}}{V'(r)^2} \right)
% e^{-V(r)} dr\\
% & =
% - \int_0^\infty 
% \frac{(n-1)r}{V'(r)} e^{-V(r)} r^{n-1}dr
% - \int_0^\infty \left(
% \frac{2r}{V'(r)} - \frac{V''(r)r^2}{V'(r)^2} \right)
% e^{-V(r)} r^{n-1}dr .
% \end{align*}
% {(\color{red}Cyril: I used that the first term (bracket term) vanishes. To be checked.)}
% Therefore,
% \begin{align*}
%   \int_0^\infty \left(\frac{1}{V''(r)} -r^2 + \frac{n-1}{V'(r)}r \right) e^{-V(r)}r^{n-1}dr
%   & =
%   \int_0^\infty \left(\frac{1}{V''(r)} -\frac{2r}{V'(r)} + \frac{V''(r)r^2}{V'(r)^2} \right) e^{-V(r)}r^{n-1}dr \\
%   & =
%   \int_0^\infty \frac{(rV''(r)-V'(r))^2}{V''(r)V'(r)^2} e^{-V(r)}r^{n-1}dr \\
%   & \geq 0
% \end{align*}
% as expected {(\color{red}Cyril: we can assume $V$ strictly convex for example here. i.e. we need to check that $\Psi$ convex increasing implies $V=\mathcal{L}\Phi_t$ is strictly convex increasing.)}
% This reproduces Cordero et al argument, after rearrangement. {\color{red}Cyril: There is no real gain with respect to Cordero et al, just that we are dealing with convex functions on the positive axis only, which probably makes life easier for regularity/technicalities.}


% \bigskip 

% \begin{center}
% \underline{\hspace{5cm}}
% \end{center}

% {\color{red}Cyril: I don't really like the fact that, in the proof of Cordero et al, we need, in the end, to rescale, which, at the level of the semi-groups means to go from the heat flow to the Focker Planck flow. 

% In this section I would like to use and introduce only one semi-group in order to avoid rescalings etc. The natural candidate is the Focker Planck semi-group, after passing to polar coordinate. I now reproduce the proof above without details (there are some technical issues I guess one would need to fix if one wants to make the argument rigorous. The idea is to see what could be implemented for the polar version of Baschke Santalo functional inequality: we could try to mimick the proof below with the Laguerre semi-group).
% }

% On $(0,\infty)$, define the following generator (the projection of the Ornstein-Uhlenbeck semi-group on radial functions),
% $$
% Lf(r)=f''(r) +(\frac{n-1}{r} - r)f'(r), \qquad r>0
% $$
% where $f \colon (0,\infty) \to \mathbb{R}$ is smooth enough. We denote $(P_t)_{t \geq 0}$
% its associated semi-group and observe that the reversible measure is $\nu(dr)=n\omega_n r^{n-1}e^{-r^2/2}/(2 \pi)^{n/2}dr$, $r \in (0,\infty)$.

% Next, we introduce the "Polar Focker-Franck" operator, that is the dual of $L$ in 
% $\mathbb{L}^2(r^{n-1}dr)$. Namely the operator $M$ satisfying for all $f,g  \colon (0,\infty) \to \mathbb{R}$ smooth enough,
% $$
% \int_0^\infty f(r)Lg(r) r^{n-1}dr 
% = 
% \int_0^\infty g(r)Mf(r) r^{n-1}dr .
% $$
% Integrating by parts one easily gets
% $$
% Mf(r)=f''(r) + (\frac{n-1}{r} + r)f'(r) + nf(r), \qquad r >0 .
% $$
% The associated semi-group we denote $(Q_t)_{t \geq 0}$.
% Consider
% $$
% \Phi_t(x)=-\log Q_t(e^{-\Psi})(|x|), 
% $$
% for $\Psi$ defined on the positive axis, convex increasing, with $\Psi(0)=0$. 
% Set, passing to polar coordinates,
% \begin{align*}
% \alpha(t)
% & \coloneqq
% \int_{\mathbb{R}^n} e^{-\Phi_t(x)}dx 
%     \int_{\mathbb{R}^n} e^{-\mathcal{L}\Phi_t(x)}dx \\
% & =
% \int_{\mathbb{R}^n} Q_t(e^{-\Psi})(|x|) dx 
%     \int_{\mathbb{R}^n} e^{-\mathcal{L}\Phi_t(|x|)}dx \\
%     & =
% (n\omega_n)^2    \int_0^\infty Q_t(e^{-\Psi})(r)r^{n-1}dr  
%     \int_0^\infty e^{-\mathcal{L}\Phi_t(r)}r^{n-1}dr 
%     %\\
%     %& =
%    %(n\omega_n)^2    \int_0^\infty e^{-\Psi(r)}r^{n-1}dr  
%    % \int_0^\infty e^{-\mathcal{L}\Phi_t(r)}r^{n-1}dr  
% \end{align*}
% Our aim is to prove that $\alpha$ is non-decreasing. To see this, observe first that
% (by duality)
% \begin{equation} \label{eq:0}
% \frac{d}{dt} \int_0^\infty Q_t(e^{-\Psi})(r)r^{n-1}dr 
% =
% \int_0^\infty M Q_t(e^{-\Psi})(r)r^{n-1}dr
% =
% \int_0^\infty Q_t(e^{-\Psi})(r) L \mathds{1} (r)r^{n-1}dr
% = 0 .
% \end{equation}
% Therefore
% $$
% \alpha'(t) = 
% - (n\omega_n)^2 \int_0^\infty e^{-\Psi(r)}r^{n-1}dr  
%     \int_0^\infty \partial_t \mathcal{L}\Phi_t(r) e^{-\mathcal{L}\Phi_t(r)}r^{n-1}dr .
% $$
% By Lemma \ref{lem:Cyril} it holds
% (since $\Phi_t'=-Q_t(e^{-\Psi})'/Q_t(e^{-\Psi})$ and
% $\Phi_t''=-Q_t(e^{-\Psi})''/Q_t(e^{-\Psi})+(Q_t(e^{-\Psi})'/Q_t(e^{-\Psi}))^2$), for $s:=(\mathcal{L}\Phi_t)'(r)$,
% \begin{align*}
%    \partial_t \mathcal{L}\Phi_t(r) 
%    & = 
%     \frac{M Q_t(e^{-\Psi})(s)}{Q_t(e^{-\Psi})(s)} \\
%     & =
%     \frac{Q_t(e^{-\Psi})''(s)}{Q_t(e^{-\Psi})(s)}
%     + \left( \frac{n-1}{s} + s \right) \frac{Q_t(e^{-\Psi})'(s)}{Q_t(e^{-\Psi})(s)}
%     + n \\
%     & =
%     \Phi_t'(s)^2 - \Phi_t''(r) -\left( \frac{n-1}{s}+s \right) \Phi_t'(s) +n  .
% \end{align*}
% In turn, since $(\mathcal{L}\Phi_t)'(r)={\Phi_t'}^{-1}(r)$, 
% \begin{align*}
% \partial_t \mathcal{L}\Phi_t(r) 
% & =
% r^2 - \Phi_t''({\Phi_t'}^{-1}(r)) -\left( \frac{n-1}{(\mathcal{L}\Phi_t)'(r)}+(\mathcal{L}\Phi_t)'(r) \right) r + n , \qquad r >0 .
% \end{align*}
% Now $\Phi_t'( {\mathcal{L}\Phi_t'}(r))=r$ guarantees that
% $\Phi_t''({\mathcal{L}\Phi_t'}(r)){\mathcal{L}\Phi_t''}(r)=1$, so that
% $\Phi_t''({\Phi_t'}^{-1}(r))=\frac{1}{(\mathcal{L}\Phi_t)''(r)}$.
% It follows that
% $$
% \partial_t \mathcal{L}\Phi_t(r) 
% =
% r^2 - \frac{1}{(\mathcal{L}\Phi_t)''(r)} -\left( \frac{n-1}{(\mathcal{L}\Phi_t)'(r)}+(\mathcal{L}\Phi_t)'(r) \right)r + n  , \qquad r >0 .
% $$
% It follows that $\alpha'(t) \geq 0$ if we can prove that
% $$
% \int_0^\infty \left(\frac{1}{(\mathcal{L}\Phi_t)''(r)} - r^2 + \left( \frac{n-1}{(\mathcal{L}\Phi_t)'(r)}+(\mathcal{L}\Phi_t)'(r) \right) - n  \right) e^{-(\mathcal{L}\Phi_t)(r)}r^{n-1}dr \geq 0  
% $$

% Set $V=\mathcal{L}\Phi_t$ for simplicity. Then the latter is recast
% $$
% \int_0^\infty \left(\frac{1}{V''(r)} -r^2 + \left( \frac{n-1}{V'(r)}+V'(r) \right)r  - n\right) e^{-V(r)}r^{n-1}dr \geq 0 .
% $$
% Integrating by part, we have {(\color{red}Cyril: we need $V'>0$ here)}
% \begin{align*}
% -\int e^{-V(r)}r^{n+1}dr
% & =
% \int -V'(r)e^{-V(r)} \frac{r^{n+1}}{V'(r)}dr \\
% & =
% \left[ e^{-V(r)}r^{n+1}/V'(r)\right]_0^\infty 
% - \int_0^\infty \left(
% \frac{(n+1)r^n}{V'(r)} - \frac{V''(r)r^{n+1}}{V'(r)^2} \right)
% e^{-V(r)} dr\\
% & =
% - \int_0^\infty 
% \frac{(n-1)r}{V'(r)} e^{-V(r)} r^{n-1}dr
% - \int_0^\infty \left(
% \frac{2r}{V'(r)} - \frac{V''(r)r^2}{V'(r)^2} \right)
% e^{-V(r)} r^{n-1}dr .
% \end{align*}
% {(\color{red}Cyril: I used that the first term (bracket term) vanishes. To be checked.)}
% Similarly
% \begin{align*}
% \int_0^\infty e^{-V(r)} V'(r) r^{n}dr
% & = 
% \left[ -e^{-V(r)} r^n \right]_0^\infty
% - \int_0^\infty n e^{-V(r)} r^{n-1}dr     \\
% & =
% - \int_0^\infty n e^{-V(r)} r^{n-1}dr .
% \end{align*}
% Therefore,
% \begin{align*}
%  & \int_0^\infty \left(\frac{1}{V''(r)} -r^2 + \left( \frac{n-1}{V'(r)}+V'(r) \right)r + n \right)  e^{-V(r)}r^{n-1}dr \\
%   & \qquad \qquad \qquad\qquad=
%   \int_0^\infty \left(\frac{1}{V''(r)} -\frac{2r}{V'(r)} + \frac{V''(r)r^2}{V'(r)^2} \right) e^{-V(r)}r^{n-1}dr \\
%   &  \qquad \qquad \qquad\qquad=
%   \int_0^\infty \frac{(rV''(r)-V'(r))^2}{V''(r)V'(r)^2} e^{-V(r)}r^{n-1}dr \\
%   & \qquad \qquad \qquad\qquad 
%   \geq 0
% \end{align*}
% as expected. As a consequence $t \mapsto \alpha(t)$
% is non-decreasing so that
% $$
% \alpha(0) = \int_{\mathbb{R}^n} e^{-\Psi} \int e^{\mathcal{L}\Psi}
% \leq 
% \lim_{t \to \infty} \alpha(t) 
% $$
% and we are left with computing 
% \begin{align*}
% \lim_{t \to \infty} \alpha(t) 
% & =
% \lim_{t \to \infty} (n\omega_n)^2    \int_0^\infty Q_t(e^{-\Psi})(r)r^{n-1}dr  
%     \int_0^\infty e^{-\mathcal{L}\Phi_t(r)}r^{n-1}dr \\
% &    =
% \int_{\mathbb{R}^n} e^{-\Psi}(x)dx  
% \lim_{t \to \infty}  \int_{\mathbb{R}^n} e^{-\mathcal{L}\Phi_t(x)}dx 
% \end{align*}
% where in the last line we used \eqref{eq:0}.

% By construction, when $t$ tends to infinity, $Q_tg(r) \to n\omega_n \frac{e^{-r^2/2}}{(2\pi)^{n/2}} \int_0^\infty g(r)r^{n-1}dr$ (to see this one can use the relation $\int_0^\infty f(r)Q_t g(r) r^{n-1}dr = \int_0^\infty g(r)P_tf(r)r^{n-1}$ with $f=\delta_r$,
% the Dirac mass at some fixed point $r \in (0,\infty)$, and the fact that $P_tf \to n\omega_n \int f(s)e^{-s^2/2}s^{n-1}dr /(2\pi)^{n/2} = n\omega_n e^{-r^2/2}r^{n-1}/(2\pi)^{n/2}$ in the limit $t \to \infty$). 

% Hence,
% $$
% \lim_{t \to \infty} \Phi_t (x)
% =
% - \log \left( n\omega_n \frac{e^{-|x|^2/2}}{(2\pi)^{n/2}} \int_0^\infty e^{-\Psi(r)}r^{n-1}dr \right)
% =
% - \log \left( \frac{1}{(2\pi)^{n/2}} \int_{\mathbb{R}^n} e^{-\Psi(x)}dx \right)
% + \frac{|x|^2}{2}
% $$
% so that {(\color{red}Cyril: again there are some technical issues here)}
% $$
% \lim_{t \to \infty}  \mathcal{L}\Phi_t(x)
% =
% \log \left( \frac{1}{(2\pi)^{n/2}} \int_{\mathbb{R}^n} e^{-\Psi(x)}dx  \right) 
% +
% \frac{|x|^2}{2}  
% .
% $$
% In turn 
% $$
% \lim_{t \to \infty} 
%  \int_{\mathbb{R}^n} e^{-\mathcal{L}\Phi_t(x)}dx 
%  =
%  (2\pi)^{n/2} \int_{\mathbb{R}^n} e^{-|x|^2/2}dx \left( \int_{\mathbb{R}^n} e^{-\Psi(x)}dx  \right)^{-1} 
%  =
% (2\pi)^{n}  \left( \int_{\mathbb{R}^n} e^{-\Psi(x)}dx  \right)^{-1}  .
% $$
% This leads to
% $$
% \lim_{t \to \infty} \alpha(t) 
% =  (2\pi)^{n}  
% $$
% therefore proving the functional form of Blaschke Santalo inequality.

% This reproduces Cordero et al argument, after rearrangement. {\color{red}Cyril: There is no real gain with respect to Cordero et al, just that we are dealing with convex functions on the positive axis only, which probably makes life easier for regularity/technicalities. 

% Also, I wonder if our approach could say anything about robustness (quantitative version of BS)?}


\subsection{Polar Bascke Santalo}

Let us present a tentative of proving Artstein-Avidan-Shlomka's result related to  polar Baschke Santalo. Namely, we want to maximize the following quantity
$$
\int e^{-f} \int e^{-f^\circ}
=
\int e^{-f} 
\int \Gamma(n+1, \frac{1}{\mathcal{L}f}) .
$$
which, after rearrangement and change of variable, reduces to maximizing
$$
\int_0^\infty e^{-f(r)}r^{n-1}dr 
\int_0^\infty \Gamma(n+1, \frac{1}{\mathcal{L}f(r)})r^{n-1}dr 
$$
over all convex increasing function $f \colon [0,\infty) \to [0,\infty)$ satisfying $f(0)=0$.

We want to use the semi-group approach developped by Cordero et al. To that aim, consider the Laguerre operator
$$
Lf(r)=rf''(r)+(n-r)f'(r), \qquad r >0
$$
symmetric in $\mathbb{L}^2(\nu)$, with $\nu$ the Gamma distribution with density
$$
\varphi(r) \coloneqq \frac{r^{n-1}e^{-r} }{\Gamma(n)}
$$
on the positive axis. This means that, for any functions $f,g$ smooth enough
$$
\int_0^\infty fLg d\nu = \int g Lf d\nu .
$$
Then define the dual $M$ of $L$ in $\mathbb{L}^2(r^{n-1}dr)$. Namely, the operator acting on smooth functions as
$$
Mf(r) = rf''(r)+(n+r)f'(r)+nf(r), \qquad r >0
$$
satisfies
$$
\int_0^\infty f(r)Lg(r)r^{n-1}dr = \int_0^\infty g(r)Mf(r)r^{n-1}dr 
$$
as one can easily check using various integration by parts.

Denote by $P_t$, $t \geq 0$ the associated Laguerre semi-group and by $Q_t$ the semi-group assocaited to $M$, and set 
$$
f_t = -\log Q_t(e^{-g}),
$$
for some fixed convex increasing function $g \colon [0,\infty) \to [0,\infty)$ satisfying $g(0)=0$. For simplicity we set
$F_t \coloneqq \mathcal{L}f_t$ that is also convex increasing and $F_t(0)=0$. In order to avoid unnecessary technicalities we assume that $g$ is strictly convex so that $F_t$ does not take infinite values.

Our aim is to prove that
\begin{align*}
\alpha(t) 
&\coloneqq
\int_0^\infty e^{-f_t(r)}r^{n-1}dr 
\int_0^\infty \Gamma(n+1, \frac{1}{F_t})r^{n-1}dr \\
& = 
\int_0^\infty Q_t(e^{-g})(r)r^{n-1}dr 
\int_0^\infty \Gamma(n+1, \frac{1}{F_t})r^{n-1}dr, \qquad t >0
\end{align*}
is non decreasing.

Observe that
\begin{align*}
\frac{d}{dt} 
\int_0^\infty Q_t(e^{-g})(r)r^{n-1}dr
& =
\int_0^\infty MP_t(e^{-g})(r)r^{n-1}dr \\
& =
\int_0^\infty P_t(e^{-g})(r) L \mathds{1}(r) e^r d\nu(r) & \mbox{(by duality)}\\
& =
0 
\end{align*}
so that $t \mapsto \int_0^\infty Q_t(e^{-g})(r)r^{n-1}dr$ is a constant function. 
Differentiating in $t$, we therefore obtain that
\begin{align*}
\alpha'(t)
& = 
- \int_0^\infty e^{-g(r)}r^{n-1}dr 
\int_0^\infty \Gamma'(n+1, \frac{1}{F_t}) \frac{\frac{d}{dt}F_t(r)}{F_t(r)^2}r^{n-1}dr
\end{align*}
where $\Gamma'(a,x)=\frac{d}{dx}\Gamma(a,x)=-x^{a-1}e^{-x}$ is the derivative of the incomplete gamma function.
% Collecting the two term, we arrive at
% $$
% \alpha'(t)
%  = 
% \int_0^\infty P_t(e^{-g})(r)r^{n-1}dr 
% \int_0^\infty \left( n \Gamma(n+1, \frac{1}{F_t})  -\Gamma'(n+1, \frac{1}{F_t}) \frac{\frac{d}{dt}F_t(r)}{F_t(r)^2} \right) r^{n-1}dr .
% $$
Hence the sign of $\alpha'$ is governed by 
$$
I \coloneqq 
\int_0^\infty \left(   -\Gamma'(n+1, \frac{1}{F_t}) \frac{\frac{d}{dt}F_t(r)}{F_t(r)^2} \right) r^{n-1}dr
$$
we now analyze.
By Lemma \ref{lem:Cyril} we observe that
\begin{align*}
\frac{d}{dt}F_t 
& = - \frac{d}{dt} f_t(F_t') \\
& = 
-\frac{MQ_t(e^{-g})(F_t')}{Q_t(e^{-g})(F_t')} \\
& =
-F_t'\frac{Q_t(e^{-g})''(F_t')}{Q_t(e^{-g})(F_t')} -\left(n+F_t'\right) \frac{Q_t(e^{-g})'(F_t')}{Q_t(e^{-g})(F_t')} - n .
\end{align*}
\\{\color{blue}Devraj: The second line should be $$\frac{MQ_t(e^{-g})(F_t')}{Q_t(e^{-g})(F_t')}$$
}
However, 
$$
f_t'= - \frac{Q_t(e^{-g})'}{Q_t(e^{-g})}
\qquad \mbox{and} \qquad 
f_t''=-\frac{Q_t(e^{-g})''}{Q_t(e^{-g})} + 
\frac{{Q_t(e^{-g})'}^2}{Q_t(e^{-g})^2} .
$$
Hence
\begin{align*}
\frac{d}{dt}F_t 
& =
-F_t'\left( {f_t'}^2(F_t') - f_t''(F_t') \right)
+\left(n+F_t'\right) f_t'(F_t') -n \\
& =
- F_t'(r)\left( r^2 - \frac{1}{F_t''(r)} \right)
+\left(n+F_t'(r)\right) r -n 
\end{align*}
since $f_t'(F_t'(r))=r$, which also leads to
$f_t''(F_t')F_t''(r)=1$ and therefore to 
$f_t''(F_t')=\frac{1}{F_t''}$.
Changing notation ($V=F_t$), we obtain
$$
I = \int_0^\infty \frac{e^{-1/V(r)}}{V^{n+2}(r)}
\left(  \frac{V'(r)}{V''(r)}-r^2V'(r) + nr+V'(r) r -n \right)r^{n-1}dr .
$$
{\color{red}Cyril: the scope would be to prove that $I \geq 0$, maybe like in the previous section showing that, after some appropriate change of variables, the bracket is nothing but a square (or anything trivially positive)...}
\\{\color{blue}Devraj: I then end up with $$I = -\int_0^\infty \frac{e^{-1/V(r)}}{V^{n+2}(r)}
\left(  \frac{V'(r)}{V''(r)}-r^2V'(r) + nr+V'(r) r -n \right)r^{n-1}dr$$, so with the opposite sign, evidently, the difficult tasks remains}

% Later on we will use the following trivial fact
% $$
% L e^r=ne^r .
% $$
% {\color{red}Cyril: There might be a problem here since $r \mapsto e^r$ does not belong to $\mathbb{L}^2(\nu)$.}

% Denote by $P_t$, $t \geq 0$ the associated Laguerre semi-group and set 
% $$
% f_t = -\log P_t(e^{-g}),
% $$
% for some fixed convex increasing function $g \colon [0,\infty) \to [0,\infty)$ satisfying $g(0)=0$. For simplicity we set
% $F_t \coloneqq \mathcal{L}f_t$ that is also convex increasing and $F_t(0)=0$. In order to avoid unnecessary technicalities we assume that $g$ is strictly convex so that $F_t$ does not take infinite values.

% Our aim is to prove that
% \begin{align*}
% \alpha(t) 
% &\coloneqq
% \int_0^\infty e^{-f_t(r)}r^{n-1}dr 
% \int_0^\infty \Gamma(n+1, \frac{1}{F_t})r^{n-1}dr \\
% & = 
% \int_0^\infty P_t(e^{-g})(r)r^{n-1}dr 
% \int_0^\infty \Gamma(n+1, \frac{1}{F_t})r^{n-1}dr, \qquad t >0
% \end{align*}
% is non decreasing.

% Observe that
% \begin{align*}
% \frac{d}{dt} 
% \int_0^\infty P_t(e^{-g})(r)r^{n-1}dr
% & =
% \int_0^\infty LP_t(e^{-g})(r)r^{n-1}dr \\
% & =
% \Gamma(n) \int_0^\infty LP_t(e^{-g})(r) e^r d\nu(r) \\
% & =
% \Gamma(n) \int_0^\infty P_t(e^{-g})(r) L e^r d\nu(r) & \mbox{(by symmetry)}\\
% & =
% n \int_0^\infty P_t(e^{-g})(r) r^{n-1}dr &  \mbox{(by the trivial fact above)} .
% \end{align*}
% Differentiating in $t$, we therefore obtain that
% \begin{align*}
% \alpha'(t)
% & = 
% n \alpha(t)
% - 
% \int_0^\infty P_t(e^{-g})(r)r^{n-1}dr 
% \int_0^\infty \Gamma'(n+1, \frac{1}{F_t}) \frac{\frac{d}{dt}F_t(r)}{F_t(r)^2}r^{n-1}dr
% \end{align*}
% where $\Gamma'(a,x)=\frac{d}{dx}\Gamma(a,x)=-x^{a-1}e^{-x}$ is the derivative of the incomplete gamma function.
% Collecting the two term, we arrive at
% $$
% \alpha'(t)
%  = 
% \int_0^\infty P_t(e^{-g})(r)r^{n-1}dr 
% \int_0^\infty \left( n \Gamma(n+1, \frac{1}{F_t})  -\Gamma'(n+1, \frac{1}{F_t}) \frac{\frac{d}{dt}F_t(r)}{F_t(r)^2} \right) r^{n-1}dr .
% $$
% Hence the sign of $\alpha'$ is governed by the sign of the second integral
% $$
% I \coloneqq 
% \int_0^\infty \left( n \Gamma(n+1, \frac{1}{F_t})  -\Gamma'(n+1, \frac{1}{F_t}) \frac{\frac{d}{dt}F_t(r)}{F_t(r)^2} \right) r^{n-1}dr
% $$
% we now analyze.

% By Lemma \ref{lem:Cyril} we observe that
% \begin{align*}
% \frac{d}{dt}F_t 
% & = - \frac{d}{dt} f_t(F_t') \\
% & = 
% \frac{LP_t(e^{-g})(F_t')}{P_t(e^{-g})(F_t')} \\
% & =
% F_t'\frac{P_t(e^{-g})''(F_t')}{P_t(e^{-g})(F_t')} +\left(n-F_t'\right) \frac{P_t(e^{-g})'(F_t')}{P_t(e^{-g})(F_t')} .
% \end{align*}






% \newpage

% Recall that we "proved" earlier that
% $$
%     \sup_{f\in \tilde C(\mathbb{R}^n)} 
%     \int e^{-f}\int e^{-\mathcal{L}f}
%     =
%     (n\omega_{n})^2 
%     \sup_{\Psi} 
%     \int_0^\infty e^{-\Psi(r)}r^{n-1}dr 
%     \int_0^\infty e^{-\mathcal{L}\Psi}r^{n-1}dr 
% $$
% with $\Psi$ convex, non-decreasing with $\Psi(0)=0$.
% Changing $\Psi$ into $\Psi+c$ for any constant $c$ doesn't affect the volume product, therefore the constraint $\Psi(0)=0$ might be removed if needed.

% Since we will use the heat flow on the line, we extend $\Psi$ on the negative axis by parity. Moreover by an approximation argument
% we can assume that $\Psi$ is strictly convex. Our aim is therefore to compute
% $$
% M = \sup_{\Psi} 
%     \int_{-\infty}^\infty e^{-\Psi(r)}|r|^{n-1}dr 
%     \int_{-\infty}^\infty e^{-\mathcal{L}\Psi}|r|^{n-1}dr 
% $$
% for $\Psi \colon \mathbb{R} \to [0,\infty)$  even, non-decreasing  and strictly convex on $(0,\infty)$ (whatever the value of $\Psi(0)$). 
% Then, we would get
% $$
%     \sup_{f\in \tilde C(\mathbb{R}^n)} 
%     \int e^{-f}\int e^{-\mathcal{L}f}
%     =
% \frac{1}{4}(n\omega_{n})^2 M 
% $$
% and we are left with an estimation of $M$.
% To that aim, and following Cordero et al., consider $\Psi_t \coloneqq - \log P_tf$
% where $P_t$ is the one dimensional heat flow and $f(r)=e^{-\Psi(r)}$, $r \in \mathbb{R}$ for a given $\Psi$ as above. 
% %Since we can change the value of $\Psi(0)$, we assume that $\int e^{-\Psi(r)}dr=1$, turning $f$ into a probability density on the line.
% Now, for $x \in \mathbb{R}$ {\color{red}Cyril: I'm not sure about the normalization here, the $4t$ and $\sqrt{2t}$}
% $$
% P_tf(x) \coloneqq \mathbb{E}(x+\sqrt{2t}Z) = \frac{1}{\sqrt{4\pi t}} \int_{-\infty}^\infty f(y)e^{-|x-y|^2/4t}dy
% =
% \frac{1}{\sqrt{2t+1}}
% Q_{\log\sqrt{2t+1}}f(\frac{x}{\sqrt{2t+1}})
% $$
% with $Z \sim \mathcal{N}(0,1)$ a one dimensional standard normal. It is generated by the Laplacien $L$ which, in dimension 1, takes the form $Lf=f''$ for $f$ smooth enough. Here $Q$ is the Focker-Planck semi-group, that is, the adjoint of the Ornstein-Uhlenbeck semi-group in $\mathbb{L}^2(dx)$ (for the Lebesgue measure), see \eqref{eq:Focker-Planck}. One important fact is that $\lim_{t \to \infty} Q_tf = \frac{1}{\sqrt{2\pi}}e^{-r^2/2} \int_{-\infty}^\infty f(r)dr$.
% For $t \geq 0$, consider the following quantity
% $$
% \alpha(t) \coloneqq 
%     \int_{-\infty}^\infty e^{-\Psi_t(r)}|r|^{n-1}dr 
%     \int_{-\infty}^\infty e^{-\mathcal{L}\Psi_t}|r|^{n-1}dr 
%     =
%     \int_{-\infty}^\infty P_tf(r)|r|^{n-1}dr 
%     \int_{-\infty}^\infty e^{-\mathcal{L}\Psi_t}|r|^{n-1}dr
% $$
% and observe that 
% $$
% \alpha(0)= \int_{-\infty}^\infty e^{-\Psi(r)}|r|^{n-1}dr 
%     \int_{-\infty}^\infty e^{-\mathcal{L}\Psi}|r|^{n-1}dr
% $$
% while, (after change of variable $s=r/\sqrt{2t+1}$), {\color{red}Cyril: and some convergence theorem that I'm skipping but that should be controlled properly}
% \begin{align*}
% \lim_{t \to \infty} \alpha(t)
% & = \int_{-\infty}^\infty Q_\infty f(s) |s|^{n-1}ds
%    \lim_{t \to \infty}  \int_{-\infty}^\infty e^{-\mathcal{L}\Psi_t}|r|^{n-1}dr \\
%    & =
%   \frac{1}{\sqrt{2\pi}}\int_{-\infty}^\infty e^{-s^2/2} |s|^{n-1}ds \int_{-\infty}^\infty f(r)dr   \int_{-\infty}^\infty \sqrt{2\pi}e^{-s^2/2} |s|^{n-1}ds \left( \int_{-\infty}^\infty f(r)dr \right)^{-1} \\
%   & = 
%   4 (2^{\frac{n}{2}-1} \Gamma(n/2))^2 
% \end{align*}
% where we used, in the second line, similar computations than in \eqref{eq:Devraj1}, \eqref{eq:Devraj2}, \eqref{eq:Devraj3}.

% Assume one can prove that $\alpha$ is non-decreasing, one would get that 
% $M \leq 4 (2^{\frac{n}{2}-1} \Gamma(n/2))^2$, and therefore, after some algebra, the functional Blashke Santalo inequality
% $$
%  \sup_{f\in \tilde C(\mathbb{R}^n)} 
%     \int e^{-f}\int e^{-\mathcal{L}f}
%     = (2 \pi)^n .
% $$
% Therefore, to fully recover the proof by Cordero et al. it remains to prove that $\alpha$ is non decreasing.


% %Using Brascamp-Lieb Cordero et al prove that $\Phi_t''>0$
% %{\color{red}Cyril: computations are in the appendix of their paper, would be better to have a self-contained easy argument of this fact}. 
 
% {\color{red}Cyril: what is about to come is essentially wrong. $r^{n-1}$ is the bad guy. I don't know how to deal with it so far.}
 
% % Therefore $\Phi_t'$ has an inverse function we denote ${\Phi_t'}^{-1}$ and it is easy to check that
% % ${\Phi_t'}^{-1}=(\mathcal{L}\Phi_t)'$ and
% % \begin{align*}
% % \mathcal{L}\Phi_t (x)
% % & = x {\Phi_t'}^{-1}(x) - \Phi_t({\Phi_t'}^{-1}(x)) \\
% % & = x (\mathcal{L}\Phi_t)'(x) - \Phi_t((\mathcal{L}\Phi_t)'(x))
% % , \qquad x \in \mathbb{R}.
% %\end{align*}
% % Taking the derivative in $t$, it holds
% % $$
% % \partial_t \mathcal{L}\Phi_t (x) = 
% % x \partial_t (\mathcal{L}\Phi_t)'(x)
% % -\partial_t \Phi_t((\mathcal{L}\Phi_t)'(x))
% % - \Phi_t'((\mathcal{L}\Phi_t)'(x))\partial_t \mathcal{L}\Phi_t)'(x) .
% % $$
% % Now the identity ${\Phi_t'}((\mathcal{L}\Phi_t)'(x))=x$ leads to
% % $$
% % \partial_t \mathcal{L}\Phi_t (x) = 
% % -\partial_t \Phi_t((\mathcal{L}\Phi_t)'(x)) ,
% % \qquad x \in \mathbb{R}, \quad t >0  .
% %  $$
% For the heat flow,
% \begin{align*}
%     \partial_t \Phi_t
%     =
%     - \frac{LP_t(f)}{P_t(f)}
%     =
%     - \frac{P_t(f)''}{P_t(f)}
%     = 
%     \Phi_t'' - {\Phi_t'}^2 .
% \end{align*}
% We conclude that
% \begin{align*}
% \partial_t \mathcal{L}\Phi_t (x) 
% & = 
% -\Phi_t''((\mathcal{L}\Phi_t)'(x)) + {\Phi_t'}^2 ((\mathcal{L}\Phi_t)'(x)) \\
% & =
% -\Phi_t''({\Phi_t'}^{-1}(x)) + x^2 \\
% & =
% -\frac{1}{(\mathcal{L}\Phi_t)''(x)} + x^2 .
% \end{align*}
% Now, given $\Psi$ set
% $$
% \alpha(t) = \int_{-\infty}^\infty e^{-\Phi_t(r)}|r|^{n-1}dr 
%     \int_{-\infty}^\infty e^{-\mathcal{L}\Phi_t}|r|^{n-1}dr
% $$
% By construction of $\Phi_t$, and since the heat semi-group is invariant for the Lebesgue measure, in fact
% $$
% \alpha(t) = 
% \int_{-\infty}^\infty e^{-\Psi(r)}|r|^{n-1}dr 
%     \int_{-\infty}^\infty e^{-\mathcal{L}\Phi_t}|r|^{n-1}dr
% $$
% Therefore
% \begin{align*}
% \alpha'(t)
% & =
% - \int_{-\infty}^\infty e^{-\Psi(r)}|r|^{n-1}dr 
%     \int_{-\infty}^\infty \partial_t \mathcal{L}\Phi_t (x)  e^{-\mathcal{L}\Phi_t}|r|^{n-1}dr \\
%     & =
% \int_{-\infty}^\infty e^{-\Psi(r)}|r|^{n-1}dr 
%     \int_{-\infty}^\infty \left( \frac{1}{(\mathcal{L}\Phi_t)''(x)} - x^2 \right)  e^{-\mathcal{L}\Phi_t}|r|^{n-1}dr .
% \end{align*}


% \newpage

% {\color{red}Cyril: what comes next is Devraj's atempt to proving the polar functional Blashle-Santalo.}

% In this section we would like to show that $$
% \int_{\mathbb{{R}}^{n}}e^{-\phi(x)}dx\int_{\mathbb{{R}}^{n}}e^{-\phi(x)^\circ}dx\leq\int_{\mathbb{R}^{n}}e^{-|x|}dx\int_{\mathbb{{R}}^{n}}e^{-|x|^\circ}dx=\left[|S_{n-1}|(n-1)!\right]^{2}
% $$ where $\phi\in\Tilde{C}(\mathbb{R}^{n})$
% Using arguments shown above we have the following 
% \begin{equation}
%     \begin{split}
%         &\sup_{\phi\in\Tilde{C}(\mathbb{R}^{n})}\int_{\mathbb{{R}}^{n}}e^{-\phi(x)}dx\int_{\mathbb{{R}}^{n}}e^{-\phi(x)^\circ}dx\\
%         \leq&\sup_{\phi\in\Tilde{C}(\mathbb{R}^{n})}\int_{\mathbb{{R}}^{n}}e^{-\phi(|x|)}dx\int_{\mathbb{{R}}^{n}}e^{-\phi(|x|)^\circ}dx\\
%         \leq&\sup_{\phi\in\Tilde{C}(\mathbb{R}^{n})}|S_{n-1}|^2\int_{0}^{\infty}e^{-\phi(r)}r^{n-1}dr\int_{0}^{\infty}e^{-\phi(r)^\circ}r^{n-1}dr\\
%         =&\sup_{\phi\in\Tilde{C}(\mathbb{R}^{n}),\int_{0}^{\infty}e^{-\phi(r)}dr=1}|S_{n-1}|^2\int_{0}^{\infty}e^{-\phi(r)}r^{n-1}dr\int_{0}^{\infty}e^{-\phi(r)^\circ}r^{n-1}dr\\
%     \end{split}
% \end{equation}
% Note that if $\phi\in\Tilde{C}(\mathbb{R}^{n})$ and $\int_{0}^{\infty}e^{-\phi(r)}dr=1$ we must have that $$
% \int_{0}^{\infty}e^{-\phi(r)}r^{n-1}dr\leq\int_{0}^{\infty}e^{-r}r^{n-1}dr\leq(n-1)!
% $$ This gives the following $$
% \sup_{\phi\in\Tilde{C}(\mathbb{R}^{n})}\int_{\mathbb{{R}}^{n}}e^{-\phi(x)}dx\int_{\mathbb{{R}}^{n}}e^{-\phi(x)^\circ}dx\leq\sup_{\phi\in\Tilde{C}(\mathbb{R}^{n}),\int_{0}^{\infty}e^{-\phi(r)}dr=1}|S_{n-1}|^{2}(n-1)!\int_{0}^{\infty}e^{-\phi(r)^\circ}r^{n-1}dr
% $$ It is rather evident that one much show that $$
% \sup_{\phi\in\Tilde{C}(\mathbb{R}^{n}),\int_{0}^{\infty}e^{-\phi(r)}dr=1}\int_{0}^{\infty}e^{-\phi(r)^\circ}r^{n-1}dr\leq (n-1)!
% $$ To simplify notation, for $f=e^{-\phi}$, define $$
% M(f):=\int_{0}^{\infty}exp\left[-\left(-log f\right)^\circ\right]r^{n-1}dr
% $$ To this aim, we introduce the Laguerre semigroup $P_{t}$ that is associated to the one sided exponential measure $\lambda,\frac{d\lambda}{dx}=e^{-x}, x>0$. We need the dual semigroup (with respect to the Lebesgue measure on the half line) $\Tilde{P}_{t}$ of this semigroup which satisfies the following $$
% \int_{0}^{\infty}g(x)P_{t}f(x)dx=\int_{0}^{\infty}f(x)\Tilde{P}_{t}g(x)dx
% $$ From the above, we see that $$
% \Tilde{P}_{0}f=f, \Tilde{P}_{\infty}f=e^{-r}\int_{0}^{\infty}f(r)dr
% $$ Note that $\Tilde{P}_{t}f(x)$ is the solution to the following differential equation $$
% \partial_{t}q(t,x)=q(t,x)+(x+1)\frac{d}{dx}q(t,x)+x\frac{d^2}{dx^2}q(t,x),q(0,x)=f(x)
% $$ We need to show that $\alpha(t)=log M(\Tilde{P}_{t}f)$ is non-decreasing in time for $-log(f)\in\Tilde{C}(R^{n}),\int_{0}^{\infty}f(r)dr=1$. This would yield our result. $$
% \alpha'(t)=\int_{0}^{\infty}-\frac{d}{dt}[-log \Tilde{P}_{t}(f)(r)]^\circ\frac{\mu_{t}}{\int_{0}^{\infty}\mu_{t}(r)dr}dr
% $$ where $$
% \mu_{t}(r)=exp\left[-(-log \Tilde{P}_{t}f)^\circ\right]r^{n-1}
% $$ 
% {\color{blue}Devraj:The above is precisely what was done in the article by Cordero et Al. Regrettably, we've to work with this strange semigroup as the scaling argument won't work well with the polar transform (meaning we don't have $(h+c)^\circ=h^\circ-c$ in stark contrast to the Legendre transform. To follow their article, we would have to show \begin{enumerate}
%     \item $\mu_{t}$ is strictly log concave (so we can use some brascamp-lieb inequality). For this it may be better to use the relationship with the Orstein-Uhlenbeck semigroup as things are much more tractable. 
%     \item Compute $\frac{d}{dt}[-log \Tilde{P}_{t}(f)(r)]^\circ$ to see if for some choice of function in the brascamp lieb we can bound this by something nice.
% \end{enumerate}
% But to be quite honest...the computations are proving to be rather atrocious :( 
% }



\section{Comparison between polar and Legendre transform}

{\color{red}Cyril: what comes in this section is some sort of attempt to proving the polar functional Blaschke-Santalo... Will most probably be removed later on...}

Our first aim is to prove the following general relationship between the polar transform and the Legendre transform, consequence of the layer cake formula and the following fact, proved earlier: 
$\forall \lambda>0, \{f^\circ>\lambda\}=\lambda\{\mathcal{L}f>\frac{1}{\lambda}\}$.



\begin{proposition} \label{prop:polar-legendre}
Let $J \colon (0,\infty) \to (0,\infty)$ be a $\mathcal{C}^1$ non-increasing function converging to $0$ at infinity and $H \colon (0,\infty) \to (0,\infty)$ be the primitive of $u \in (0,\infty) \mapsto  J'(1/u)/u^{2+n}$ that converges to $0$ at infinity. By convention $H(\infty)=J(\infty)=0$.
Then, for any
$f\in \mathrm{Cvx}_0(\mathbb{R}^{n})$, it holds (in $[0,\infty]$)
$$
\int H( f^\circ) = \int J(\mathcal{L}f) .
$$
\end{proposition}

\begin{proof}
In order to deal with finite integrals, fix $0 < \varepsilon < M$ and set $B=B_{\varepsilon,M}\coloneqq\{x : \varepsilon \leq |x| \leq M \}$. Then, it holds by Fubini-Tonelli's theorem
\begin{align*}
\int_{B} H(f^\circ(x))dx 
& =
\int_{\mathbb{R}^n} \int_{f^\circ(x)}^\infty -H'(\lambda)d\lambda \mathds{1}_B dx \\
& =
\int_{\mathbb{R}^n}
\int_{0}^\infty -H'(\lambda)  \mathds{1}_{f^\circ(x) \leq  \lambda }  \mathds{1}_B d\lambda dx \\
& =
\int_{0}^\infty -H'(\lambda) |\{ f^\circ(x) \leq  \lambda \} \cap B|d\lambda
\end{align*}
where as usual $|\cdot|$ denote the Lebesgue measure.    

Since $\{f^\circ>\lambda\}=\lambda\{\mathcal{L}f>\frac{1}{\lambda}\}$ we infer that 
$\{f^\circ \leq \lambda\}=\lambda\{\mathcal{L}f \leq \frac{1}{\lambda}\}$, $\lambda >0$. Therefore, by homogeneity of the Lebesgue measure, 
$$
|\{f^\circ \leq \lambda\} \cap B| =|\lambda\{\mathcal{L}f \leq \frac{1}{\lambda}\} \cap B| = \lambda^n |\{\mathcal{L}f \leq \frac{1}{\lambda}\} \cap B_\lambda|$$ 
with 
$B_\lambda \coloneqq\{x : \frac{\varepsilon}{\lambda} \leq |x| \leq \frac{M}{\lambda} \}$. 
Hence
\begin{align*}
\int_{B} H(f^\circ(x))dx 
& =
\int_{0}^\infty -H'(\lambda) \lambda^n |\{\mathcal{L}f \leq \frac{1}{\lambda}\} \cap B_\lambda|d\lambda \\
& =
\int_{0}^\infty J'(1/\lambda) |\{\mathcal{L}f \leq \frac{1}{\lambda}\} \cap B_\lambda| (-\frac{d\lambda}{\lambda^2}) .
\end{align*}
Changing variable $t = \frac{1}{\lambda}$, it follows by Fubini Tonelli's theorem that 
\begin{align*}
\int_{B} H(f^\circ(x))dx 
& =
- \int_{0}^\infty J'(t) |\{\mathcal{L}f \leq t\} \cap B_{1/t}| dt \\
& =
\int_{\mathbb{R}^n}
\int_{0}^\infty -J'(t)  \mathds{1}_{\mathcal{L}f(x) \leq  t}  \mathds{1}_{B_{1/t}}(x) dtdx  
.
\end{align*}
Taking the limit $\varepsilon \to 0$ and $M \to \infty$, by the monotone convergence theorem, we end up with
\begin{align*}
\int_{\mathbb{R}^n} H(f^\circ(x))dx 
& =
\int_{\mathbb{R}^n}
\int_{0}^\infty -J'(t)  \mathds{1}_{\mathcal{L}f(x) \leq  t}  dtdx  \\
& =
\int_{\mathbb{R}^n}
J(\mathcal{L}f(x))dx 
.
\end{align*}
This ends the proof of the proposition.
\end{proof}


Consider as an example of illustration 
$J(u)=\frac{1}{u^p}$ with $p>0$, $u>0$. 
Then
$H'(u)=-\frac{p}{u^{p+n+3}}$ and in turn
$H(u)=\frac{p}{n+p+2} \frac{1}{u^{n+p+2}}$, so that
$$
p \int \frac{1}{(f^\circ)^{n+p+2}} = (n+p+2) \int \frac{1}{(\mathcal{L}f)^p} .
$$
As a second example consider $H(u)=e^{-u}$, $u>0$. Then $J'(u)=H'(1/u)/u^{n+2}=-\frac{e^{-1/u}}{u^{n+2}}$, $u>0$. In particular changing variable $v=1/u$ we recognize the incomplete gamma function 
$\Gamma(a,x)=\int_x^\infty t^{a-1}e^{-t}dt$ with $a=n+1$. Hence
$$
J(u)=\Gamma(n+1,1/u)
$$
so that
$J(u)=n!e^{-1/u} \sum_{k=0}^n \frac{1}{u^kk!}$,$u>0$, 
and we have
$$
\int e^{-f^\circ} = \int \Gamma(n+1, \frac{1}{\mathcal{L}f}) .
$$
The functional polar product (Blashke Santalo) introduced by Artstein-Avidan and Slomka can therefore be restated, for $f \in \tilde C(\mathbb{R}^n)$, as
$$
\int e^{-f} \int e^{-f^\circ}
=
\int e^{-f} 
\int \Gamma(n+1, \frac{1}{\mathcal{L}f}) .
$$



